% Generated mechanically from the frozen spaces-properties.tex target.
% Do not edit this generated file; edit the bound locale source instead.

% OK, start here.
%

\setcounter{chapter}{65}
\chapter{代数空间的性质}



\phantomsection
\label{spaces-properties-section-phantom}


\section{引言}
\label{spaces-properties-section-introduction}

\noindent
关于代数空间的简要介绍，请参见
“空间”第 \ref{spaces-section-introduction} 节；关于代数空间的基本定义与约定，
也请阅读该章的有关部分。本章开始介绍代数空间的一些基本概念和性质。
对于拟分离代数空间的情形，基础参考文献是 \cite{Kn}。

\medskip\noindent
由于我们在编排上决定先讨论代数空间本身的性质，稍后再讨论代数空间态射的性质，
有些地方的论述会略显不便。代数空间的{\it \'etale 态射}是这条规则的例外，
我们将在第 \ref{spaces-properties-section-etale-morphisms} 节中引入它。
但在到达该节之前，每当我们说某个态射具有某种性质时，均自动假定该态射的源是概形
（或者该态射是可表示的）。

\medskip\noindent
本章的部分内容（尤其是关于点的内容）将在“良好代数空间”一章中得到加强。


\section{约定}
\label{spaces-properties-section-conventions}

\noindent
贯穿本章的约定是：所有概形都包含在一个大 fppf 站点
$\Sch_{fppf}$ 中；所考虑的每个环 $A$ 都满足
$\Spec(A)$（同构于）该大站点的一个对象。

\medskip\noindent
设 $S$ 为概形，$X$ 为 $S$ 上的代数空间。在本章及后续各章中，
我们把 $X$ 与自身的积（在 $S$ 上代数空间的范畴中）记作
$X \times_S X$，而不记作 $X \times X$。这样做是为了避免在变更基概形时
产生混淆；参见“空间”第 \ref{spaces-section-change-base-scheme} 节。


\section{分离公理}
\label{spaces-properties-section-separation}

\noindent
本节汇集代数空间的所有“绝对”分离条件。在我们的表述中，任一代数空间都是某个
确定基概形上的代数空间。因此，$S$ 上代数空间 $X$ 的任一绝对性质，都对应于把
$X$ 视作 $\Spec(\mathbf{Z})$ 上的代数空间时所满足的一个条件。精确定义如下。

\begin{definition}
\label{spaces-properties-definition-separated}
（比较“空间”定义 \ref{spaces-definition-separated}。）
考虑大 fppf 站点
$\Sch_{fppf} = (\Sch/\Spec(\mathbf{Z}))_{fppf}$.
设 $X$ 是 $\Spec(\mathbf{Z})$ 上的代数空间，并设
$\Delta : X \to X \times X$ 为对角态射。
\begin{enumerate}
\item 若 $\Delta$ 是闭浸入，则称 $X$ {\it 分离}。
\item 若 $\Delta$ 是浸入，则称 $X$
{\it 局部分离}\footnote{文献中，这通常指同时拟分离且局部分离的代数空间。}。
\item 若 $\Delta$ 拟紧，则称 $X$ {\it 拟分离}。
\item 若存在 Zariski 覆盖
$X = \bigcup_{i \in I} X_i$（见“空间”定义
\ref{spaces-definition-Zariski-open-covering}），使得每个 $X_i$ 都拟分离，
则称 $X$ {\it Zariski 局部拟分离}\footnote{这一概念由 B.\ Conrad 建议。}。
\end{enumerate}
设 $S$ 是包含在 $\Sch_{fppf}$ 中的概形，$X$ 是 $S$ 上的代数空间。
若把 $X$ 视作 $\Spec(\mathbf{Z})$ 上的代数空间（见“空间”定义
\ref{spaces-definition-base-change}）时具有相应性质，则称 $X$
{\it 分离}、{\it 局部分离}、{\it 拟分离}或
{\it Zariski 局部拟分离}。
\end{definition}

\noindent
若 $S$ 上的代数空间 $X$ 按上述绝对意义分离，则它在 $S$ 上分离；
对上述其他绝对分离性质亦同。我们将在“空间的态射”第
\ref{spaces-morphisms-section-separation-axioms} 节中详细讨论这一点。
引理 \ref{spaces-properties-lemma-quasi-separated-quasi-compact-pieces} 将表明，
Zariski 局部分离与基概形无关（因而等价于绝对概念）。

\begin{lemma}
\label{spaces-properties-lemma-trivial-implications}
设 $S$ 为概形，$X$ 为 $S$ 上的代数空间。定义
\ref{spaces-properties-definition-separated} 的各分离公理之间有下列蕴含关系：
\begin{enumerate}
\item 分离蕴含其余所有条件；
\item 拟分离蕴含 Zariski 局部拟分离。
\end{enumerate}
\end{lemma}

\begin{proof}
略。
\end{proof}

\begin{lemma}
\label{spaces-properties-lemma-characterize-quasi-separated}
设 $S$ 为概形，$X$ 为 $S$ 上的代数空间。下列条件等价：
\begin{enumerate}
\item $X$ 是拟分离代数空间；
\item 对任意 $U \to X$、$V \to X$，若 $U$、$V$ 是拟紧概形，
则纤维积 $U \times_X V$ 拟紧；
\item 对任意 $U \to X$、$V \to X$，若 $U$、$V$ 是仿射概形，
则纤维积 $U \times_X V$ 拟紧。
\end{enumerate}
\end{lemma}

\begin{proof}
由“空间”引理
\ref{spaces-lemma-category-of-spaces-over-smaller-base-scheme}
可设 $S = \Spec(\mathbf{Z})$。由于
$U \times_X V = X \times_{X \times X} (U \times V)$，
且当 $U$、$V$ 拟紧时 $U \times V$ 也拟紧，所以 (1) 蕴含 (2)。
显然 (2) 蕴含 (3)。现假设 (3)。选取概形 $W$ 以及满 \'etale 态射
$W \to X$。于是 $W \times W \to X \times X$ 是满 \'etale 态射。
因此，只需证明
$$
j : W \times_X W = X \times_{(X \times X)} (W \times W) \to W \times W
$$
拟紧；见“空间”引理
\ref{spaces-lemma-descent-representable-transformations-property}。
若 $U \subset W$ 和 $V \subset W$ 是仿射开集，则由假设，
$j^{-1}(U \times V) = U \times_X V$ 拟紧。仿射开集 $U \times V$
构成 $W \times W$ 的仿射开覆盖（“概形”引理
\ref{schemes-lemma-affine-covering-fibre-product}），故结论由“概形”引理
\ref{schemes-lemma-quasi-compact-affine} 得出。
\end{proof}

\begin{lemma}
\label{spaces-properties-lemma-characterize-separated}
设 $S$ 为概形，$X$ 为 $S$ 上的代数空间。下列条件等价：
\begin{enumerate}
\item $X$ 是分离代数空间；
\item 对任意 $U \to X$、$V \to X$，若 $U$、$V$ 是仿射概形，
则纤维积 $U \times_X V$ 仿射，并且
$$
\mathcal{O}(U) \otimes_\mathbf{Z} \mathcal{O}(V)
\longrightarrow
\mathcal{O}(U \times_X V)
$$
是满射。
\end{enumerate}
\end{lemma}

\begin{proof}
由“空间”引理
\ref{spaces-lemma-category-of-spaces-over-smaller-base-scheme}
可设 $S = \Spec(\mathbf{Z})$。由于
$U \times_X V = X \times_{X \times X} (U \times V)$，
且当 $U$、$V$ 仿射时 $U \times V$ 也仿射，所以 (1) 蕴含 (2)。
现假设 (2)。选取概形 $W$ 以及满 \'etale 态射 $W \to X$。
于是 $W \times W \to X \times X$ 是满 \'etale 态射。因此，只需证明
$$
j : W \times_X W = X \times_{(X \times X)} (W \times W) \to W \times W
$$
是闭浸入；见“空间”引理
\ref{spaces-lemma-descent-representable-transformations-property}。
若 $U \subset W$ 和 $V \subset W$ 是仿射开集，则由假设，
$j^{-1}(U \times V) = U \times_X V$ 仿射；又因相应的环同态为满射，
映射 $U \times_X V \to U \times V$ 是闭浸入。仿射开集
$U \times V$ 构成 $W \times W$ 的仿射开覆盖
（“概形”引理 \ref{schemes-lemma-affine-covering-fibre-product}），
故结论由“态射”引理 \ref{morphisms-lemma-closed-immersion} 得出。
\end{proof}







\section{代数空间的点}
\label{spaces-properties-section-points}

\noindent
由“空间”例 \ref{spaces-example-affine-line-translation} 可清楚看出，
不应把代数空间的点定义为从域谱出发的单态射。
相反，我们完全按照“概形”第 \ref{schemes-section-points} 节的说明，
把它们定义为从域谱出发的态射之等价类。

\medskip\noindent
设 $S$ 为概形，$F$ 是 $(\Sch/S)_{fppf}$ 上的预层，$K$ 为域。
考虑态射
$$
\Spec(K) \longrightarrow F.
$$
由 Yoneda 引理，该态射由元素 $p \in F(\Spec(K))$ 给出。
若存在第三个域 $\Omega$ 以及交换图表
$$
\xymatrix{
\Spec(\Omega) \ar[r] \ar[d] &
\Spec(L) \ar[d]^q \\
\Spec(K) \ar[r]^p &
F.
}
$$
则称这样的两个二元组 $(\Spec(K), p)$ 和 $(\Spec(L), q)$
{\it 等价}。换言之，存在域扩张 $K \to \Omega$ 和 $L \to \Omega$，
使 $p$ 与 $q$ 映到 $F(\Spec(\Omega))$ 中的同一元素。
略去该关系确为等价关系的验证。

\begin{definition}
\label{spaces-properties-definition-points}
设 $S$ 为概形，$X$ 为 $S$ 上的代数空间。$X$ 的一个{\it 点}，
是从域谱到 $X$ 的态射之等价类。$X$ 的点集记作 $|X|$。
\end{definition}

\noindent
注意，若 $f : X \to Y$ 是 $S$ 上代数空间的态射，则它诱导映射
$|f| : |X| \to |Y|$，把代表元 $x : \Spec(K) \to X$
送到代表元 $f \circ x : \Spec(K) \to Y$。

\begin{lemma}
\label{spaces-properties-lemma-scheme-points}
设 $S$ 为概形，$X$ 为 $S$ 上的概形。把 $X$ 视作概形所得的点，
与把 $X$ 视作代数空间所得的点之间有典范一一对应。
\end{lemma}

\begin{proof}
这就是“概形”引理 \ref{schemes-lemma-characterize-points}。
\end{proof}

\begin{lemma}
\label{spaces-properties-lemma-points-cartesian}
设 $S$ 为概形，并设
$$
\xymatrix{
Z \times_Y X \ar[r] \ar[d] & X \ar[d] \\
Z \ar[r] & Y
}
$$
是 $S$ 上代数空间的 Cartesian 图表。则点集映射
$$
|Z \times_Y X|
\longrightarrow
|Z| \times_{|Y|} |X|
$$
是满射。
\end{lemma}

\begin{proof}
具体地，设给定域 $K$、$L$ 以及态射
$\Spec(K) \to X$、$\Spec(L) \to Z$。它们作为 $|Y|$ 中的元素相同，
意味着存在公共扩张 $M/K$ 和 $M/L$，使得
$\Spec(M) \to \Spec(K) \to X \to Y$ 与
$\Spec(M) \to \Spec(L) \to Z \to Y$ 相同。
这恰好是得到态射 $\Spec(M) \to Z \times_Y X$ 的条件。
\end{proof}

\begin{lemma}
\label{spaces-properties-lemma-characterize-surjective}
设 $S$ 为概形，$X$ 为 $S$ 上的代数空间，并设
$f : T \to X$ 是从概形到 $X$ 的态射。下列条件等价：
\begin{enumerate}
\item 按“空间”定义
\ref{spaces-definition-relative-representable-property}，
$f : T \to X$ 是满射；
\item $|f| : |T| \to |X|$ 是满射。
\end{enumerate}
\end{lemma}

\begin{proof}
假设 (1)。设 $x : \Spec(K) \to X$ 是从域谱到 $X$ 的态射。
由假设，概形态射 $\Spec(K) \times_X T \to \Spec(K)$ 是满射。
因此存在域扩张 $K'/K$ 以及态射
$\Spec(K') \to \Spec(K) \times_X T$，使下图左侧方块交换：
$$
\xymatrix{
\Spec(K') \ar[r] \ar[d] &
\Spec(K) \times_X T \ar[d] \ar[r] &
T \ar[d] \\
\Spec(K) \ar@{=}[r] &
\Spec(K) \ar[r]^-x & X
}
$$
这表明 $|f| : |T| \to |X|$ 是满射。

\medskip\noindent
假设 (2)。设 $Z \to X$ 为态射，其中 $Z$ 是概形。我们必须证明概形态射
$Z \times_X T \to T$ 是满射，即 $|Z \times_X T| \to |Z|$ 是满射。
这由 (2) 与引理 \ref{spaces-properties-lemma-points-cartesian} 得出。
\end{proof}

\begin{lemma}
\label{spaces-properties-lemma-points-presentation}
设 $S$ 为概形，$X$ 为 $S$ 上的代数空间，并设 $X = U/R$
为 $X$ 的一个表示；见“空间”定义 \ref{spaces-definition-presentation}。
则 $|R| \to |U| \times |U|$ 的像是等价关系，并且 $|X|$
是 $|U|$ 对该等价关系的商。
\end{lemma}

\begin{proof}
该假设意味着：$U$ 是概形，$p : U \to X$ 是满 \'etale 态射，
$R = U \times_X U$ 是概形，并在 $U$ 上定义一个 \'etale 等价关系，
使得作为层有 $X = U/R$。由引理
\ref{spaces-properties-lemma-characterize-surjective}，$|U| \to |X|$ 是满射。
由引理 \ref{spaces-properties-lemma-points-cartesian}，映射
$$
|R| \longrightarrow |U| \times_{|X|} |U|
$$
是满射。因此，$|R| \to |U| \times |U|$ 的像恰好是所有二元组
$(u_1, u_2) \in |U| \times |U|$ 所成的集合，其中 $u_1$ 与 $u_2$
在 $|X|$ 中的像相同。合并这两个陈述即得本引理。
\end{proof}

\begin{lemma}
\label{spaces-properties-lemma-topology-points}
设 $S$ 为概形。$S$ 上代数空间的点集上存在唯一的拓扑，满足下列性质：
\begin{enumerate}
\item 若 $X$ 是 $S$ 上的概形，则 $|X|$ 上的拓扑是通常拓扑
（经由引理 \ref{spaces-properties-lemma-scheme-points} 的识别）；
\item 对每个 $S$ 上代数空间的态射 $X \to Y$，映射
$|X| \to |Y|$ 连续；
\item 对每个 \'etale 态射 $U \to X$，若 $U$ 是概形，则拓扑空间映射
$|U| \to |X|$ 连续且开。
\end{enumerate}
\end{lemma}

\begin{proof}
设 $X$ 是 $S$ 上的代数空间，并设 $p : U \to X$ 为满 \'etale 态射，
其中 $U$ 是 $S$ 上的概形。定义 $W \subset |X|$ 为开集，当且仅当
$|p|^{-1}(W)$ 是 $|U|$ 的开子集。这在 $|X|$ 上定义了一个拓扑
（即 $|X|$ 上的商拓扑；见“拓扑”引理 \ref{topology-lemma-quotient}）。

\medskip\noindent
下面证明该拓扑与表示的选择无关。为此，只需说明：若 $U'$ 是概形，
且 $U' \to X$ 是 \'etale 态射，则映射 $|U'| \to |X|$
（其中 $|X|$ 的拓扑按上述 $U \to X$ 定义）开且连续。
这还将证明 (3) 成立。令 $U'' = U \times_X U'$，于是有交换图表
$$
\xymatrix{
U'' \ar[r] \ar[d] & U' \ar[d] \\
U \ar[r] & X
}
$$
由于 $U \to X$ 和 $U' \to X$ 都是 \'etale 态射，
$U'' \to U$ 与 $U'' \to U'$ 都是概形的 \'etale 态射。
此外，$U'' \to U'$ 是满射。因此得到点集映射的交换图表
$$
\xymatrix{
|U''| \ar[r] \ar[d] & |U'| \ar[d] \\
|U| \ar[r] & |X|
}
$$
下方横箭头是满射（见引理 \ref{spaces-properties-lemma-characterize-surjective}
或引理 \ref{spaces-properties-lemma-points-presentation}），并且按 $|X|$ 上拓扑的定义连续。
由“态射”引理 \ref{morphisms-lemma-etale-open}，上方横箭头满、连续且开；
左侧竖箭头也连续且开（再次使用“态射”引理
\ref{morphisms-lemma-etale-open}）。于是形式地可知，
右侧竖箭头连续且开。

\medskip\noindent
最后证明 (2)。设 $a : X \to Y$ 为代数空间的态射。根据“空间”引理
\ref{spaces-lemma-lift-morphism-presentations}，可以找到图表
$$
\xymatrix{
U \ar[d]_p \ar[r]_\alpha & V \ar[d]^q \\
X \ar[r]^a & Y
}
$$
其中 $U$ 与 $V$ 是概形，$p$ 与 $q$ 是满 \'etale 态射。它给出图表
$$
\xymatrix{
|U| \ar[d]_p \ar[r]_\alpha & |V| \ar[d]^q \\
|X| \ar[r]^a & |Y|
}
$$
其中除下方横箭头外的所有箭头均已知连续，并且两个竖箭头都满且开。
由此得到所需的下方横箭头连续。
\end{proof}

\begin{definition}
\label{spaces-properties-definition-topological-space}
设 $S$ 为概形，$X$ 为 $S$ 上的代数空间。$X$ 的底层{\it 拓扑空间}
是点集 $|X|$ 连同引理 \ref{spaces-properties-lemma-topology-points} 所构造的拓扑。
\end{definition}

\noindent
事实证明，该拓扑空间携带的信息与“空间”定义
\ref{spaces-definition-small-Zariski-site} 的小 Zariski 站点
$X_{Zar}$ 相同。

\begin{lemma}
\label{spaces-properties-lemma-open-subspaces}
设 $S$ 为概形，$X$ 为 $S$ 上的代数空间。
\begin{enumerate}
\item 规则 $X' \mapsto |X'|$ 在 $X$ 的开子空间 $X'$
（见“空间”定义 \ref{spaces-definition-immersion}）与拓扑空间 $|X|$
的开集之间定义一个保持包含关系的双射。
\item $X$ 的开子空间族 $\{X_i \subset X\}_{i \in I}$
是 Zariski 覆盖（见“空间”定义
\ref{spaces-definition-Zariski-open-covering}），当且仅当
$|X| = \bigcup |X_i|$。
\end{enumerate}
换言之，$X$ 的小 Zariski 站点 $X_{Zar}$ 典范地等同于与拓扑空间
$|X|$ 相关联的站点（见“站点”例 \ref{sites-example-site-topological}）。
\end{lemma}

\begin{proof}
为证明 (1)，构造该规则的逆。设 $W \subset |X|$ 为开集。
选取表示 $X = U/R$，它对应于满 \'etale 映射 $p : U \to X$
及 \'etale 映射 $s, t : R \to U$。由构造，$|p|^{-1}(W)$ 是 $U$ 的开集。
以 $W' \subset U$ 表示相应的开子概形。显然，
$R' = s^{-1}(W') = t^{-1}(W')$ 是 $R$ 的 Zariski 开集，
并在 $W'$ 上定义 \'etale 等价关系。由“空间”引理
\ref{spaces-lemma-finding-opens}，态射 $X' = W'/R' \to X$ 是开浸入。
再由“空间”引理 \ref{spaces-lemma-representable-over-space}，
$X'$ 是代数空间。按构造 $|X'| = W$，即 $X'$ 是与 $W$ 对应的
$X$ 的子空间。这证明了 (1)。

\medskip\noindent
为证明 (2)，注意：若 $\{X_i \subset X\}_{i \in I}$ 是开子空间族，
则它是 Zariski 覆盖，当且仅当
$U = \bigcup U \times_X X_i$ 是开覆盖。这由 Zariski 覆盖的定义以及
态射 $U \to X$ 作为 $(\Sch/S)_{fppf}$ 上预层的映射是满射这一事实得出。
另一方面，由引理 \ref{spaces-properties-lemma-points-presentation}（以及投影
$U \times_X X_i \to X_i$ 满且 \'etale 这一事实），
$|X| = \bigcup |X_i|$ 当且仅当
$U = \bigcup U \times_X X_i$。于是得到 (2) 的等价性。
\end{proof}

\begin{lemma}
\label{spaces-properties-lemma-factor-through-open-subspace}
设 $S$ 为概形，$X$、$Y$ 为 $S$ 上的代数空间，$X' \subset X$
为开子空间，并设 $f : Y \to X$ 为 $S$ 上代数空间的态射。
则 $f$ 经过 $X'$ 分解，当且仅当 $|f| : |Y| \to |X|$
经过 $|X'| \subset |X|$ 分解。
\end{lemma}

\begin{proof}
由“空间”引理 \ref{spaces-lemma-base-change-immersions}，
$Y' = Y \times_X X' \to Y$ 是开浸入。若 $|f|(|Y|) \subset |X'|$，
则显然 $|Y'| = |Y|$。故由引理 \ref{spaces-properties-lemma-open-subspaces}，
$Y' = Y$。
\end{proof}

\begin{lemma}
\label{spaces-properties-lemma-etale-image-open}
设 $S$ 为概形，$X$ 为 $S$ 上的代数空间。设 $U$ 为概形，
$f : U \to X$ 为 \'etale 态射。令 $X' \subset X$ 为通过引理
\ref{spaces-properties-lemma-open-subspaces} 与开集 $|f|(|U|) \subset |X|$ 对应的开子空间。
则 $f$ 经过满 \'etale 态射 $f' : U \to X'$ 分解。
此外，若 $R = U \times_X U$，则 $R = U \times_{X'} U$，
且 $X'$ 有表示 $X' = U/R$。
\end{lemma}

\begin{proof}
分解的存在性由引理 \ref{spaces-properties-lemma-factor-through-open-subspace} 得出。
根据引理 \ref{spaces-properties-lemma-characterize-surjective}，态射 $f'$ 是满射。
为说明 $f'$ 是 \'etale 态射，设 $T \to X'$ 为态射，其中 $T$ 是概形。
由于 $X' \to X$ 是层的单态射，有
$T \times_X U = T \times_{X'} U$。又因假设 $f$ 为 \'etale 态射，
投影 $T \times_{X'} U \to T$ 是 \'etale 态射。
由于 $X' \to X$ 是单态射，还有
$U \times_X U = U \times_{X'} U$。最后，
$X' = U/R$ 由“空间”引理 \ref{spaces-lemma-space-presentation} 得出。
\end{proof}

\begin{lemma}
\label{spaces-properties-lemma-equivalence-class-point-monomorphism}
设 $S$ 为概形，$X$ 为 $S$ 上的代数空间。设
$p : \Spec(K) \to X$ 与 $q : \Spec(L) \to X$ 为态射，
其中 $K$ 和 $L$ 是域。假设 $p$ 与 $q$ 确定 $|X|$ 中的同一点，
且 $p$ 是单态射。则 $q$ 唯一地经过 $p$ 分解。
\end{lemma}

\begin{proof}
由于 $p$ 与 $q$ 定义 $|X|$ 中的同一点，概形
$$
Y = \Spec(K) \times_{p, X, q} \Spec(L)
$$
非空。单态射的基变换仍是单态射，故投影态射 $Y \to \Spec(L)$
是单态射。因此 $Y = \Spec(L)$；见“概形”引理
\ref{schemes-lemma-mono-towards-spec-field}。由此可知 $q$ 经过 $p$ 分解。
唯一性来自 $p$ 是单态射这一事实。
\end{proof}

\begin{lemma}
\label{spaces-properties-lemma-points-monomorphism}
设 $S$ 为概形，$X$ 为 $S$ 上的代数空间。考虑映射
$$
\{\Spec(k) \to X \text{ 为单态射，其中 }k\text{ 是域}\}
\longrightarrow
|X|
$$
该映射是单射。
\end{lemma}

\begin{proof}
这由引理 \ref{spaces-properties-lemma-equivalence-class-point-monomorphism} 得出。
\end{proof}

\noindent
我们将在“良好空间”引理
\ref{decent-spaces-lemma-decent-points-monomorphism} 中看到：
当 $X$ 良好时，引理 \ref{spaces-properties-lemma-points-monomorphism} 的映射是双射。

















\section{拟紧代数空间}
\label{spaces-properties-section-quasi-compact}

\begin{definition}
\label{spaces-properties-definition-quasi-compact}
设 $S$ 为概形，$X$ 为 $S$ 上的代数空间。若存在满 \'etale 态射
$U \to X$，其中 $U$ 是拟紧概形，则称 $X$ {\it 拟紧}。
\end{definition}

\begin{lemma}
\label{spaces-properties-lemma-quasi-compact-space}
设 $S$ 为概形，$X$ 为 $S$ 上的代数空间。则 $X$ 拟紧，当且仅当
$|X|$ 拟紧。
\end{lemma}

\begin{proof}
选取概形 $U$ 以及满 \'etale 态射 $U \to X$。我们将使用引理
\ref{spaces-properties-lemma-characterize-surjective}。若 $U$ 拟紧，则因
$|U| \to |X|$ 是满射，可知 $|X|$ 拟紧。若 $|X|$ 拟紧，
则因 $|U| \to |X|$ 是开映射，存在拟紧开集 $U' \subset U$，
使 $|U'| \to |X|$ 是满射（并且仍为 \'etale 态射）。结论得证。
\end{proof}

\begin{lemma}
\label{spaces-properties-lemma-finite-disjoint-quasi-compact}
有限个拟紧代数空间的不交并是拟紧代数空间。
\end{lemma}

\begin{proof}
这由引理 \ref{spaces-properties-lemma-quasi-compact-space} 及相应的拓扑事实立即得出。
\end{proof}

\begin{example}
\label{spaces-properties-example-quasi-compact-not-very-reasonable}
空间 $\mathbf{A}^1_{\mathbf{Q}}/\mathbf{Z}$ 是拟紧代数空间。
\end{example}

\begin{lemma}
\label{spaces-properties-lemma-space-locally-quasi-compact}
设 $S$ 为概形，$X$ 为 $S$ 上的代数空间。$|X|$ 的每个点都有一个
由拟紧开邻域组成的基本邻域系。特别地，$|X|$ 按“拓扑”定义
\ref{topology-definition-locally-quasi-compact} 的意义局部拟紧。
\end{lemma}

\begin{proof}
这形式地来自下述事实：存在概形 $U$ 以及拓扑空间的满、开、连续映射
$U \to |X|$。更精确地说，若 $u \in U$ 映到 $x \in |X|$，
则 $u$ 的仿射邻域之像给出 $x$ 的一个拟紧开基本邻域系。
\end{proof}







\section{特殊覆盖}
\label{spaces-properties-section-special-coverings}

\noindent
本节汇集若干关于代数空间存在满 \'etale 覆盖的直接引理。

\begin{lemma}
\label{spaces-properties-lemma-cover-by-union-affines}
设 $S$ 为概形，$X$ 为 $S$ 上的代数空间。存在满 \'etale 态射
$U \to X$，其中 $U$ 是仿射概形的不交并。还可以假设其中每个仿射概形
都映入 $S$ 的某个仿射开集。
\end{lemma}

\begin{proof}
设 $V \to X$ 为满 \'etale 态射。取 Zariski 开覆盖
$V = \bigcup_{i \in I} V_i$，使每个 $V_i$ 都映入 $S$ 的某个仿射开集。
令 $U = \coprod_{i \in I} V_i$，并取诱导态射 $U \to V \to X$。
它是 \'etale 且满的，因为它是函子的 \'etale、满、可表示变换之复合；
这里使用一般原理“空间”引理
\ref{spaces-lemma-composition-representable-transformations-property}，
以及“态射”引理 \ref{morphisms-lemma-composition-surjective} 和
\ref{morphisms-lemma-composition-etale}。
\end{proof}

\begin{lemma}
\label{spaces-properties-lemma-union-of-quasi-compact}
设 $S$ 为概形，$X$ 为 $S$ 上的代数空间。存在 Zariski 覆盖
$X = \bigcup X_i$，使每个代数空间 $X_i$ 都有一个由仿射概形给出的
满 \'etale 覆盖。还可以假设每个 $X_i$ 都映入 $S$ 的某个仿射开集。
\end{lemma}

\begin{proof}
由引理 \ref{spaces-properties-lemma-cover-by-union-affines}，可找到满 \'etale 态射
$U = \coprod U_i \to X$，其中 $U_i$ 仿射且映入 $S$ 的某个仿射开集。
令 $X_i \subset X$ 为 $X$ 的开子空间，使 $U_i \to X$ 经过满
\'etale 态射 $U_i \to X_i$ 分解；见引理
\ref{spaces-properties-lemma-etale-image-open}。由于 $U = \bigcup U_i$，
可知 $X = \bigcup X_i$。又因 $U_i \to X_i$ 是满射，
$X_i \to S$ 的像位于 $S$ 的某个仿射开集中。
\end{proof}

\begin{lemma}
\label{spaces-properties-lemma-quasi-compact-affine-cover}
设 $S$ 为概形，$X$ 为 $S$ 上的代数空间。则 $X$ 拟紧，当且仅当
存在满 \'etale 态射 $U \to X$，其中 $U$ 是仿射概形。
\end{lemma}

\begin{proof}
若存在满 \'etale 态射 $U \to X$，且 $U$ 仿射，则由定义
\ref{spaces-properties-definition-quasi-compact}，$X$ 拟紧。反之，若 $X$ 拟紧，
则 $|X|$ 拟紧。令 $U = \coprod_{i \in I} U_i$ 为仿射概形的不交并，
并带有满 \'etale 映射 $\varphi : U \to X$
（引理 \ref{spaces-properties-lemma-cover-by-union-affines}）。于是
$|X| = \bigcup \varphi(|U_i|)$。由拟紧性，存在有限子集
$i_1, \ldots, i_n$，使 $|X| = \bigcup \varphi(|U_{i_j}|)$。
因此 $U_{i_1} \cup \ldots \cup U_{i_n}$ 是带有到 $X$
的有限满态射的仿射概形。
\end{proof}

\noindent
下面的引理将被“代数空间的态射”一章中关于分离态射的讨论取代。

\begin{lemma}
\label{spaces-properties-lemma-separated-cover}
设 $S$ 为概形，$X$ 为 $S$ 上的代数空间。设 $U$ 是分离概形，
且 $U \to X$ 为 \'etale 态射。则 $U \to X$ 分离，并且
$R = U \times_X U$ 是分离概形。
\end{lemma}

\begin{proof}
令 $X' \subset X$ 为开子概形，使 $U \to X$ 经过满 \'etale 态射
$U \to X'$ 分解；见引理 \ref{spaces-properties-lemma-etale-image-open}。
若 $U \to X'$ 分离，则 $U \to X$ 也分离；见“空间”引理
\ref{spaces-lemma-composition-representable-transformations-property}。
这里，开浸入 $X' \to X$ 分离，由“空间”引理
\ref{spaces-lemma-representable-transformations-property-implication}
及“概形”引理 \ref{schemes-lemma-immersions-monomorphisms} 可知。
此外，由于 $U \times_{X'} U = U \times_X U$，只需用 $X'$ 替换 $X$
后证明结论；换言之，可以假设 $U \to X$ 满。考虑交换图表
$$
\xymatrix{
R = U \times_X U \ar[r] \ar[d] & U \ar[d] \\
U \ar[r] & X
}
$$
在“空间”引理 \ref{spaces-lemma-properties-diagonal} 的证明中，
我们已看到 $j : R \to U \times_S U$ 分离。由于 $U$ 是分离概形，
概形态射 $U \to S$ 分离；见“概形”引理
\ref{schemes-lemma-compose-after-separated}。因此，作为基变换，
$U \times_S U \to U$ 分离；见“概形”引理
\ref{schemes-lemma-separated-permanence}。故概形 $U \times_S U$ 分离
（仍由该引理）。又因 $j$ 分离，以同样方式可知 $R$ 分离。
于是 $R \to U$ 是分离态射（再次使用“概形”引理
\ref{schemes-lemma-compose-after-separated}）。因此，由“空间”引理
\ref{spaces-lemma-representable-morphisms-spaces-property} 与上图，
可知 $U \to X$ 分离。
\end{proof}

\begin{lemma}
\label{spaces-properties-lemma-quasi-separated}
设 $S$ 为概形，$X$ 为 $S$ 上的代数空间。若存在拟分离概形 $U$
以及满 \'etale 态射 $U \to X$，使两个投影
$U \times_X U \to U$ 中的任一个拟紧，则 $X$ 拟分离。
\end{lemma}

\begin{proof}
可把 $X$ 视作 $\mathbf{Z}$ 上的代数空间。考虑 Cartesian 图表
$$
\xymatrix{
U \times_X U \ar[r] \ar[d]_j & X \ar[d]^\Delta \\
U \times U \ar[r] & X \times X
}
$$
由于 $U$ 拟分离，投影 $U \times U \to U$ 拟分离
（它是概形的拟分离态射的基变换；见“概形”引理
\ref{schemes-lemma-separated-permanence}）。因此，由本引理的假设与
“概形”引理 \ref{schemes-lemma-quasi-compact-permanence}，$j$ 拟紧。
再由“空间”引理 \ref{spaces-lemma-representable-morphisms-spaces-property}，
得到所需的 $\Delta$ 拟紧。
\end{proof}

\begin{lemma}
\label{spaces-properties-lemma-quasi-separated-quasi-compact-pieces}
设 $S$ 为概形，$X$ 为 $S$ 上的代数空间。下列条件等价：
\begin{enumerate}
\item $X$ 在 $S$ 上 Zariski 局部拟分离；
\item $X$ Zariski 局部拟分离；
\item 存在 Zariski 开覆盖 $X = \bigcup X_i$，使对每个 $i$，
都存在仿射概形 $U_i$ 以及拟紧的满 \'etale 态射 $U_i \to X_i$；
\item 存在 Zariski 开覆盖 $X = \bigcup X_i$，使对每个 $i$，
都存在映入 $S$ 的某个仿射开集的仿射概形 $U_i$，
以及拟紧的满 \'etale 态射 $U_i \to X_i$。
\end{enumerate}
\end{lemma}

\begin{proof}
假设 $U_i \to X_i \subset X$ 如 (3) 所述。为证明 (4)，
对每个 $i$ 选取有限仿射开覆盖
$U_i = U_{i1} \cup \ldots \cup U_{in_i}$，使每个 $U_{ij}$
都映入 $S$ 的某个仿射开集。复合
$U_{ij} \to U_i \to X_i$ 是 \'etale 且拟紧的
（见“空间”引理
\ref{spaces-lemma-composition-representable-transformations-property}）。
令 $X_{ij} \subset X_i$ 为与 $|U_{ij}| \to |X_i|$ 的像相对应的开子空间；
见引理 \ref{spaces-properties-lemma-etale-image-open}。注意，$X_{ij} \subset X_i$
是单态射，而 $U_{ij} \to X$ 拟紧，故 $U_{ij} \to X_{ij}$ 拟紧。
于是 $X = \bigcup X_{ij}$ 是 (4) 所述的覆盖。
(4) $\Rightarrow$ (3) 是立即的。

\medskip\noindent
假设 (4)。为证明 $X$ 在 $S$ 上 Zariski 局部拟分离，
只需证明 $X_i$ 在 $S$ 上拟分离。因此可假设存在映入 $S$
某个仿射开集的仿射概形 $U$，以及拟紧的满 \'etale 态射
$U \to X$。考虑纤维积方块
$$
\xymatrix{
U \times_X U \ar[r] \ar[d] & U \times_S U \ar[d] \\
X \ar[r]^-{\Delta_{X/S}} & X \times_S X
}
$$
右侧竖箭头是满 \'etale 态射（见“空间”引理
\ref{spaces-lemma-product-representable-transformations-property}）。
又因 $U$ 映入 $S$ 的某个仿射开集，$U \times_S U$ 仿射
（见“概形”第 \ref{schemes-section-fibre-products} 节）。
此外，投影 $U \times_X U \to U$ 是 $U \to X$ 的基变换，因而拟紧；
故 $U \times_X U$ 拟紧。由“空间”引理
\ref{spaces-lemma-representable-morphisms-spaces-property}，
得到所需的 $\Delta_{X/S}$ 拟紧。

\medskip\noindent
假设 (1)。为证明 (3)，可立即归约到 $X$ 在 $S$ 上拟分离的情形。
由引理 \ref{spaces-properties-lemma-union-of-quasi-compact}，可找到 Zariski 开覆盖
$X = \bigcup X_i$，使每个 $X_i$ 都映入 $S$ 的某个仿射开集，
并且存在仿射概形 $U_i$ 以及满 \'etale 态射 $U_i \to X_i$。
由于 $U_i \to S$ 映入 $S$ 的某个仿射开集，
$U_i \times_S U_i$ 仿射；见“概形”第
\ref{schemes-section-fibre-products} 节。由于 $X$ 在 $S$ 上拟分离，态射
$$
R_i = U_i \times_{X_i} U_i = U_i \times_X U_i
\longrightarrow
U_i \times_S U_i
$$
作为 $\Delta_{X/S}$ 的基变换均拟紧。因此 $R_i$ 是拟紧概形，
进而每个投影 $R_i \to U_i$ 拟紧。于是，对覆盖 $U_i \to X_i$
及态射 $U_i \to X_i$ 应用“空间”引理
\ref{spaces-lemma-representable-morphisms-spaces-property}，
可知态射 $U_i \to X_i$ 按所需拟紧。

\medskip\noindent
至此，(1)、(3) 与 (4) 等价。由于 (3) 不涉及基概形，
可知它们也与 (2) 等价。
\end{proof}

\noindent
下面的引理将十分有用。

\begin{lemma}
\label{spaces-properties-lemma-finite-fibres-presentation}
设 $S$ 为概形，$X$ 为 $S$ 上的代数空间，$U$ 为概形。
设 $\varphi : U \to X$ 为 \'etale 态射，并且投影
$R = U \times_X U \to U$ 拟紧；例如，$\varphi$ 拟紧时即如此。
则下列映射的纤维
$$
|U| \to |X|
\quad\text{且}\quad
|R| \to |X|
$$
都是有限集。
\end{lemma}

\begin{proof}
记 $R = U \times_X U$，并以 $s, t : R \to U$ 表示两个投影。
设 $u \in U$ 为点，$x \in |X|$ 为其像。由引理
\ref{spaces-properties-lemma-points-cartesian}，$|U| \to |X|$ 在 $x$ 上的纤维等于
$s(t^{-1}(\{u\}))$，而 $|R| \to |X|$ 在 $x$ 上的纤维是
$t^{-1}(s(t^{-1}(\{u\})))$。由于 $t : R \to U$ 是拟紧
\'etale 态射，其纤维有限；事实上，由“态射”引理
\ref{morphisms-lemma-etale-over-field}，其纤维是域谱的不交并，
并且拟紧。结论得证。
\end{proof}








\section{由概形性质定义的空间性质}
\label{spaces-properties-section-types-properties}

\noindent
概形的任何 \'etale 局部性质都通过下述引理给出代数空间的相应性质。

\begin{lemma}
\label{spaces-properties-lemma-type-property}
设 $S$ 为概形，$X$ 为 $S$ 上的代数空间。设 $\mathcal{P}$
是对 \'etale 拓扑为局部的概形性质；见“下降”定义
\ref{descent-definition-property-local}。下列条件等价：
\begin{enumerate}
\item 存在概形 $U$ 及满 \'etale 态射 $U \to X$，使概形 $U$
具有性质 $\mathcal{P}$；
\item 对每个概形 $U$ 及每个 \'etale 态射 $U \to X$，概形 $U$
都具有性质 $\mathcal{P}$。
\end{enumerate}
若 $X$ 可表示，则这些条件等价于 $\mathcal{P}(X)$。
\end{lemma}

\begin{proof}
(2) $\Rightarrow$ (1) 是立即的。反之，选取满 \'etale 态射
$U \to X$，其中 $U$ 是具有 $\mathcal{P}$ 的概形，并设 $V$
为 \'etale $X$-概形。于是 $U \times_X V \rightarrow V$
是概形的满 \'etale 态射，故 $V$ 从 $U \times_X V$ 继承
$\mathcal{P}$，后者又从 $U$ 继承 $\mathcal{P}$（见“下降”定义
\ref{descent-definition-property-local} 之后的讨论）。最后的断言由
(1) 及“下降”定义 \ref{descent-definition-property-local} 即得。
\end{proof}

\begin{definition}
\label{spaces-properties-definition-type-property}
设 $\mathcal{P}$ 为对 \'etale 拓扑局部的概形性质。设 $S$ 为概形，
$X$ 为 $S$ 上的代数空间。若引理 \ref{spaces-properties-lemma-type-property}
中的任一等价条件成立，则称 $X$ {\it 具有性质 $\mathcal{P}$}。
\end{definition}

\begin{remark}
\label{spaces-properties-remark-list-properties-local-etale-topology}
下列性质对 \'etale 拓扑是局部的（应注意，fpqc、fppf、syntomic
及光滑拓扑都强于 \'etale 拓扑）：
\begin{enumerate}
\item 局部 Noether；见“下降”引理
\ref{descent-lemma-Noetherian-local-fppf}；
\item Jacobson；见“下降”引理
\ref{descent-lemma-Jacobson-local-fppf}；
\item 局部 Noether 且满足 $(S_k)$；见“下降”引理
\ref{descent-lemma-Sk-local-syntomic}；
\item Cohen--Macaulay；见“下降”引理
\ref{descent-lemma-CM-local-syntomic}；
\item Gorenstein；见“概形上的对偶性”引理
\ref{duality-lemma-gorenstein-local-syntomic}；
\item 既约；见“下降”引理
\ref{descent-lemma-reduced-local-smooth}；
\item 正规；见“下降”引理
\ref{descent-lemma-normal-local-smooth}；
\item 局部 Noether 且满足 $(R_k)$；见“下降”引理
\ref{descent-lemma-Rk-local-smooth}；
\item 正则；见“下降”引理
\ref{descent-lemma-regular-local-smooth}；
\item Nagata；见“下降”引理
\ref{descent-lemma-Nagata-local-smooth}。
\end{enumerate}
\end{remark}

\noindent
概形芽的任何 \'etale 局部性质都给出代数空间的相应性质。下面是必需的引理。

\begin{lemma}
\label{spaces-properties-lemma-local-source-target-at-point}
设 $\mathcal{P}$ 是概形芽的 \'etale 局部性质；见“下降”定义
\ref{descent-definition-local-at-point}。设 $S$ 为概形，$X$ 为
$S$ 上的代数空间，$x \in |X|$ 为 $X$ 的一点。考虑以概形 $U$
为源的 \'etale 态射 $a : U \to X$。下列条件等价：
\begin{enumerate}
\item 对上述任意 $U \to X$ 及满足 $a(u) = x$ 的 $u \in U$，
都有 $\mathcal{P}(U, u)$；
\item 对上述某个 $U \to X$ 及某个满足 $a(u) = x$ 的 $u \in U$，
有 $\mathcal{P}(U, u)$。
\end{enumerate}
若 $X$ 可表示，则这些条件等价于 $\mathcal{P}(X, x)$。
\end{lemma}

\begin{proof}
从略。
\end{proof}

\begin{definition}
\label{spaces-properties-definition-property-at-point}
设 $S$ 为概形，$X$ 为 $S$ 上的代数空间，$x \in |X|$。设
$\mathcal{P}$ 为概形芽的 \'etale 局部性质。若引理
\ref{spaces-properties-lemma-local-source-target-at-point} 中任一等价条件成立，
则称 $X$ {\it 在 $x$ 处具有性质 $\mathcal{P}$}。
\end{definition}

\begin{remark}
\label{spaces-properties-remark-list-properties-local-ring-local-etale-topology}
设 $P$ 为局部环的性质。假设对任意 \'etale 环映射 $A \to B$，
以及 $B$ 中位于 $A$ 的素理想 $\mathfrak p$ 之上的素理想
$\mathfrak q$，都有
$P(A_\mathfrak p) \Leftrightarrow P(B_\mathfrak q)$.
规定 $\mathcal{P}(U, u) = P(\mathcal{O}_{U, u})$，便得到概形芽
$(U, u)$ 的一个 \'etale 局部性质。在这种情形下，我们用“$X$
在 $x$ 处的局部环具有 $P$”表示 $X$ 在 $x$ 处具有性质
$\mathcal{P}$。下列各项都是这样的性质 $P$：
\begin{enumerate}
\item Noether；见“代数进阶”引理
\ref{more-algebra-lemma-Noetherian-etale-extension}；
\item 维数为 $d$；见“代数进阶”引理
\ref{more-algebra-lemma-dimension-etale-extension}；
\item 正则；见“代数进阶”引理
\ref{more-algebra-lemma-regular-etale-extension}；
\item 离散赋值环；这由 (2)、(3) 及“代数”引理
\ref{algebra-lemma-characterize-dvr} 得出；
\item 既约；见“代数进阶”引理
\ref{more-algebra-lemma-henselization-reduced}；
\item 正规；见“代数进阶”引理
\ref{more-algebra-lemma-henselization-normal}；
\item Noether 且深度为 $k$；见“代数进阶”引理
\ref{more-algebra-lemma-henselization-depth}；
\item Noether 且 Cohen--Macaulay；见“代数进阶”引理
\ref{more-algebra-lemma-henselization-CM}；
\item Noether 且 Gorenstein；见“对偶复形”引理
\ref{dualizing-lemma-flat-under-gorenstein}。
\end{enumerate}
还有更多满足这一条件的性质，例如 G-环与 Nagata。若我们以后需要，
会将它们以及相应代数事实之详细证明的参考文献一并加入这里。
\end{remark}







\section{可构造集}
\label{spaces-properties-section-constructible}

\begin{lemma}
\label{spaces-properties-lemma-locally-constructible}
设 $S$ 为概形，$X$ 为 $S$ 上的代数空间，$E \subset |X|$
为子集。下列条件等价：
\begin{enumerate}
\item 对源 $U$ 为概形的每个 \'etale 态射 $U \to X$，$E$
在 $U$ 中的逆像都是 $U$ 的局部可构造子集；
\item 对源 $U$ 为仿射概形的每个 \'etale 态射 $U \to X$，$E$
在 $U$ 中的逆像都是 $U$ 的可构造子集；
\item 存在源 $U$ 为概形的满 \'etale 态射 $U \to X$，使 $E$
在 $U$ 中的逆像为 $U$ 的局部可构造子集。
\end{enumerate}
\end{lemma}

\begin{proof}
由“性质”引理 \ref{properties-lemma-locally-constructible}，
(1) 与 (2) 等价。(1) 蕴含 (3) 是立即的。因此，设我们有满
\'etale 态射 $\varphi : U \to X$，其中 $U$ 为概形且
$\varphi^{-1}(E)$ 局部可构造。再设
$\varphi' : U' \to X$ 为另一个以概形 $U'$ 为源的 \'etale 态射。
于是
$$
E'' = \text{pr}_1^{-1}(\varphi^{-1}(E)) = \text{pr}_2^{-1}((\varphi')^{-1}(E))
$$
其中 $\text{pr}_1 : U \times_X U' \to U$ 与
$\text{pr}_2 : U \times_X U' \to U'$ 为投影。由“态射”引理
\ref{morphisms-lemma-inverse-image-constructible}，$E''$
在 $U \times_X U'$ 中局部可构造。设 $W' \subset U'$ 为仿射开集。
由于 $\text{pr}_2$ 是 \'etale 的，因而是开映射，故可选取拟紧开集
$W'' \subset U \times_X U'$，使 $\text{pr}_2(W'') = W'$。
于是 $\text{pr}_2|_{W''} : W'' \to W'$ 拟紧。因为 $\varphi$
满，由引理 \ref{spaces-properties-lemma-points-cartesian}，
$W' \cap (\varphi')^{-1}(E) = \text{pr}_2(E'' \cap W'')$。
因此再由“态射”定理 \ref{morphisms-theorem-chevalley}，
$W' \cap (\varphi')^{-1}(E) = \text{pr}_2(E'' \cap W'')$
局部可构造，正合所需。
\end{proof}

\begin{definition}
\label{spaces-properties-definition-locally-constructible}
设 $S$ 为概形，$X$ 为 $S$ 上的代数空间，$E \subset |X|$
为子集。若引理 \ref{spaces-properties-lemma-locally-constructible} 中的等价条件成立，
则称 $E$ {\it \'etale 局部可构造}。
\end{definition}

\noindent
当然，若 $X$ 可表示，即 $X$ 是概形，则这恰好表示 $E$
是其底拓扑空间的局部可构造子集。







\section{一点处的维数}
\label{spaces-properties-section-dimension}

\noindent
可用“下降”引理 \ref{descent-lemma-dimension-at-point-local}
定义代数空间 $X$ 在一点 $x$ 处的维数。所得概念不同于拓扑意义下的概念
（即 $|X|$ 在 $x$ 处的维数）。

\begin{definition}
\label{spaces-properties-definition-dimension-at-point}
设 $S$ 为概形，$X$ 为 $S$ 上的代数空间，$x \in |X|$
为 $X$ 的一点。定义 {\it $X$ 在 $x$ 处的维数}为元素
$\dim_x(X) \in \{0, 1, 2, \ldots, \infty\}$，它满足：
对任意（等价地，某个）二元组 $(a : U \to X, u)$，都有
$\dim_x(X) = \dim_u(U)$；这里 $a : U \to X$ 是从概形到 $X$
的 \'etale 态射，$u \in U$ 且 $a(u) = x$。见定义
\ref{spaces-properties-definition-property-at-point}、引理
\ref{spaces-properties-lemma-local-source-target-at-point} 及“下降”引理
\ref{descent-lemma-dimension-at-point-local}。
\end{definition}

\noindent
警告：一般而言，{\bf 并非} $\dim_x(X) = \dim_x(|X|)$。
“空间”例 \ref{spaces-example-infinite-product} 中的代数空间 $X$
给出一个反例。具体地，设 $x \in |X|$ 是不等于 $|X|$ 的一般点
$x_0$ 的一点。此时 $\dim_x(X) = 0$，但 $\dim_x(|X|) = 1$。
特别地，$X$ 的维数（定义如下）不同于 $|X|$ 的维数。

\begin{definition}
\label{spaces-properties-definition-dimension}
设 $S$ 为概形，$X$ 为 $S$ 上的代数空间。$X$ 的{\it 维数}
$\dim(X)$ 定义为
$$
\dim(X) = \sup\nolimits_{x \in |X|} \dim_x(X)
$$
\end{definition}

\noindent
由“性质”引理 \ref{properties-lemma-dimension}，若 $X$ 为概形，
这就是通常的概念。还有另一个整数度量概形在一点处的维数，即局部环的维数。
这一不变量也与 \'etale 态射相容；见第
\ref{spaces-properties-section-dimension-local-ring} 节。






\section{局部环的维数}
\label{spaces-properties-section-dimension-local-ring}

\noindent
代数空间的局部环之维数是一个良定义的概念。

\begin{lemma}
\label{spaces-properties-lemma-pre-dimension-local-ring}
设 $S$ 为概形，$X$ 为 $S$ 上的代数空间，$x \in |X|$ 为一点，
$d \in \{0, 1, 2, \ldots, \infty\}$。下列条件等价：
\begin{enumerate}
\item 存在概形 $U$、\'etale 态射 $a : U \to X$ 及满足
$a(u) = x$ 的点 $u \in U$，使 $\dim(\mathcal{O}_{U, u}) = d$；
\item 对任意概形 $U$、任意 \'etale 态射 $a : U \to X$
及任意满足 $a(u) = x$ 的点 $u \in U$，都有
$\dim(\mathcal{O}_{U, u}) = d$。
\end{enumerate}
若 $X$ 为概形，则这些条件等价于 $\dim(\mathcal{O}_{X, x}) = d$。
\end{lemma}

\begin{proof}
结合引理 \ref{spaces-properties-lemma-local-source-target-at-point} 与“下降”引理
\ref{descent-lemma-dimension-local-ring-local}。
\end{proof}

\begin{definition}
\label{spaces-properties-definition-dimension-local-ring}
设 $S$ 为概形，$X$ 为 $S$ 上的代数空间，$x \in |X|$ 为一点。
{\it $X$ 在 $x$ 处之局部环的维数}是满足引理
\ref{spaces-properties-lemma-pre-dimension-local-ring} 中等价条件的元素
$d \in \{0, 1, 2, \ldots, \infty\}$。在这种情形下，我们也称
{\it $x$ 是 $X$ 上余维为 $d$ 的点}。
\end{definition}

\noindent
除下述引理外，还请读者参见引理 \ref{spaces-properties-lemma-dimension-local-ring}
与 \ref{spaces-properties-lemma-dimension-decent-invariant-under-etale}。

\begin{lemma}
\label{spaces-properties-lemma-dimension}
设 $S$ 为概形，$X$ 为 $S$ 上的代数空间。下列量相等：
\begin{enumerate}
\item $X$ 的维数；
\item $X$ 各局部环之维数的上确界；
\item 当 $x \in |X|$ 时，$\dim_x(X)$ 的上确界。
\end{enumerate}
\end{lemma}

\begin{proof}
由定义 \ref{spaces-properties-definition-dimension}，(1) 与 (3) 中的数相等。
设 $U \to X$ 为以概形 $U$ 为源的满 \'etale 态射。由定义，
当 $x \in |X|$ 时 $\dim_x(X)$ 的上确界，等于当 $u$ 遍历 $U$
各点时 $\dim_u(U)$ 的上确界。由“性质”引理
\ref{properties-lemma-dimension}，后者又等于
$\dim(\mathcal{O}_{U, u})$ 的上确界；而由定义，这就是 (2)。
\end{proof}




\section{一般点}
\label{spaces-properties-section-generic-points}

\noindent
设 $T$ 为拓扑空间。依照 EGA I 第二版，{\it $T$ 的极大点}是
$T$ 的某个不可约分支的一般点。若 $T = |X|$ 是与代数空间 $X$
相伴的拓扑空间，则至少有两种极大点概念：可以把 $T$ 看作拓扑空间
并取其极大点，也可以取 $U$ 的极大点在 \'etale 态射
$U \to X$ 下的像，其中 $U$ 是概形。第二种概念对应于余维 $0$ 的点集（引理
\ref{spaces-properties-lemma-codimension-0-points}）。对于一般的代数空间，
余维 $0$ 的点更便于使用；对拟分离代数空间，以及更一般的良好代数空间，
这两个概念一致（“良好空间”引理
\ref{decent-spaces-lemma-decent-generic-points}）。

\begin{lemma}
\label{spaces-properties-lemma-codimension-0-points}
设 $S$ 为概形，$X$ 为 $S$ 上的代数空间，$x \in |X|$。
考虑以概形 $U$ 为源的 \'etale 态射 $a : U \to X$。下列条件等价：
\begin{enumerate}
\item $x$ 是 $X$ 上余维为 $0$ 的点；
\item 对上述某个 $U \to X$ 及满足 $a(u) = x$ 的某个 $u \in U$，
点 $u$ 是 $U$ 的某个不可约分支的一般点；
\item 对上述任意 $U \to X$ 及映到 $x$ 的任意 $u \in U$，
点 $u$ 都是 $U$ 的某个不可约分支的一般点。
\end{enumerate}
若 $X$ 可表示，则这些条件等价于 $x$ 是 $|X|$
某个不可约分支的一般点。
\end{lemma}

\begin{proof}
注意，概形 $U$ 的一点 $u$ 是 $U$ 某个不可约分支的一般点，当且仅当
$\dim(\mathcal{O}_{U, u}) = 0$（“性质”引理
\ref{properties-lemma-generic-point}）。故结论由 $X$
上一点之余维的定义（定义 \ref{spaces-properties-definition-dimension-local-ring}）得出。
\end{proof}

\begin{lemma}
\label{spaces-properties-lemma-codimension-0-points-dense}
设 $S$ 为概形，$X$ 为 $S$ 上的代数空间。$X$ 的余维 $0$ 点集
在 $|X|$ 中稠密。
\end{lemma}

\begin{proof}
若 $U$ 为概形，则其各不可约分支之一般点的集合在 $U$ 中稠密
（这对任意拟清醒拓扑空间都成立）。因此，若 $U \to X$
为满 \'etale 态射，则 $X$ 的余维 $0$ 点集是 $|U|$ 中某个稠密子集的像
（引理 \ref{spaces-properties-lemma-codimension-0-points}）。由于 $|X|$
对 $|U| \to |X|$ 带商拓扑，结论得证。
\end{proof}





\section{既约空间}
\label{spaces-properties-section-reduced}

\noindent
我们已在第 \ref{spaces-properties-section-types-properties} 节定义了既约代数空间。
这里仅证明关于既约代数空间的几个简单引理。

\begin{lemma}
\label{spaces-properties-lemma-subspace-induced-topology}
设 $S$ 为概形，$Z \to X$ 为代数空间的浸入。则 $|Z| \to |X|$
是 $|Z|$ 到 $|X|$ 的某个局部闭子集上的同胚。
\end{lemma}

\begin{proof}
设 $U$ 为概形，$U \to X$ 为满 \'etale 态射。于是
$Z \times_X U \to U$ 是概形的浸入，因而给出 $|Z \times_X U|$
到 $|U|$ 的某个局部闭子集 $T'$ 上的同胚。由引理
\ref{spaces-properties-lemma-points-cartesian}，子集 $T'$ 是 $|Z| \to |X|$ 之像
$T$ 的逆像。映射 $|Z| \to |X|$ 为单射，因为函子变换
$Z \to X$ 是单射；见“空间”第 \ref{spaces-section-Zariski} 节。
由“拓扑”引理 \ref{topology-lemma-open-morphism-quotient-topology}，
$T$ 在 $|X|$ 中局部闭。此外，连续映射 $|Z| \to T$ 是同胚，
因为 $|Z \times_X U| \to T'$ 是同胚，且
$|Z \times_Y U| \to |Z|$ 是商映射。
\end{proof}

\noindent
下述引理将帮助我们构造（局部）闭子空间。

\begin{lemma}
\label{spaces-properties-lemma-subspaces-presentation}
设 $S$ 为概形，$j : R \to U \times_S U$ 为 \'etale 等价关系，
$X = U/R$ 为相伴的代数空间（“空间”定理
\ref{spaces-theorem-presentation}）。存在典范双射
$$
R\text{-不变的局部闭子概形 }Z'\text{，所在概形为 }U
\leftrightarrow
\text{局部闭子空间 }Z\text{，所在空间为 }X
$$
此外，$Z \to X$ 为闭（相应地，开）浸入，当且仅当
$Z' \to U$ 为闭（相应地，开）浸入。
\end{lemma}

\begin{proof}
以 $\varphi : U \to X$ 表示典范映射。该双射把 $Z \to X$
送到 $Z' = Z \times_X U \to U$。由定义立即可知，若 $Z \to X$
为浸入、闭浸入或开浸入，则 $Z' \to U$ 分别也是浸入、闭浸入或开浸入。
$Z'$ 为 $R$-不变的也是清楚的（见“群胚”定义
\ref{groupoids-definition-invariant-open}）。

\medskip\noindent
反之，设 $Z' \to U$ 为 $R$-不变的浸入。令 $R'$ 为 $R$
在 $Z'$ 上的限制；见“群胚”定义
\ref{groupoids-definition-restrict-groupoid}。在本情形中
$R' = R \times_{s, U} Z' = Z' \times_{U, t} R$，故 $R'$
是 $Z'$ 上的 \'etale 等价关系。由“空间”定理
\ref{spaces-theorem-presentation}，$Z = Z'/R'$ 是代数空间。
按构造有 $U \times_X Z = Z'$，故 $U \times_X Z \to Z$ 是浸入。
注意，性质“浸入”在基变换下保持，并且在基底上对 fppf 拓扑是局部的
（见“空间”第 \ref{spaces-section-lists} 节）。此外，浸入是分离且局部
拟有限的（见“概形”引理 \ref{schemes-lemma-immersions-monomorphisms}
与“态射”引理 \ref{morphisms-lemma-immersion-locally-quasi-finite}）。
故由“态射进阶”引理
\ref{more-morphisms-lemma-separated-locally-quasi-finite-morphisms-fppf-descend}，
浸入对 fppf 覆盖满足下降。这意味着对 $Z \to X$、
$\mathcal{P}=$“浸入”以及满 \'etale 态射 $U \to X$，“空间”引理
\ref{spaces-lemma-morphism-sheaves-with-P-effective-descent-etale}
的全部假设均成立。因此 $Z \to X$ 可表示且为浸入；这正是子空间的定义
（见“空间”定义 \ref{spaces-definition-immersion}）。

\medskip\noindent
显然，这两个构造互逆，结论得证。
\end{proof}

\begin{lemma}
\label{spaces-properties-lemma-reduced-closed-subspace}
设 $S$ 为概形，$X$ 为 $S$ 上的代数空间，$T \subset |X|$
为闭子集。存在唯一的闭子空间 $Z \subset X$，满足：
(a) $|Z| = T$；(b) $Z$ 既约。
\end{lemma}

\begin{proof}
设 $U \to X$ 为满 \'etale 态射，其中 $U$ 为概形。令
$R = U \times_X U$，从而 $X = U/R$；见“空间”引理
\ref{spaces-lemma-space-presentation}。照例以 $s, t : R \to U$
表示两个投影态射。由引理 \ref{spaces-properties-lemma-points-presentation}，
$T$ 对应于闭子集 $T' \subset |U|$，且
$s^{-1}(T') = t^{-1}(T')$。令 $Z' \subset U$ 为 $T'$
上诱导的既约概形结构。在此情形中，纤维积
$Z' \times_{U, t} R$ 与 $Z' \times_{U, s} R$ 是 $R$ 的闭子概形
（“概形”引理 \ref{schemes-lemma-base-change-immersion}），
并且在 $Z'$ 上 \'etale（“态射”引理
\ref{morphisms-lemma-base-change-etale}），故为既约的
（因为既约性对 \'etale 拓扑是局部的；见注
\ref{spaces-properties-remark-list-properties-local-etale-topology}）。由于二者具有相同的
底拓扑空间（见上），可知
$Z' \times_{U, t} R = Z' \times_{U, s} R$。于是可应用引理
\ref{spaces-properties-lemma-subspaces-presentation}，得到闭子空间 $Z \subset X$，
其到 $U$ 的拉回为 $Z'$。按构造，$|Z| = T$ 且 $Z$ 既约，
这证明了存在性。唯一性的证明从略。
\end{proof}

\begin{lemma}
\label{spaces-properties-lemma-map-into-reduction}
设 $S$ 为概形，$X$、$Y$ 为 $S$ 上的代数空间，$Z \subset X$
为闭子空间，并假设 $Y$ 既约。态射 $f : Y \to X$ 通过 $Z$
分解，当且仅当 $f(|Y|) \subset |Z|$。
\end{lemma}

\begin{proof}
假设 $f(|Y|) \subset |Z|$。选取图表
$$
\xymatrix{
V \ar[d]_b \ar[r]_h & U \ar[d]^a \\
Y \ar[r]^f & X
}
$$
其中 $U$、$V$ 为概形，两个竖直箭头均满且 \'etale。由引理
\ref{spaces-properties-lemma-type-property}，概形 $V$ 既约。因此，由“概形”引理
\ref{schemes-lemma-map-into-reduction}，$h$ 通过 $a^{-1}(Z)$ 分解，
故 $a \circ h$ 通过 $Z$ 分解。由于 $Z \subset X$ 是子层，而
$V \to Y$ 是 $(\Sch/S)_{fppf}$ 上层的满射，可知
$X \to Y$ 通过 $Z$ 分解。
\end{proof}

\begin{definition}
\label{spaces-properties-definition-reduced-induced-space}
设 $S$ 为概形，$X$ 为 $S$ 上的代数空间，$Z \subset |X|$
为闭子集。$Z$ 上的一个{\it 代数空间结构}由 $X$ 的闭子空间
$Z'$ 给出，其中 $|Z'|$ 等于 $Z$。$Z$ 上的{\it 诱导既约代数空间结构}
是引理 \ref{spaces-properties-lemma-reduced-closed-subspace} 中构造的结构。
{\it $X$ 的既约化 $X_{red}$}是 $|X|$ 上的诱导既约代数空间结构。
\end{definition}











\section{概形轨迹}
\label{spaces-properties-section-schematic}

\noindent
每个代数空间都有一个最大的、为概形的开子空间；这多少是清楚的，
但我们仍在下面写出证明。当然，这个子空间可能为空，例如
$X = \mathbf{A}^1_{\mathbf{Q}}/\mathbf{Z}$（普适反例）。
另一方面，例如当 $X$ 拟分离时，这个最大开子概形事实上在 $X$ 中稠密！

\begin{lemma}
\label{spaces-properties-lemma-subscheme}
设 $S$ 为概形，$X$ 为 $S$ 上的代数空间。存在一个最大的开子空间
$X' \subset X$，它是概形。
\end{lemma}

\begin{proof}
设 $U \to X$ 为满 \'etale 态射，其中 $U$ 为概形，并令
$R = U \times_X U$。$X$ 的开子空间与 $U$ 的 $R$-不变开子概形
按 $1 - 1$ 对应，因而它们组成一个集合。设 $X_i$（$i \in I$）为 $X$
中所有为概形（即可表示）的开子空间。考虑开子空间 $X' \subset X$，
其底点集为 $|X|$ 的开集 $\bigcup |X_i|$。由引理
\ref{spaces-properties-lemma-characterize-surjective}，
$$
\coprod X_i \longrightarrow X'
$$
是 $(\Sch/S)_{fppf}$ 上层的满射。但由于每个 $X_i \to X'$
均可由开浸入表示，该映射事实上在 Zariski 拓扑中满。具体地，
若 $T \to X'$ 是从概形到 $X'$ 的态射，则
$X_i \times_{X'} T$ 是 $T$ 的开子概形。因此可应用“概形”引理
\ref{schemes-lemma-glue-functors}，得出 $X'$ 是概形。
\end{proof}

\noindent
在本节余下部分中，若代数空间 $X$ 的开子空间 $X'$ 所对应的开子集
$|X'| \subset |X|$ 稠密，则称该开子空间{\it 稠密}。

\begin{lemma}
\label{spaces-properties-lemma-quasi-separated-finite-etale-cover-dense-open-scheme}
设 $S$ 为概形，$X$ 为 $S$ 上的代数空间。若存在有限、\'etale
且满的态射 $U \to X$，其中 $U$ 为拟分离概形，则 $X$
存在一个为概形的稠密开子空间 $X'$。更精确地，$X$
中每个余维为 $0$ 的点 $x \in |X|$ 都包含于 $X'$。
\end{lemma}

\begin{proof}
设 $X' \subset X$ 为最大的、是概形的开子空间（引理
\ref{spaces-properties-lemma-subscheme}）。设 $x \in |X|$ 为 $X$ 上余维为 $0$
的一点。由引理 \ref{spaces-properties-lemma-codimension-0-points-dense}，
只需证明 $x \in X'$。取引理陈述中的 $U \to X$。记
$R = U \times_X U$，并照例以 $s, t : R \to U$ 表示投影。
注意，$s,t$ 满、有限且 \'etale。由引理
\ref{spaces-properties-lemma-finite-fibres-presentation}，$|U| \to |X|$
在 $x$ 上的纤维有限，记为 $\{\eta_1, \ldots, \eta_n\}$。
由引理 \ref{spaces-properties-lemma-codimension-0-points}，每个 $\eta_i$
都是 $U$ 某个不可约分支的一般点。由“性质”引理
\ref{properties-lemma-maximal-points-affine}，可找到包含
$\{\eta_1, \ldots, \eta_n\}$
的仿射开集 $W \subset U$（这里用到了 $U$ 拟分离）。由“群胚”引理
\ref{groupoids-lemma-find-invariant-affine}，可假设 $W$ 为
$R$-不变。由于 $W \subset U$ 是 $R$-不变仿射开集，$R$
在 $W$ 上的限制 $R_W$ 等于
$R_W = s^{-1}(W) = t^{-1}(W)$（见“群胚”定义
\ref{groupoids-definition-invariant-open} 及其后的讨论）。
特别地，映射 $R_W \to W$ 也有限且 \'etale，故 $R_W$ 仿射。
于是由“群胚”命题 \ref{groupoids-proposition-finite-flat-equivalence}，
$W/R_W$ 是概形。另一方面，由“空间”引理
\ref{spaces-lemma-finding-opens}，$W/R_W$ 是 $X$ 的开子空间，
并且按构造包含 $x$。
\end{proof}

\noindent
我们将在“良好空间”定理
\ref{decent-spaces-theorem-decent-open-dense-scheme}
中把下述命题推广到良好代数空间。

\begin{proposition}
\label{spaces-properties-proposition-locally-quasi-separated-open-dense-scheme}
设 $S$ 为概形，$X$ 为 $S$ 上的代数空间。若 $X$ 是 Zariski
局部拟分离的（例如 $X$ 拟分离），则 $X$ 存在为概形的稠密开子空间
$X'$。更精确地，$X$ 上每个余维为 $0$ 的点 $x \in |X|$
都包含于 $X'$。
\end{proposition}

\begin{proof}
由引理 \ref{spaces-properties-lemma-subscheme}，该问题在 $X$ 上是局部的。
故由引理 \ref{spaces-properties-lemma-quasi-separated-quasi-compact-pieces}，
可假设存在仿射概形 $U$ 及满、拟紧、\'etale 态射 $U \to X$。
此外，$U \to X$ 分离（引理 \ref{spaces-properties-lemma-separated-cover}）。
令 $R = U \times_X U$，并照例以 $s, t : R \to U$ 表示投影。
则 $s,t$ 满、拟紧、分离且 \'etale。因此 $s,t$ 也拟有限且具有有限纤维
（“态射”引理 \ref{morphisms-lemma-etale-locally-quasi-finite}、
\ref{morphisms-lemma-quasi-finite-locally-quasi-compact} 及
\ref{morphisms-lemma-quasi-finite}）。
由“态射”引理 \ref{morphisms-lemma-generically-finite}，对作为 $U$
某个不可约分支一般点的每个 $\eta \in U$，存在 $\eta$ 的开邻域
$V \subset U$，使 $s^{-1}(V) \to V$ 有限。由“下降”引理
\ref{descent-lemma-descending-property-finite}，有限性在目标上是
fpqc 局部的（特别是 \'etale 局部的）。因此可应用“群胚进阶”引理
\ref{more-groupoids-lemma-property-invariant}：使 $s$ 有限的最大开集
$W \subset U$ 是 $R$-不变的。由上可知，$W$ 包含 $U$
每个不可约分支的一般点。$R$ 在 $W$ 上的限制 $R_W$ 等于
$R_W = s^{-1}(W) = t^{-1}(W)$（见“群胚”定义
\ref{groupoids-definition-invariant-open} 及其后的讨论）。
按构造，$s_W, t_W : R_W \to W$ 有限且 \'etale。考虑开子空间
$X' = W/R_W \subset X$（见“空间”引理
\ref{spaces-lemma-finding-opens}）。按构造，包含映射 $X' \to X$
在余维 $0$ 的点上诱导双射。于是问题归结为引理
\ref{spaces-properties-lemma-quasi-separated-finite-etale-cover-dense-open-scheme}。
\end{proof}




\section{得到概形}
\label{spaces-properties-section-getting-a-scheme}

\noindent
上一节中我们用到了如下事实：若态射 $s, t : R \to U$ 有限且
\'etale，则仿射概形 $U$ 对等价关系 $R$ 的商 $U/R$ 是概形。
这是下述结果的特殊情形。

\begin{proposition}
\label{spaces-properties-proposition-finite-flat-equivalence-global}
设 $S$ 为概形，$(U, R, s, t, c)$ 为 $S$ 上的群胚概形。假设：
\begin{enumerate}
\item $s, t : R \to U$ 有限局部自由；
\item $j = (t, s)$ 是等价关系；
\item $U$ 的每个非空闭子集 $Z$ 都包含一点 $u$，其 $R$-等价类
$t(s^{-1}(\{u\}))$ 包含于 $U$ 的某个仿射开集\footnote{令
$E \subset U$ 为满足如下条件的点 $u$ 之集：
$t(s^{-1}(\{u\}))$ 包含于 $U$ 的某个仿射开集。若 $E = U$，
或 $U$ 的每个有限型点都属于 $E$，或每个 $u \in U$
都特化到 $E$ 的一点，则条件 (3) 成立。}。
\end{enumerate}
则存在 $S$ 上概形之间的有限局部自由态射 $U \to M$，使得
$R = U \times_M U$，并且 $M$ 表示 fppf 拓扑中的商层 $U/R$。
\end{proposition}

\begin{proof}
由假设 (3) 与“群胚”引理 \ref{groupoids-lemma-find-invariant-affine}，
可找到开覆盖 $U = \bigcup U_i$，使每个 $U_i$ 都是 $U$
的 $R$-不变仿射开集。令 $R_i = R|_{U_i}$。考虑 fppf 层
$F = U/R$ 与 $F_i = U_i/R_i$。由“空间”引理
\ref{spaces-lemma-finding-opens}，态射 $F_i \to F$ 可表示且为开浸入。
由“群胚”命题 \ref{groupoids-proposition-finite-flat-equivalence}，
层 $F_i$ 可由仿射概形表示。若 $T$ 为概形且 $T \to F$ 为态射，
则 $V_i = F_i \times_F T$ 在 $T$ 中开；我们断言
$T = \bigcup V_i$。事实上，在 $T$ 上 fppf 局部地，可将
$T \to F$ 提升为态射 $f : T \to U$，此时
$f^{-1}(U_i) \subset V_i$。因此 $F$ 可由概形表示；见“概形”引理
\ref{schemes-lemma-glue-functors}。
\end{proof}

\noindent
例如，若 $U$ 同构于某个仿射概形的局部闭子概形，或同构于某个分次环
$A$ 的 $\text{Proj}(A)$ 的局部闭子概形，则由“性质”引理
\ref{properties-lemma-ample-finite-set-in-affine}，第三个假设成立。
特别地，这可应用于有限群及有限群概形在拟仿射或拟投射概形上的自由作用。
例如，有限群 $G$ 在拟投射簇 $X$ 上自由作用所得的商 $X/G$
是概形。下面给出详细陈述。

\begin{lemma}
\label{spaces-properties-lemma-quotient-scheme}
设 $S$ 为概形，$G \to S$ 为群概形，$X \to S$ 为概形的态射，
$a : G \times_S X \to X$ 为一个作用。假设：
\begin{enumerate}
\item $G \to S$ 有限局部自由；
\item 作用 $a$ 自由；
\item $X \to S$ 仿射、拟仿射、投射或拟投射，或者 $X$
同构于某个仿射概形的开子概形，或者 $X$ 同构于某个分次环 $A$
之 $\text{Proj}(A)$ 的开子概形，或者 $G \to S$ 是 radicial 的。
\end{enumerate}
则 fppf 商层 $X/G$ 是概形，且 $X \to X/G$ 是 fppf $G$-挠子。
\end{lemma}

\begin{proof}
先证明 $X/G$ 是概形。因为作用自由，态射
$j = (a, \text{pr}) : G \times_S X \to X \times_S X$
是单态射，因而是等价关系；见“群胚”引理
\ref{groupoids-lemma-free-action}。由于假设 $G \to S$ 有限局部自由，
映射 $s, t : G \times_S X \to X$ 有限局部自由。现在只需证明命题
\ref{spaces-properties-proposition-finite-flat-equivalence-global} 的最后一个假设成立。
由于 $G$ 的作用在 $S$ 上，只需证明 $X \to S$ 的任一纤维中的有限点集
都包含于 $X$ 的某个仿射开集。若 $X$ 同构于仿射概形的开子概形，
或同构于某个分次环 $A$ 之 $\text{Proj}(A)$ 的开子概形，
则这由“性质”引理 \ref{properties-lemma-ample-finite-set-in-affine}
得出。若 $X \to S$ 仿射、拟仿射、投射或拟投射，可将 $S$
换为仿射开集，从而回到刚处理的情形。若 $G \to S$ 是 radicial 的，
则 $G$ 作用在 $X$ 各点上的轨道均为单点集，条件显然成立。
若干细节从略。

\medskip\noindent
为证明 $X \to X/G$ 是 fppf $G$-挠子（“群胚”定义
\ref{groupoids-definition-principal-homogeneous-space}），须证明
$G \times_S X \to X \times_{X/G} X$ 是同构，且 $X \to X/G$
在 fppf 局部有截面。第二点由 $X \to X/G$ 作为 fppf 层的映射为满
（按构造）立即得出。第一点来自命题
\ref{spaces-properties-proposition-finite-flat-equivalence-global} 结论中的同构
$R = U \times_M U$（注意，在本情形中 $R = G \times_S X$）。
\end{proof}

\begin{lemma}
\label{spaces-properties-lemma-quotient-separated}
记号与假设同命题 \ref{spaces-properties-proposition-finite-flat-equivalence-global}。则：
\begin{enumerate}
\item 若 $U$ 在 $S$ 上拟分离，则 $U/R$ 在 $S$ 上拟分离；
\item 若 $U$ 拟分离，则 $U/R$ 拟分离；
\item 若 $U$ 在 $S$ 上分离，则 $U/R$ 在 $S$ 上分离；
\item 若 $U$ 分离，则 $U/R$ 分离；
\item 在此添加更多内容。
\end{enumerate}
在引理 \ref{spaces-properties-lemma-quotient-scheme} 的情形中也有类似结果。
\end{lemma}

\begin{proof}
由于 $M$ 表示商层，有概形的 Cartesian 图表
$$
\xymatrix{
R \ar[r]_-j \ar[d] & U \times_S U \ar[d] \\
M \ar[r] & M \times_S M
}
$$
由于 $U \times_S U \to M \times_S M$ 满且有限局部自由，为证明
$M \to M \times_S M$ 拟紧（相应地，为闭浸入），只需证明
$j : R \to U \times_S U$ 拟紧（相应地，为闭浸入）；见“下降”引理
\ref{descent-lemma-descending-property-quasi-compact} 与
\ref{descent-lemma-descending-property-closed-immersion}。
由于 $j : R \to U \times_S U$ 是 $U$ 上的态射，而 $R$ 在 $U$
上有限，只要投影 $U \times_S U \to U$ 拟分离，$j$ 便拟紧
（“概形”引理 \ref{schemes-lemma-quasi-compact-permanence}）。
由于 $j$ 为单态射且局部有限型，只要它固有，$j$ 便是闭浸入
（“\'Etale 态射”引理 \ref{etale-lemma-characterize-closed-immersion}）；
而只要投影 $U \times_S U \to U$ 分离，这一条件便成立
（“态射”引理 \ref{morphisms-lemma-image-proper-scheme-closed}）。
这证明了 (1) 与 (3)。为证明 (2) 与 (4)，将 $S$
换为 $\Spec(\mathbf{Z})$；见定义 \ref{spaces-properties-definition-separated}。
由于引理 \ref{spaces-properties-lemma-quotient-scheme} 是通过应用命题
\ref{spaces-properties-proposition-finite-flat-equivalence-global} 证明的，
最后的断言也清楚。
\end{proof}




\section{拟分离空间上的点}
\label{spaces-properties-section-points-quasi-separated}

\noindent
在 Stacks 项目所采用的一般性下，代数空间上的点可能表现得很差。
然而，对拟分离空间，其行为大多与概形上的点相似。本节证明几个相关结果；
“良好空间”一章还包含更多结果，例如见“良好空间”第
\ref{decent-spaces-section-points} 节。

\begin{lemma}
\label{spaces-properties-lemma-quasi-separated-sober}
设 $S$ 为概形，$X$ 为 $S$ 上 Zariski 局部拟分离的代数空间。
则拓扑空间 $|X|$ 是清醒的（见“拓扑”定义
\ref{topology-definition-generic-point}）。
\end{lemma}

\begin{proof}
结合“拓扑”引理 \ref{topology-lemma-sober-local} 与引理
\ref{spaces-properties-lemma-quasi-separated-quasi-compact-pieces}，可假设存在仿射概形
$U$ 及满、拟紧、\'etale 态射 $U \to X$。令
$R = U \times_X U$，投影映射为 $s, t : R \to U$。应用引理
\ref{spaces-properties-lemma-finite-fibres-presentation}，可知 $s,t$ 的纤维有限。
于是“拓扑”引理 \ref{topology-lemma-quotient-kolmogorov}
的全部假设均满足，从而 $|X|$ 是 Kolmogorov 空间\footnote{
事实上，这里还用到了“概形”引理
\ref{schemes-lemma-scheme-sober}（概形的清醒性）、
“态射”引理 \ref{morphisms-lemma-etale-flat} 与
\ref{morphisms-lemma-generalizations-lift-flat}
（沿 \'etale 态射可提升一般化）、引理
\ref{spaces-properties-lemma-points-presentation}（用表示描述代数空间上的点），以及引理
\ref{spaces-properties-lemma-topology-points}（商映射的开性）。}。

\medskip\noindent
还需证明每个不可约闭子集 $T \subset |X|$ 都有一般点。由引理
\ref{spaces-properties-lemma-reduced-closed-subspace}，存在闭子空间 $Z \subset X$，
使 $|Z| = |T|$。注意，$U \times_X Z \to Z$ 是从仿射概形到 $Z$
的拟紧、满、\'etale 态射，故由引理
\ref{spaces-properties-lemma-quasi-separated-quasi-compact-pieces}，$Z$ 是 Zariski
局部拟分离的。由命题
\ref{spaces-properties-proposition-locally-quasi-separated-open-dense-scheme}，
存在为概形的开稠密子空间 $Z' \subset Z$。这意味着
$|Z'| \subset T$ 开且稠密。因此拓扑空间 $|Z'|$ 不可约，
即 $Z'$ 是不可约概形。由“概形”引理
\ref{schemes-lemma-scheme-sober}，$|Z'|$ 是单个点
$\eta \in |Z'| \subset T$ 的闭包，因而也有
$T = \overline{\{\eta\}}$，结论得证。
\end{proof}

\begin{lemma}
\label{spaces-properties-lemma-quasi-compact-quasi-separated-spectral}
设 $S$ 为概形，$X$ 为 $S$ 上拟紧且拟分离的代数空间。
拓扑空间 $|X|$ 是谱空间。
\end{lemma}

\begin{proof}
由“拓扑”定义 \ref{topology-definition-spectral-space}，须验证
$|X|$ 清醒、拟紧、具有由拟紧开集组成的基，并且任意两个拟紧开集
之交拟紧。由引理 \ref{spaces-properties-lemma-quasi-separated-sober}，$|X|$ 清醒；
由引理 \ref{spaces-properties-lemma-quasi-compact-space}，$|X|$ 拟紧。
由引理 \ref{spaces-properties-lemma-quasi-compact-affine-cover}，存在仿射概形 $U$
及满 \'etale 态射 $f : U \to X$。因为
$|f| : |U| \to |X|$ 开且连续，而 $|U|$ 具有由拟紧开集组成的基，
故 $|X|$ 也有这样的基。最后，设 $A, B \subset |X|$ 为拟紧开集。
则存在开子空间 $X', X'' \subset X$，使
$A = |X'|$ 且 $B = |X''|$（引理 \ref{spaces-properties-lemma-open-subspaces}）；
并且可选取仿射概形 $V$、$W$ 及满 \'etale 态射
$V \to X'$、$W \to X''$（引理
\ref{spaces-properties-lemma-quasi-compact-affine-cover}）。于是 $A \cap B$
是 $|V \times_X W| \to |X|$ 的像（引理
\ref{spaces-properties-lemma-points-cartesian}）。由于 $X$ 拟分离，
$V \times_X W$ 拟紧（引理
\ref{spaces-properties-lemma-characterize-quasi-separated}），故 $A \cap B$ 拟紧，
证明完成。
\end{proof}

\noindent
下述引理可用于证明一个代数空间同构于某个域的谱。

\begin{lemma}
\label{spaces-properties-lemma-point-like-spaces}
设 $S$ 为概形，$k$ 为域，$X$ 为 $S$ 上的代数空间，并假设存在
满 \'etale 态射 $\Spec(k) \to X$。若 $X$ 拟分离，则
$X \cong \Spec(k')$，其中 $k/k'$ 为有限可分扩张。
\end{lemma}

\begin{proof}
令 $R = \Spec(k) \times_X \Spec(k)$，于是有纤维积图表
$$
\xymatrix{
R \ar[r]_-s \ar[d]_-t & \Spec(k) \ar[d] \\
\Spec(k) \ar[r] & X
}
$$
由“空间”引理 \ref{spaces-lemma-space-presentation}，
$X = \Spec(k)/R$ 是商层。由于 $\Spec(k) \to X$ 为 \'etale，
态射 $s$ 与 $t$ 均为 \'etale。因此
$R = \coprod_{i \in I} \Spec(k_i)$ 是域谱的不交并，并且 $s$ 与 $t$
都诱导有限可分域扩张 $s, t : k \subset k_i$；见“态射”引理
\ref{morphisms-lemma-etale-over-field}。因为
$$
R = \Spec(k) \times_X \Spec(k)
= (\Spec(k) \times_S \Spec(k)) \times_{X \times_S X, \Delta} X
$$
并且按假设 $\Delta$ 拟紧，可知
$R \to \Spec(k) \times_S \Spec(k)$ 拟紧。由于
$\Spec(k) \times_S \Spec(k)$ 仿射，$R$ 拟紧，故 $I$ 有限。
这蕴含 $s$ 与 $t$ 是有限局部自由态射。于是由“群胚”命题
\ref{groupoids-proposition-finite-flat-equivalence}，
$\Spec(k)/R$ 可由 $\Spec(k')$ 表示，其中 $k' \subset k$
有限局部自由，且
$$
k' = \{x \in k \mid s_i(x) = t_i(x)\text{ 对所有 }i \in I\}
$$
容易看出 $k'$ 是域。
\end{proof}

\begin{remark}
\label{spaces-properties-remark-cannot-decide-yet}
引理 \ref{spaces-properties-lemma-point-like-spaces} 对良好代数空间也成立；见“良好空间”
引理 \ref{decent-spaces-lemma-decent-point-like-spaces}。事实上，
只有一个点的良好代数空间是概形；见“良好空间”引理
\ref{decent-spaces-lemma-when-field}。当 $X$ 局部分离时，这也成立，
因为局部分离的代数空间是良好的；见“良好空间”引理
\ref{decent-spaces-lemma-locally-separated-decent}。
\end{remark}







\section{代数空间的 \'Etale 态射}
\label{spaces-properties-section-etale-morphisms}

\noindent
本节本应属于“代数空间的态射”一章，但为了定义代数空间的小 \'etale
站点，需要一个代数空间在另一个之上 \'etale 的概念。因此必须先对从概形
到代数空间的 \'etale 态射以及代数空间之间的 \'etale 态射做一些准备。
关于代数空间的 \'etale 态射的更多内容，见“空间的态射”第
\ref{spaces-morphisms-section-etale} 节。

\begin{lemma}
\label{spaces-properties-lemma-etale-over-space}
设 $S$ 为概形，$X$ 为 $S$ 上的代数空间，$U$、$U'$ 为 $S$ 上的概形。
\begin{enumerate}
\item 若 $U \to U'$ 是概形的 \'etale 态射，而 $U' \to X$
是从 $U'$ 到 $X$ 的 \'etale 态射，则复合 $U \to X$
是从 $U$ 到 $X$ 的 \'etale 态射。
\item 若 $\varphi : U \to X$ 与 $\varphi' : U' \to X$
是到 $X$ 的 \'etale 态射，而概形态射 $\chi : U \to U'$
满足 $\varphi = \varphi' \circ \chi$，则 $\chi$ 是概形的
\'etale 态射。
\item 若 $\chi : U \to U'$ 是概形的满 \'etale 态射，而
$\varphi' : U' \to X$ 是满足 $\varphi = \varphi' \circ \chi$
为 \'etale 的态射，则 $\varphi'$ 为 \'etale。
\end{enumerate}
\end{lemma}

\begin{proof}
回忆：从概形到代数空间的任意态射均可表示；据此，我们从“空间”定义
\ref{spaces-definition-relative-representable-property}
得到从概形到代数空间的 \'etale 态射之定义。

\medskip\noindent
引理的 (1) 由此、\'etale 态射在复合下保持这一事实
（“态射”引理 \ref{morphisms-lemma-composition-etale}）以及“空间”引理
\ref{spaces-lemma-composition-representable-transformations-property} 与
\ref{spaces-lemma-morphism-schemes-gives-representable-transformation-property}
（它们是形式的）得出。

\medskip\noindent
为证明 (2)，选取 $S$ 上的概形 $W$ 及满 \'etale 态射 $W \to X$。
考虑 $\chi$ 的基变换
$\chi_W : W \times_X U \to W \times_X U'$。由于
$W \times_X U$ 与 $W \times_X U'$ 在 $W$ 上 \'etale，
由“态射”引理 \ref{morphisms-lemma-etale-permanence}，
$\chi_W$ 为 \'etale。另一方面，在交换图表
$$
\xymatrix{
W \times_X U \ar[r] \ar[d] & W \times_X U' \ar[d] \\
U \ar[r] & U'
}
$$
中，两个竖直箭头都 \'etale 且满。故由“下降”引理
\ref{descent-lemma-syntomic-smooth-etale-permanence}，
$U \to U'$ 为 \'etale。

\medskip\noindent
为证明 (3)，选取 $S$ 上的概形 $W$ 及态射 $W \to X$。
如上考虑图表
$$
\xymatrix{
W \times_X U \ar[r] \ar[d] & W \times_X U' \ar[d] \ar[r] & W \ar[d] \\
U \ar[r] & U' \ar[r] & X
}
$$
现在知道 $W \times_X U \to W \times_X U'$ 满且 \'etale
（因为它是 $U \to U'$ 的基变换），而
$W \times_X U \to W$ 为 \'etale。因此，由“下降”引理
\ref{descent-lemma-syntomic-smooth-etale-permanence}，
$W \times_X U' \to W$ 为 \'etale。按定义，这意味着
$\varphi'$ 为 \'etale。
\end{proof}

\begin{definition}
\label{spaces-properties-definition-etale}
设 $S$ 为概形。称 $S$ 上代数空间之间的态射 $f : X \to Y$
为 {\it \'etale} 的，当且仅当对每个以概形 $U$ 为源的 \'etale
态射 $\varphi : U \to X$，复合 $f \circ \varphi$ 也为 \'etale。
\end{definition}

\noindent
若 $X$ 与 $Y$ 是概形，则这与概形的 \'etale 态射之通常概念一致。
事实上，只要 $X \to Y$ 是代数空间的可表示态射，这一概念便与通过
“空间”定义 \ref{spaces-definition-relative-representable-property}
给出的概念一致。这由下述引理 \ref{spaces-properties-lemma-etale-local} 与“空间”引理
\ref{spaces-lemma-representable-morphisms-spaces-property} 结合得出。

\begin{lemma}
\label{spaces-properties-lemma-etale-local}
设 $S$ 为概形，$f : X \to Y$ 为 $S$ 上代数空间的态射。
下列条件等价：
\begin{enumerate}
\item $f$ 为 \'etale；
\item 存在满 \'etale 态射 $\varphi : U \to X$，其中 $U$ 为概形，
使复合 $f \circ \varphi$ 为 \'etale（作为代数空间的态射）；
\item 存在满 \'etale 态射 $\psi : V \to Y$，其中 $V$ 为概形，
使基变换 $V \times_Y X \to V$ 为 \'etale（作为代数空间的态射）；
\item 存在交换图表
$$
\xymatrix{
U \ar[d] \ar[r] & V \ar[d] \\
X \ar[r] & Y
}
$$
其中 $U$、$V$ 为概形，两个竖直箭头为 \'etale，左竖直箭头满，
并且上方的水平箭头为 \'etale。
\end{enumerate}
\end{lemma}

\begin{proof}
先证明 (4) 蕴含 (1)。假设给定 (4) 中的图表。设 $W \to X$
为以概形 $W$ 为源的 \'etale 态射。则
$W \times_X U \to U$ 为 \'etale。由于
$W \times_X U \to U$ 与 $U \to V$ 均为概形的 \'etale 态射，
其复合 $W \times_X U \to V$ 为 \'etale。因此，由引理
\ref{spaces-properties-lemma-etale-over-space} (1)，$W \times_X U \to Y$ 为
\'etale。投影 $W \times_X U \to W$ 也满且 \'etale，故由引理
\ref{spaces-properties-lemma-etale-over-space} (3)，$W \to Y$ 为 \'etale。

\medskip\noindent
再证明 (1) 蕴含 (4)。假设 (1)，并选取交换图表
$$
\xymatrix{
U \ar[d] \ar[r] & V \ar[d] \\
X \ar[r] & Y
}
$$
其中 $U \to X$ 与 $V \to Y$ 满且 \'etale；见“空间”引理
\ref{spaces-lemma-lift-morphism-presentations}。按假设，
态射 $U \to Y$ 为 \'etale，故由引理
\ref{spaces-properties-lemma-etale-over-space} (2)，$U \to V$ 为 \'etale。

\medskip\noindent
(2) 与 (3) 也等价于 (1) 的证明从略。
\end{proof}

\begin{lemma}
\label{spaces-properties-lemma-composition-etale}
代数空间的两个 \'etale 态射之复合为 \'etale。
\end{lemma}

\begin{proof}
这由定义立即得出。
\end{proof}

\begin{lemma}
\label{spaces-properties-lemma-base-change-etale}
代数空间的 \'etale 态射沿代数空间的任意态射作基变换后仍为 \'etale。
\end{lemma}

\begin{proof}
设 $X \to Y$ 为 $S$ 上代数空间的 \'etale 态射，$Z \to Y$
为代数空间的态射。选取概形 $U$ 及满 \'etale 态射 $U \to X$；
再选取概形 $W$ 及满 \'etale 态射 $W \to Z$。于是 $U \to Y$
为 \'etale，故在图表
$$
\xymatrix{
W \times_Y U \ar[d] \ar[r] & W \ar[d] \\
Z \times_Y X \ar[r] & Z
}
$$
中，上方的水平箭头为 \'etale。此外，左侧竖直箭头满且 \'etale
（验证从略）。故由引理 \ref{spaces-properties-lemma-etale-local}，下方水平箭头为
\'etale。
\end{proof}

\begin{lemma}
\label{spaces-properties-lemma-etale-permanence}
设 $S$ 为概形，$X,Y,Z$ 为代数空间。设
$g : X \to Z$、$h : Y \to Z$ 为 \'etale 态射，而
$f : X \to Y$ 为满足 $h \circ f = g$ 的态射。则 $f$ 为 \'etale。
\end{lemma}

\begin{proof}
选取交换图表
$$
\xymatrix{
U \ar[d] \ar[r]_\chi & V \ar[d] \\
X \ar[r] & Y
}
$$
其中 $U \to X$ 与 $V \to Y$ 满且 \'etale；见“空间”引理
\ref{spaces-lemma-lift-morphism-presentations}。按假设，态射
$\varphi : U \to X \to Z$ 与 $\psi : V \to Y \to Z$ 为
\'etale。此外，由对 $f,g,h$ 的假设，
$\psi \circ \chi = \varphi$。故由引理
\ref{spaces-properties-lemma-etale-over-space} (2)，$U \to V$ 为 \'etale。
\end{proof}

\begin{lemma}
\label{spaces-properties-lemma-etale-open}
设 $S$ 为概形。若 $X \to Y$ 为 $S$ 上代数空间的 \'etale 态射，
则相伴的拓扑空间映射 $|X| \to |Y|$ 是开映射。
\end{lemma}

\begin{proof}
这由引理 \ref{spaces-properties-lemma-etale-local} 中的图表及引理
\ref{spaces-properties-lemma-topology-points} 即得。
\end{proof}

\noindent
最后给出一个有趣的引理。带有到概形之 \'etale 态射的代数空间不一定是概形；
见“空间”例 \ref{spaces-example-non-representable-descent}。
但若目标是某个域的谱，结论成立。

\begin{lemma}
\label{spaces-properties-lemma-etale-over-field-scheme}
设 $S$ 为概形，$X \to \Spec(k)$ 为 $S$ 上的 \'etale 态射，
其中 $k$ 为域。则 $X$ 是概形。
\end{lemma}

\begin{proof}
设 $U$ 为仿射概形，$U \to X$ 为 \'etale 态射。由定义
\ref{spaces-properties-definition-etale}，$U \to \Spec(k)$ 为 \'etale 态射。
故 $U = \coprod_{i = 1, \ldots, n} \Spec(k_i)$ 是有限个域谱的不交并，
其中 $k_i$ 为 $k$ 的有限可分扩张；见“态射”引理
\ref{morphisms-lemma-etale-over-field}。态射
$R = U \times_X U \to U \times_{\Spec(k)} U$ 为单态射，而
$U \times_{\Spec(k)} U$ 也为有限个 $k$ 的有限可分扩张之谱的不交并。
故由“概形”引理 \ref{schemes-lemma-mono-towards-spec-field}，$R$
同样是有限个 $k$ 的有限可分扩张之谱的不交并。这里 $U$ 与 $R$
均仿射，且两个投影 $R \to U$ 都有限局部自由。因此，由“群胚”命题
\ref{groupoids-proposition-finite-flat-equivalence}，$U/R$ 是概形。
由“空间”引理 \ref{spaces-lemma-finding-opens}，它也是 $X$
的开子空间。再由引理 \ref{spaces-properties-lemma-subscheme}，$X$ 是概形。
\end{proof}













\section{空间与 fpqc 覆盖}
\label{spaces-properties-section-fpqc}

\noindent
设 $S$ 为概形。$S$ 上的代数空间定义为 fppf 拓扑中满足若干附加性质的层。
因此，它是否满足 fpqc 拓扑的层性质并非立即清楚（见“拓扑”定义
\ref{topologies-definition-sheaf-property-fpqc}）。本节给出 Gabber
的论证，证明答案是肯定的。不过，当我们说代数空间 $X$ 满足 fpqc
拓扑的层性质时，实际上只考虑这样的 fpqc 覆盖
$\{f_i : T_i \to T\}_{i \in I}$：$T,T_i$ 是大站点
$(\Sch/S)_{fppf}$ 的对象（依照我们的约定，见第
\ref{spaces-properties-section-conventions} 节）。

\begin{proposition}[Gabber]
\label{spaces-properties-proposition-sheaf-fpqc}
设 $S$ 为概形，$X$ 为 $S$ 上的代数空间。则 $X$ 满足 fpqc
拓扑的层性质。
\end{proposition}

\begin{proof}
由于 $X$ 是 Zariski 拓扑的层，只需证明如下事实：给定仿射概形的
满平坦态射 $f : T' \to T$，$X(T)$ 是两个映射
$X(T') \to X(T' \times_T T')$ 的等化子。见“拓扑”引理
\ref{topologies-lemma-sheaf-property-fpqc}（这里略去了一个小论证，
因为所引引理是对全体概形范畴上定义的函子表述的）。

\medskip\noindent
设 $a, b : T \to X$ 为满足 $a \circ f = b \circ f$ 的两个态射。
须证明 $a = b$。考虑纤维积
$$
E = X \times_{\Delta_{X/S}, X \times_S X, (a, b)} T.
$$
由“空间”引理 \ref{spaces-lemma-properties-diagonal}，态射
$\Delta_{X/S}$ 是可表示单态射，故 $E \to T$ 是概形的单态射。
假设 $a \circ f = b \circ f$ 蕴含 $T' \to T$
（唯一地）通过 $E$ 分解。考虑交换图表
$$
\xymatrix{
T' \times_T E \ar[r] \ar[d] & E \ar[d] \\
T' \ar[r] \ar@/^5ex/[u] \ar[ru] & T
}
$$
投影 $T' \times_T E \to T'$ 是带截面的单态射，故为同构。
于是由“下降”引理 \ref{descent-lemma-descending-property-isomorphism}，
$E \to T$ 为同构。这意味着 $a = b$，正合所需。

\medskip\noindent
接着，设 $c : T' \to X$ 为一个态射，使得两个复合
$T' \times_T T' \to T' \to X$ 相同。须找到态射 $a : T \to X$，
使它与 $T' \to T$ 的复合为 $c$。选取仿射概形 $U$
及 \'etale 态射 $U \to X$，使 $|U| \to |X|$ 的像包含
$|c| : |T'| \to |X|$ 的像。这可由引理
\ref{spaces-properties-lemma-topology-points}、\ref{spaces-properties-lemma-cover-by-union-affines}、
有限个仿射概形的不交并仍仿射这一事实以及 $|T'|$ 拟紧这一事实做到
（略去一个小论证）。由于 $U \to X$ 分离（引理
\ref{spaces-properties-lemma-separated-cover}），
$$
V = U \times_{X, c} T' \longrightarrow T'
$$
是概形的满、\'etale、分离态射（要看出它满，使用引理
\ref{spaces-properties-lemma-points-cartesian} 及对 $U \to X$ 的选择）。
等式 $c \circ \text{pr}_0 = c \circ \text{pr}_1$ 意味着
$V/T'/T$ 上得到一个下降数据（“下降”定义
\ref{descent-definition-descent-datum}），因为
\begin{align*}
V \times_{T'} (T' \times_T T')
& =
U \times_{X, c \circ \text{pr}_0} (T' \times_T T') \\
& =
(T' \times_T T') \times_{c \circ \text{pr}_1, X} U \\
& =
(T' \times_T T') \times_{T'} V
\end{align*}
由“态射进阶”引理
\ref{more-morphisms-lemma-etale-separated-ind-quasi-affine}，
态射 $V \to T'$ 是 ind-拟仿射的（因为 \'etale 态射局部拟有限；
见“态射”引理 \ref{morphisms-lemma-etale-locally-quasi-finite}）。
由“群胚进阶”引理 \ref{more-groupoids-lemma-ind-quasi-affine}，
该下降数据有效。设 $W \to T$ 为态射，并有同构
$\alpha : T' \times_T W \to V$，它与 $V$ 上给定的下降数据以及
$T' \times_T W$ 上的典范下降数据相容。则 $W \to T$ 满且
\'etale（“下降”引理 \ref{descent-lemma-descending-property-surjective}
与 \ref{descent-lemma-descending-property-etale}）。考虑复合
$$
b' : T' \times_T W \longrightarrow V = U \times_{X, c} T' \longrightarrow U
$$
两个复合
$b' \circ (\text{pr}_0, 1), 
b' \circ (\text{pr}_1, 1) :
(T' \times_T T') \times_T W \to T' \times_T W \to U$
因对 $\alpha$ 的选择及 $c$ 的相应性质而相同（计算从略）。
故由“下降”引理
\ref{descent-lemma-fpqc-universal-effective-epimorphisms}，
$b'$ 下降为态射 $b : W \to U$。图表
$$
\xymatrix{
T' \times_T W \ar[r] \ar[d] & W \ar[r]_b & U \ar[d] \\
T' \ar[rr]^c &  & X
}
$$
交换。这意味着已在 $T$ 上 \'etale 局部地证明了 $a$ 的存在性，
即有 $a' : W \to X$。但由于第一段已证明唯一性，这一 \'etale
局部解满足粘合条件；也就是说，作为 $X(W \times_T W)$ 的元素，
有 $\text{pr}_0^*a' = \text{pr}_1^*a'$。由于 $X$ 是 \'etale 层，
存在唯一的 $a \in X(T)$，其到 $W$ 的限制为 $a'$。
\end{proof}









\section{代数空间的 \'etale 站点}
\label{spaces-properties-section-etale-site}

\noindent
本节定义代数空间的小 \'etale 站点。这是概形的小 \'etale 站点
$S_\etale$ 的类比。引理 \ref{spaces-properties-lemma-etale-over-space} 蕴含：
在下述定义中，$X$ 的 \'etale 站点之对象之间的任意态射都是
\'etale 的；在 $X_\etale$ 的某个对象上 \'etale 的任意概形也仍是
$X_\etale$ 的对象。

\begin{definition}
\label{spaces-properties-definition-etale-site}
设 $S$ 为概形，$\Sch_{fppf}$ 为包含 $S$ 的大 fppf 站点，
$\Sch_\etale$ 为相应的大 \'etale 站点（即有相同的底范畴）。
设 $X$ 为 $S$ 上的代数空间。$X$ 的{\it 小 \'etale 站点
$X_\etale$}定义如下：
\begin{enumerate}
\item $X_\etale$ 的对象是态射 $\varphi : U \to X$，其中
$U \in \Ob((\Sch/S)_\etale)$ 为概形，$\varphi$ 为 \'etale 态射；
\item 态射 $(\varphi : U \to X) \to (\varphi' : U' \to X)$
由满足 $\varphi = \varphi' \circ \chi$ 的概形态射
$\chi : U \to U'$ 给出；
\item $X_\etale$ 的态射族
$\{(U_i \to X) \to (U \to X)\}_{i \in I}$ 是覆盖，当且仅当
$\{U_i \to U\}_{i \in I}$ 是 $(\Sch/S)_\etale$ 的覆盖。
\end{enumerate}
\end{definition}

\noindent
这一选择的结果是：代数空间的 \'etale 站点一般没有终对象！另一方面，
若 $X$ 恰为概形，则上述定义与“拓扑”定义
\ref{topologies-definition-big-small-etale} 一致。

\medskip\noindent
上面是默认站点，但还会使用几个变体。具体地，可以考虑在 $X$ 上
\'etale 的所有{\it 代数空间} $U$，得到下述站点
$X_{spaces, \etale}$；也可以考虑在 $X$ 上 \'etale 的所有
{\it 仿射概形} $U$，得到下述站点 $X_{affine, \etale}$。
讨论小 \'etale 站点的函子性时会使用前一种概念；见引理
\ref{spaces-properties-lemma-functoriality-etale-site}。

\begin{definition}
\label{spaces-properties-definition-spaces-etale-site}
设 $S$ 为概形，$\Sch_{fppf}$ 为包含 $S$ 的大 fppf 站点，
$\Sch_\etale$ 为相应的大 \'etale 站点（即有相同的底范畴）。
设 $X$ 为 $S$ 上的代数空间。$X$ 的站点
{\it $X_{spaces, \etale}$} 定义如下：
\begin{enumerate}
\item $X_{spaces, \etale}$ 的对象是态射 $\varphi : U \to X$，
其中 $U$ 为 $S$ 上的代数空间，$\varphi$ 为 $S$
上代数空间的 \'etale 态射；
\item $X_{spaces, \etale}$ 的态射
$(\varphi : U \to X) \to (\varphi' : U' \to X)$
由满足 $\varphi = \varphi' \circ \chi$ 的代数空间态射
$\chi : U \to U'$ 给出；
\item $X_{spaces, \etale}$ 的态射族
$\{\varphi_i : (U_i \to X) \to (U \to X)\}_{i \in I}$ 是覆盖，
当且仅当 $|U| = \bigcup \varphi_i(|U_i|)$。
\end{enumerate}
照例，依“集合”引理 \ref{sets-lemma-coverings-site}，选取一组
此类覆盖，使其至少包含 $X_\etale$ 中的覆盖，从而把
$X_{spaces, \etale}$ 变为站点。
\end{definition}

\noindent
由于 $X$ 的恒等态射为 \'etale，显然 $X_{spaces, \etale}$
确有终对象。下面立即证明相应拓扑斯等于 $X$ 的小 \'etale 拓扑斯。

\begin{lemma}
\label{spaces-properties-lemma-compare-etale-sites}
函子
$$
X_\etale \longrightarrow X_{spaces, \etale}, \quad
U/X \longmapsto U/X
$$
是特殊余连续函子（“站点”定义
\ref{sites-definition-special-cocontinuous-functor}），因而诱导拓扑斯等价
$\Sh(X_\etale) \to \Sh(X_{spaces, \etale})$。
\end{lemma}

\begin{proof}
须证明该函子满足“站点”引理 \ref{sites-lemma-equivalence}
的假设 (1)--(5)。该函子连续且余连续是清楚的，这证明 (1) 与 (2)。
由于该函子全忠实，(3) 与 (4) 立即成立。按定义，代数空间有概形给出的覆盖，
故 (5) 成立。
\end{proof}

\begin{remark}
\label{spaces-properties-remark-explain-equivalence}
下面解释引理 \ref{spaces-properties-lemma-compare-etale-sites} 的含义。设 $S$ 为概形，
$X$ 为 $S$ 上的代数空间，$\mathcal{F}$ 为 $X$ 的小 \'etale
站点 $X_\etale$ 上的层。该引理说明：$X_{spaces, \etale}$
上存在唯一的层 $\mathcal{F}'$，其到子范畴 $X_\etale$
的限制为 $\mathcal{F}$。若 $U \to X$ 是代数空间的 \'etale 态射，
应如何计算 $\mathcal{F}'(U)$？由代数空间的定义，存在概形 $U'$
及满 \'etale 态射 $U' \to U$。于是 $\{U' \to U\}$ 是
$X_{spaces, \etale}$ 中的覆盖，故得到等化子图表
$$
\xymatrix{
\mathcal{F}'(U) \ar[r] &
\mathcal{F}(U') \ar@<1ex>[r] \ar@<-1ex>[r] &
\mathcal{F}(U' \times_U U').
}
$$
注意，$U' \times_U U'$ 是概形，因而这里可以写
$\mathcal{F}$ 而非 $\mathcal{F}'$。这便说明了给定层
$\mathcal{F}$ 时如何计算 $\mathcal{F}'$。
\end{remark}

\begin{definition}
\label{spaces-properties-definition-affine-etale-site}
设 $S$ 为概形，$\Sch_{fppf}$ 为包含 $S$ 的大 fppf 站点，
$\Sch_\etale$ 为相应的大 \'etale 站点（即有相同的底范畴）。
设 $X$ 为 $S$ 上的代数空间。$X$ 的站点
{\it $X_{affine, \etale}$} 定义如下：
\begin{enumerate}
\item $X_{affine, \etale}$ 的对象是态射 $\varphi : U \to X$，
其中 $U \in \Ob((\Sch/S)_\etale)$ 为仿射概形，$\varphi$
为 \'etale 态射；
\item $X_{affine, \etale}$ 的态射
$(\varphi : U \to X) \to (\varphi' : U' \to X)$
由满足 $\varphi = \varphi' \circ \chi$ 的概形态射
$\chi : U \to U'$ 给出；
\item $X_{affine, \etale}$ 的态射族
$\{\varphi_i : (U_i \to X) \to (U \to X)\}_{i \in I}$ 是覆盖，
当且仅当 $\{U_i \to U\}$ 是标准 \'etale 覆盖；见“拓扑”定义
\ref{topologies-definition-standard-etale}。
\end{enumerate}
照例，依“集合”引理 \ref{sets-lemma-coverings-site}，
选取一组此类覆盖，从而把 $X_{affine, \etale}$ 变为站点。
\end{definition}

\begin{lemma}
\label{spaces-properties-lemma-alternative}
设 $S$ 为概形，$X$ 为 $S$ 上的代数空间。函子
$X_{affine, \etale} \to X_\etale$ 特殊余连续，并诱导从
$\Sh(X_{affine, \etale})$ 到 $\Sh(X_\etale)$ 的拓扑斯等价。
\end{lemma}

\begin{proof}
从略。提示：与“拓扑”引理
\ref{topologies-lemma-affine-big-site-etale} 的证明比较。
\end{proof}

\begin{definition}
\label{spaces-properties-definition-etale-topos}
设 $S$ 为概形，$X$ 为 $S$ 上的代数空间。$X$ 的
{\it \'etale 拓扑斯}，更精确地说，$X$ 的{\it 小 \'etale 拓扑斯}，
是 $X_\etale$ 上集合层的范畴 $\Sh(X_\etale)$。
\end{definition}

\noindent
由引理 \ref{spaces-properties-lemma-compare-etale-sites}，
$\Sh(X_\etale) = \Sh(X_{spaces, \etale})$，故也可将其视为
$X_{spaces, \etale}$ 上集合层的范畴。类似地，由引理
\ref{spaces-properties-lemma-alternative}，
$\Sh(X_\etale) = \Sh(X_{affine, \etale})$。
事实证明，该拓扑斯关于代数空间的态射具有函子性。精确陈述如下。

\begin{lemma}
\label{spaces-properties-lemma-functoriality-etale-site}
设 $S$ 为概形，$f : X \to Y$ 为 $S$ 上代数空间的态射。
\begin{enumerate}
\item 连续函子
$$
Y_{spaces, \etale} \longrightarrow X_{spaces, \etale}, \quad
V \longmapsto X \times_Y V
$$
诱导站点态射
$$
f_{spaces, \etale} :
X_{spaces, \etale}
\to
Y_{spaces, \etale}.
$$
\item 规则 $f \mapsto f_{spaces, \etale}$ 与复合相容；换言之，
$(f \circ g)_{spaces, \etale}
= f_{spaces, \etale} \circ g_{spaces, \etale}$（见“站点”定义
\ref{sites-definition-composition-morphisms-sites}）。
\item 与 $f_{spaces, \etale}$ 相伴的拓扑斯态射通过引理
\ref{spaces-properties-lemma-compare-etale-sites} 诱导拓扑斯态射
$f_{small} : \Sh(X_\etale) \to \Sh(Y_\etale)$，其构造与复合相容。
\item 若 $f$ 是代数空间的可表示态射，则 $f_{small}$
来自站点态射 $X_\etale \to Y_\etale$，它对应于连续函子
$V \mapsto X \times_Y V$。
\end{enumerate}
\end{lemma}

\begin{proof}
先证明 (1) 中的函子满足“站点”命题
\ref{sites-proposition-get-morphism} 的假设。须证明
$Y_{spaces, \etale}$ 有终对象（即 $Y$），并且该函子把它变为
$X_{spaces, \etale}$ 中的终对象（即 $X$）。这很清楚，因为在任意范畴中
$X \times_Y Y = X$。接着须证明 $Y_{spaces, \etale}$ 有纤维积。
这是成立的，因为代数空间范畴有纤维积，并且若 $V$、$V'$
在 $Y$ 上 \'etale，则 $V \times_Y V'$ 也在 $Y$ 上 \'etale
（见上面的引理 \ref{spaces-properties-lemma-composition-etale} 与
\ref{spaces-properties-lemma-base-change-etale}）。因此该命题可用，并得到 (1)
所述站点态射。

\medskip\noindent
展开定义即得 (2)。利用上面引理
\ref{spaces-properties-lemma-compare-etale-sites} 对 $X$ 与 $Y$ 给出的等价，(3) 是清楚的。
(4) 也成立，因为若 $f$ 可表示，则上述函子组成交换图表
$$
\xymatrix{
X_\etale \ar[r] &
X_{spaces, \etale} \\
Y_\etale \ar[r] \ar[u] &
Y_{spaces, \etale} \ar[u]
}
$$
的范畴。
\end{proof}

\noindent
在描述 $X$ 上的层与 $Y$ 上的层之间的关系时，还可以比上述引理更进一步。
具体地，可用 $f$-映射来表述；比较“层”定义
\ref{sheaves-definition-f-map}。如下。

\begin{definition}
\label{spaces-properties-definition-f-map}
设 $S$ 为概形，$f : X \to Y$ 为 $S$ 上代数空间的态射，
$\mathcal{F}$ 为 $X_\etale$ 上的集合层，$\mathcal{G}$
为 $Y_\etale$ 上的集合层。一个
{\it $f$-映射 $\varphi : \mathcal{G} \to \mathcal{F}$}
是由交换图表
$$
\xymatrix{
U \ar[d]_g \ar[r] & X \ar[d]^f \\
V \ar[r] & Y
}
$$
指标化的一族映射
$\varphi_{(U, V, g)} : \mathcal{G}(V) \to \mathcal{F}(U)$，
其中 $U \in X_\etale$、$V \in Y_\etale$，并且只要给定扩充图表
$$
\xymatrix{
U' \ar[r] \ar[d]_{g'} & U \ar[d]_g \ar[r] & X \ar[d]^f \\
V' \ar[r] & V \ar[r] & Y
}
$$
其中 $V' \to V$ 与 $U' \to U$ 为概形的 \'etale 态射，图表
$$
\xymatrix{
\mathcal{G}(V)
\ar[rr]_{\varphi_{(U, V, g)}}
\ar[d]_{\text{限制 }\mathcal{G}} & &
\mathcal{F}(U)
\ar[d]^{\text{限制 }\mathcal{F}} \\
\mathcal{G}(V')
\ar[rr]^{\varphi_{(U', V', g')}} & &
\mathcal{F}(U')
}
$$
就交换。
\end{definition}

\begin{lemma}
\label{spaces-properties-lemma-f-map}
设 $S$ 为概形，$f : X \to Y$ 为 $S$ 上代数空间的态射，
$\mathcal{F}$ 为 $X_\etale$ 上的集合层，$\mathcal{G}$
为 $Y_\etale$ 上的集合层。下列三个集合之间存在典范双射：
\begin{enumerate}
\item 映射 $\mathcal{G} \to f_{small, *}\mathcal{F}$ 的集合；
\item 映射 $f_{small}^{-1}\mathcal{G} \to \mathcal{F}$ 的集合；
\item $f$-映射 $\varphi : \mathcal{G} \to \mathcal{F}$ 的集合。
\end{enumerate}
\end{lemma}

\begin{proof}
注意，(1) 与 (2) 相同，因为函子 $f_{small, *}$
与 $f_{small}^{-1}$ 互为伴随。设
$\alpha : f_{small}^{-1}\mathcal{G} \to \mathcal{F}$
为 $Y_\etale$ 上层的映射，并给定定义
\ref{spaces-properties-definition-f-map} 中的图表
$$
\xymatrix{
U \ar[d]_g \ar[r]_{j_U} & X \ar[d]^f \\
V \ar[r]^{j_V} & Y
}
$$
。由图表的交换性，还得到映射
$g_{small}^{-1}(j_V)^{-1}\mathcal{G} \to (j_U)^{-1}\mathcal{F}$
（关于局部化函子的描述，比较“站点”第
\ref{sites-section-localize} 节）。因此当然得到映射
$\varphi_{(V, U, g)} :
\mathcal{G}(V) = (j_V)^{-1}\mathcal{G}(V)
\to
(j_U)^{-1}\mathcal{F}(U) = \mathcal{F}(U)$.
这一规则与进一步的限制相容，并定义从 $\mathcal{G}$ 到
$\mathcal{F}$ 的 $f$-映射；其验证从略。

\medskip\noindent
反之，设给定 $f$-映射
$\varphi = (\varphi_{(U, V, g)})$。令 $\mathcal{G}'$
（相应地，$\mathcal{F}'$）表示 $\mathcal{G}$（相应地，
$\mathcal{F}$）到 $Y_{spaces, \etale}$（相应地，
$X_{spaces, \etale}$）的扩张；见引理
\ref{spaces-properties-lemma-compare-etale-sites}。于是须构造层的映射
$$
\mathcal{G}' \longrightarrow (f_{spaces, \etale})_*\mathcal{F}'
$$
为此，设 $V \to Y$ 为代数空间的 \'etale 态射。须构造集合的映射
$$
\mathcal{G}'(V) \to \mathcal{F}'(X \times_Y V)
$$
选取满 \'etale 态射 $V' \to V$，其中 $V'$ 为概形；随后选取
满 \'etale 态射 $U' \to X \times_U V'$，其中 $U'$ 为概形。
得到概形态射 $g' : U' \to V'$，以及概形态射
$$
g'' : U' \times_{X \times_Y V} U' \longrightarrow V' \times_V V'
$$
考虑图表
$$
\xymatrix{
\mathcal{F}'(X \times_Y V) \ar[r] &
\mathcal{F}(U') \ar@<1ex>[r] \ar@<-1ex>[r] &
\mathcal{F}(U' \times_{X \times_Y V} U') \\
\mathcal{G}'(X \times_Y V) \ar[r] \ar@{..>}[u] &
\mathcal{G}(V') \ar@<1ex>[r] \ar@<-1ex>[r] \ar[u]_{\varphi_{(U', V', g')}} &
\mathcal{G}(V' \times_V V') \ar[u]_{\varphi_{(U'', V'', g'')}}
}
$$
映射 $\varphi_{...}$ 与限制的相容性表明右侧两个方块交换。
$X_{spaces, \etale}$ 中覆盖的定义表明两条水平行都是等化子图表，
故得到虚线箭头。这些箭头与限制映射相容的证明留给读者。
\end{proof}

\noindent
若代数空间态射 $X \to Y$ 是 \'etale 的，则拓扑斯态射
$\Sh(X_\etale) \to \Sh(Y_\etale)$ 是一个局部化。精确地说：

\begin{lemma}
\label{spaces-properties-lemma-etale-morphism-topoi}
设 $S$ 为概形，$f : X \to Y$ 为 $S$ 上代数空间的态射。
假设 $f$ 是 \'etale 的。此时有函子
$$
j : X_\etale \to Y_\etale, \quad
(\varphi : U \to X) \mapsto (f \circ \varphi : U \to Y)
$$
它是余连续的。拓扑斯态射 $f_{small}$ 正是与 $j$ 相伴的拓扑斯态射，
参见“站点”引理 \ref{sites-lemma-cocontinuous-morphism-topoi}。
此外，$j$ 也是连续的，故可应用“站点”引理
\ref{sites-lemma-when-shriek}。特别地，对 $Y_\etale$ 上的每个层
$\mathcal{G}$，都有
$f_{small}^{-1}\mathcal{G}(U) = \mathcal{G}(jU)$。
\end{lemma}

\begin{proof}
由代数空间 \'etale 态射的定义本身（定义 \ref{spaces-properties-definition-etale}）可知，
所给规则确实定义了上述函子 $j$。$j$ 余连续且连续是清楚的，因为
$Y_\etale$ 中 $j(\varphi : U \to X)$ 的一个覆盖
$\{U_i \to U\}$，恰好就是 $X_\etale$ 中
$(\varphi : U \to X)$ 的一个覆盖。只须再证明 $j$ 诱导的拓扑斯态射
与 $f_{small}$ 相同。为此考虑范畴图
$$
\xymatrix{
X_\etale \ar[r] \ar[d]^j &
X_{spaces, \etale} \ar@/_/[d]_{j_{spaces}} \\
Y_\etale \ar[r] &
Y_{spaces, \etale} \ar@/_/[u]_{v : V \mapsto X \times_Y V}
}
$$
其中函子 $j_{spaces}$ 是 $j$ 到范畴 $X_{spaces, \etale}$ 的显然延拓。
因此内方块交换。事实上，$j_{spaces}$ 可与“站点”第
\ref{sites-section-localize} 节中讨论的局部化函子
$j_X : Y_{spaces, \etale}/X \to Y_{spaces, \etale}$
等同。于是由“站点”引理 \ref{sites-lemma-localize-given-products}，
余连续函子 $j_{spaces}$ 与图中的函子 $v$ 诱导同一个拓扑斯态射。
再由“站点”引理 \ref{sites-lemma-composition-cocontinuous}，
内方块（由站点之间的余连续函子组成）的交换性给出相伴拓扑斯态射的交换图。
最后根据引理 \ref{spaces-properties-lemma-functoriality-etale-site} 中 $f_{small}$ 的构造即得结论。
\end{proof}

\noindent
上面的引理说明，经由代数空间的 \'etale 态射 $f : X \to Y$
拉回 $\mathcal{G}$，不过是把 $\mathcal{G}$ 限制到范畴 $X_\etale$。
我们常用简记
\begin{equation}
\label{spaces-properties-equation-restrict}
\mathcal{G}|_{X_\etale} = f_{small}^{-1}\mathcal{G}
\end{equation}
来表示这一点。注意，在这种情形下，引理中的函子
$j : X_\etale \to Y_\etale$ 是忠实的，但一般并非全忠实。
第 \ref{spaces-properties-section-localize} 节将以更技术性的方式讨论此事。

\begin{lemma}
\label{spaces-properties-lemma-pushforward-etale-base-change}
设 $S$ 为概形，并设
$$
\xymatrix{
X' \ar[r] \ar[d]_{f'} & X \ar[d]^f \\
Y' \ar[r]^g & Y
}
$$
为 $S$ 上代数空间的笛卡儿方块。设 $\mathcal{F}$ 为 $X_\etale$ 上的层。
若 $g$ 是 \'etale 的，则
\begin{enumerate}
\item $f'_{small, *}(\mathcal{F}|_{X'}) = (f_{small, *}\mathcal{F})|_{Y'}$
在 $\Sh(Y'_\etale)$ 中成立\footnote{还可得
$(f')_{small}^{-1}(\mathcal{G}|_{Y'}) = (f_{small}^{-1}\mathcal{G})|_{X'}$
，这是因为该图交换且有式 (\ref{spaces-properties-equation-restrict})}；
\item 若 $\mathcal{F}$ 是阿贝尔层，则
$R^if'_{small, *}(\mathcal{F}|_{X'}) = (R^if_{small, *}\mathcal{F})|_{Y'}$.
\end{enumerate}
\end{lemma}

\begin{proof}
考虑如下函子图
$$
\xymatrix{
X'_{spaces, \etale} \ar[r]_j &
X_{spaces, \etale} \\
Y'_{spaces, \etale} \ar[r]^j \ar[u]^{V' \mapsto V' \times_{Y'} X'} &
Y_{spaces, \etale} \ar[u]_{V \mapsto V \times_Y X}
}
$$
横箭头均为局部化，而纵箭头诱导站点态射。因此“站点”引理
\ref{sites-lemma-localize-morphism} 的最后一个陈述给出 (1)。
为证明 (2)，把 (1) 应用于 $\mathcal{F}$ 的一个内射分解，并使用限制函子
正合且保持内射对象这一事实（见“站点上的上同调”引理
\ref{sites-cohomology-lemma-cohomology-of-open}）。
\end{proof}

\noindent
下面的引理表明，可把代数空间小 \'etale 站点上的一个层看作一族相容的层：
它们分别定义在该空间上 \'etale 的各概形的小 \'etale 站点上。
请注意，引理中的全部比较映射 $c_f$ 都是同构；这与“拓扑”引理
\ref{topologies-lemma-characterize-sheaf-big-etale} 以及
$X_\etale$ 的对象之间每个态射均为 \'etale 态射这一事实相符。

\begin{lemma}
\label{spaces-properties-lemma-characterize-sheaf-small-etale}
设 $S$ 为概形，$X$ 为 $S$ 上的代数空间。$X_\etale$ 上的一个层
$\mathcal{F}$ 由下列数据给出：
\begin{enumerate}
\item 对每个 $U \in \Ob(X_\etale)$，给定 $U_\etale$ 上的层
$\mathcal{F}_U$；
\item 对 $X_\etale$ 中每个 $f : U' \to U$，给定同构
$c_f : f_{small}^{-1}\mathcal{F}_U \to \mathcal{F}_{U'}$.
\end{enumerate}
这些数据须满足如下条件：对 $X_\etale$ 中任意
$f : U' \to U$ 与 $g : U'' \to U'$，复合
$c_g \circ g_{small}^{-1} c_f$ 等于 $c_{f \circ g}$。
\end{lemma}

\begin{proof}
可按引理 \ref{spaces-properties-lemma-etale-morphism-topoi} 解释 $g_{small}^{-1}$。
于是本引理由关于站点的一般事实推出，见“站点”引理
\ref{sites-lemma-glue-sheaves-absolute}。
\end{proof}

\noindent
设 $S$ 为概形，$X$ 为 $S$ 上的代数空间。取任意满的 \'etale 态射
$\varphi : U \to X$ 所给出的 $X$ 的表示 $X = U/R$，见“空间”定义
\ref{spaces-definition-presentation}。特别地，我们得到群胚
$(U, R, s, t, c, e, i)$，使得 $j = (t, s) : R \to U \times_S U$，
见“群胚”引理 \ref{groupoids-lemma-equivalence-groupoid}。

\begin{lemma}
\label{spaces-properties-lemma-descent-sheaf}
沿用上述 $S$、$\varphi : U \to X$ 及 $(U, R, s, t, c, e, i)$。
对 $X_\etale$ 上任意层 $\mathcal{F}$，层\footnote{在本引理及其证明中，
我们把 $\varphi_{small}^{-1}$ 简写为 $\varphi^{-1}$，其他拉回也同样简写。}
$\mathcal{G} = \varphi^{-1}\mathcal{F}$ 带有典范同构
$$
\alpha :
t^{-1}\mathcal{G}
\longrightarrow
s^{-1}\mathcal{G}
$$
使得图
$$
\xymatrix{
& \text{pr}_1^{-1}t^{-1}\mathcal{G} \ar[r]_-{\text{pr}_1^{-1}\alpha} &
\text{pr}_1^{-1}s^{-1}\mathcal{G} \ar@{=}[rd] & \\
\text{pr}_0^{-1}s^{-1}\mathcal{G} \ar@{=}[ru] & & &
c^{-1}s^{-1}\mathcal{G} \\
&
\text{pr}_0^{-1}t^{-1}\mathcal{G} \ar[lu]^{\text{pr}_0^{-1}\alpha} \ar@{=}[r] &
c^{-1}t^{-1}\mathcal{G} \ar[ru]_{c^{-1}\alpha}
}
$$
交换。函子 $\mathcal{F} \mapsto (\mathcal{G}, \alpha)$ 给出
$X_\etale$ 上的层与上述二元组 $(\mathcal{G}, \alpha)$ 所成范畴之间的等价。
\end{lemma}

\begin{proof}[引理 \ref{spaces-properties-lemma-descent-sheaf} 的第一种证明]
令 $\mathcal{C} = X_{spaces, \etale}$。由引理
\ref{spaces-properties-lemma-etale-morphism-topoi} 及其证明，
$U_{spaces, \etale} = \mathcal{C}/U$，且拉回函子 $\varphi^{-1}$
就是限制函子。此外，$\{U \to X\}$ 是站点 $\mathcal{C}$ 的覆盖，
且 $R = U \times_X U$。同构 $\alpha$ 就是典范等同
$$
\left(\mathcal{F}|_{\mathcal{C}/U}\right)|_{\mathcal{C}/U \times_X U}
=
\left(\mathcal{F}|_{\mathcal{C}/U}\right)|_{\mathcal{C}/U \times_X U}
$$
而该图的交换性就是胶合数据的上循环条件。因此本引理是层的胶合这一结果的
特殊情形，见“站点”第 \ref{sites-section-glueing-sheaves} 节。
\end{proof}

\begin{proof}[引理 \ref{spaces-properties-lemma-descent-sheaf} 的第二种证明]
$\alpha$ 的存在性来自 $\varphi \circ t = \varphi \circ s$，以及拉回关于态射
具有函子性这一事实，见引理 \ref{spaces-properties-lemma-functoriality-etale-site}。
完全同样地，即由拉回的函子性可见，同构 $\alpha$ 嵌入上述交换图。
构造 $\mathcal{F} \mapsto (\varphi^{-1}\mathcal{F}, \alpha)$
关于层 $\mathcal{F}$ 显然具有函子性。因此我们得到了所需函子。

\medskip\noindent
反过来，设 $(\mathcal{G}, \alpha)$ 为这样的二元组，并设 $V \to X$
为 $X_\etale$ 的对象。这时态射 $V' = U \times_X V \to V$
是概形间满的 \'etale 态射，故 $\{V' \to V\}$ 是 $V$ 的 \'etale 覆盖。
令 $\mathcal{G}' = (V' \to V)^{-1}\mathcal{G}$。由于
$R = U \times_X U$，其中 $t = \text{pr}_0$ 且 $s = \text{pr}_0$，
可见 $V' \times_V V' = R \times_X V$，而投影态射
$s', t' : V' \times_V V' \to V'$ 分别是 $t$ 与 $s$ 的拉回。
因此 $\alpha$ 拉回为同构
$\alpha' : (t')^{-1}\mathcal{G}' \to (s')^{-1}\mathcal{G}'$。
据此直接定义
$$
\xymatrix{
\mathcal{F}(V) \ar@{=}[r] &
\text{等化子}(\mathcal{G}(V') \ar@<1ex>[r] \ar@<-1ex>[r] &
\mathcal{G}(V' \times_V V').
}
$$
略去这确实定义一个层的验证。若存在态射 $V \to U$，则
$\mathcal{G}(V) = \mathcal{F}(V)$；这是因为此时上述等化子为
$H^0(\{V' \to V\}, \mathcal{G}) = \mathcal{G}(V)$。
\end{proof}







\section{小 \'etale 站点的点}
\label{spaces-properties-section-points-small-etale-site}

\noindent
本节是“\'Etale 上同调”第
\ref{etale-cohomology-section-stalks} 节的对应版本。

\begin{definition}
\label{spaces-properties-definition-geometric-point}
设 $S$ 为概形，$X$ 为 $S$ 上的代数空间。
\begin{enumerate}
\item $X$ 的一个{\it 几何点}是态射
$\overline{x} : \Spec(k) \to X$，其中 $k$ 为代数闭域。
我们常滥用记号写成 $\overline{x} = \Spec(k)$。
\item 每个几何点 $\overline{x}$ 都有相应的“像”点 $x \in |X|$。
我们称 $\overline{x}$ 为一个{\it 位于 $x$ 上方的几何点}。
\end{enumerate}
\end{definition}

\noindent
事实表明，可以完全仿照概形小 \'etale 站点的情形，在几何点处取
$X_\etale$ 上层的茎。为此定义如下 \'etale 邻域概念。

\begin{definition}
\label{spaces-properties-definition-etale-neighbourhood}
设 $S$ 为概形，$X$ 为 $S$ 上的代数空间，并设 $\overline{x}$
为 $X$ 的一个几何点。
\begin{enumerate}
\item $X$ 中 $\overline{x}$ 的一个{\it \'etale 邻域}是交换图
$$
\xymatrix{
& U \ar[d]^\varphi \\
{\bar x} \ar[r]^{\bar x} \ar[ur]^{\bar u} & X
}
$$
其中 $\varphi$ 是 $S$ 上代数空间的 \'etale 态射。我们用记号
$\varphi : (U, \overline{u}) \to (X, \overline{x})$ 表示这种情形。
\item 一个{\it \'etale 邻域态射}
$(U, \overline{u}) \to (U', \overline{u}')$
是满足 $\overline{u}' = h \circ \overline{u}$ 的 $X$-态射
$h : U \to U'$。
\end{enumerate}
\end{definition}

\noindent
注意，我们允许 $U$ 为代数空间。当取 $X_\etale$ 上层的茎时，必须限制到
属于 $X_\etale$ 的那些 $U$，故此时只考虑 $U$ 为概形的情形。
另一种做法是使用站点 $X_{space, \etale}$ 并考虑所有 \'etale 邻域。
由下述引理的最后一个断言，两种做法并无区别。

\begin{lemma}
\label{spaces-properties-lemma-cofinal-etale}
设 $S$ 为概形，$X$ 为 $S$ 上的代数空间，并设 $\overline{x}$
为 $X$ 的几何点。\'Etale 邻域范畴是余滤的。更精确地：
\begin{enumerate}
\item 设 $(U_i, \overline{u}_i)_{i = 1, 2}$ 为 $X$ 中
$\overline{x}$ 的两个 \'etale 邻域。则存在第三个 \'etale 邻域
$(U, \overline{u})$ 以及态射
$(U, \overline{u}) \to (U_i, \overline{u}_i)$, $i = 1, 2$.
\item 设 $h_1, h_2: (U, \overline{u}) \to (U', \overline{u}')$
为 $\overline{s}$ 的 \'etale 邻域之间的两个态射。则存在 \'etale 邻域
$(U'', \overline{u}'')$ 及态射
$h : (U'', \overline{u}'') \to (U, \overline{u})$
等化 $h_1$ 与 $h_2$，即满足 $h_1 \circ h = h_2 \circ h$。
\end{enumerate}
此外，给定任意 \'etale 邻域
$(U, \overline{u}) \to (X, \overline{x})$
，存在 \'etale 邻域态射
$(U', \overline{u}') \to (U, \overline{u})$
，其中 $U'$ 为概形。
\end{lemma}

\begin{proof}
对 (1)，考虑纤维积 $U = U_1 \times_X U_2$。由于 \'etale 态射在基变换及
复合下保持，它在 $U_1$ 与 $U_2$ 上都是 \'etale 的；见引理
\ref{spaces-properties-lemma-base-change-etale} 与 \ref{spaces-properties-lemma-composition-etale}。
由 $(\overline{u}_1, \overline{u}_2)$ 定义的映射 $\overline{u} \to U$
使其成为同时映到 $U_1$ 与 $U_2$ 的 \'etale 邻域。

\medskip\noindent
对 (2)，把 $U''$ 定义为纤维积
$$
\xymatrix{
U'' \ar[r] \ar[d] & U \ar[d]^{(h_1, h_2)} \\
U' \ar[r]^-\Delta & U' \times_X U'.
}
$$
由于 $\overline{u}$ 与 $\overline{u}'$ 在 $X$ 上都等于 $\overline{x}$，
可见 $\overline{u}'' = (\overline{u}, \overline{u}')$ 是 $U''$ 的几何点。
特别地，$U'' \not = \emptyset$。此外，由于 $U'$ 在 $X$ 上 \'etale，
纤维积 $U'\times_X U'$ 也同样 \'etale（如上面对
$U_1 \times_X U_2$ 的讨论）。故由引理 \ref{spaces-properties-lemma-etale-permanence}，
纵箭头 $(h_1, h_2)$ 是 \'etale 的。于是基变换表明 $U''$ 在 $U'$ 上
\'etale，进而也在 $X$ 上 \'etale（因为 \'etale 态射的复合仍为
\'etale 态射）。所以 $(U'', \overline{u}'')$ 解出了 (2) 所提问题。

\medskip\noindent
为证最后一个断言，任取满的 \'etale 态射 $U' \to U$，其中 $U'$ 为概形。
于是 $U' \times_U \overline{u}$ 是在
$\overline{u} = \Spec(k)$ 上满且 \'etale 的概形，其中 $k$ 代数闭。
由“态射”引理 \ref{morphisms-lemma-etale-over-field}，
$U' \times_U \overline{u} \to \overline{u}$ 有截面，后者给出所需的
$\overline{u}'$。
\end{proof}

\begin{lemma}
\label{spaces-properties-lemma-geometric-lift-to-usual}
设 $S$ 为概形，$X$ 为 $S$ 上的代数空间。设
$\overline{x} : \Spec(k) \to X$ 为 $X$ 的一个几何点，位于
$x \in |X|$ 上方。设 $\varphi : U \to X$ 为代数空间的 \'etale 态射，
并取满足 $\varphi(u) = x$ 的 $u \in |U|$。则存在位于 $u$ 上方的几何点
$\overline{u} : \Spec(k) \to U$，使得
$\overline{x} = \varphi \circ \overline{u}$。
\end{lemma}

\begin{proof}
选取仿射概形 $U'$、点 $u' \in U'$ 以及把 $u'$ 映到 $u$ 的 \'etale 态射
$U' \to U$。若能对 $(U', u') \to (X, x)$ 证明本引理，结论即随之成立。
故可设 $U$ 为概形；特别地，$U \to X$ 可表示。考察笛卡儿图
$$
\xymatrix{
\Spec(k) \times_{\overline{x}, X, \varphi} U
\ar[d]_{\text{pr}_1} \ar[r]_-{\text{pr}_2} & U
\ar[d]^\varphi \\
\Spec(k) \ar[r]^-{\overline{x}} & X
}
$$
投影 $\text{pr}_1$ 是 \'etale 态射的基变换，故为 \'etale；见引理
\ref{spaces-properties-lemma-base-change-etale}。因此概形
$\Spec(k) \times_{\overline{x}, X, \varphi} U$ 是若干 $k$ 的有限可分扩域
之谱的不交并，见“态射”引理 \ref{morphisms-lemma-etale-over-field}。
但 $k$ 代数闭，故所有这些扩张都平凡，于是
$\Spec(k) \times_{\overline{x}, X, \varphi} U$ 是若干
$\Spec(k)$ 的不交并；其中每一份都对应一个满足
$\varphi \circ \overline{u} = \overline{x}$ 的几何点 $\overline{u}$。
由引理 \ref{spaces-properties-lemma-points-cartesian}，映射
$$
|\Spec(k) \times_{\overline{x}, X, \varphi} U|
\longrightarrow
|\Spec(k)| \times_{|X|} |U|
$$
是满的，故可选取位于 $u$ 上方的 $\overline{u}$。
\end{proof}

\begin{lemma}
\label{spaces-properties-lemma-geometric-lift-to-cover}
设 $S$ 为概形，$X$ 为 $S$ 上的代数空间，$\overline{x}$ 为 $X$
的几何点，且 $(U, \overline{u})$ 为 $\overline{x}$ 的 \'etale 邻域。
设 $\{\varphi_i : U_i \to U\}_{i \in I}$ 为
$X_{spaces, \etale}$ 中的一个 \'etale 覆盖。则存在 $i \in I$ 及
$\overline{u}_i : \overline{x} \to U_i$，使得
$\varphi_i : (U_i, \overline{u}_i) \to (U, \overline{u})$
为 \'etale 邻域态射。
\end{lemma}

\begin{proof}
令 $u \in |U|$ 为 $\overline{u}$ 的像。由于
$|U| = \bigcup_{i \in I} \varphi_i(|U_i|)$，存在某个 $i$ 以及映到
$x$ 的点 $u_i \in U_i$。对 $(U_i, u_i) \to (U, u)$ 与
$\overline{u}$ 应用引理 \ref{spaces-properties-lemma-geometric-lift-to-usual}，即得所需几何点。
\end{proof}

\begin{definition}
\label{spaces-properties-definition-stalk}
设 $S$ 为概形，$X$ 为 $S$ 上的代数空间，$\mathcal{F}$ 为
$X_\etale$ 上的预层，并设 $\overline{x}$ 为 $X$ 的几何点。
$\mathcal{F}$ 在 $\overline{x}$ 处的{\it 茎}是
$$
\mathcal{F}_{\bar x}
=
\colim_{(U, \overline{u})} \mathcal{F}(U)
$$
其中 $(U, \overline{u})$ 遍历 $X$ 中 $\overline{x}$ 的所有 \'etale 邻域，
且 $U \in \Ob(X_\etale)$。
\end{definition}

\noindent
由引理 \ref{spaces-properties-lemma-cofinal-etale}，此余极限的指标范畴是滤过的，
它就是 $X_\etale$ 中 \'etale 邻域范畴的反范畴。更精确地，
引理 \ref{spaces-properties-lemma-cofinal-etale} 说明所有 \'etale 邻域所成范畴的反范畴
是滤过的，而由属于 $X_\etale$ 的那些邻域组成的全子范畴是共尾子范畴，
因而也滤过。

\medskip\noindent
这意味着，可把 $\mathcal{F}_{\overline{x}}$ 的一个元素看作三元组
$(U, \overline{u}, \sigma)$，其中 $U \in \Ob(X_\etale)$ 且
$\sigma \in \mathcal{F}(U)$。两个三元组
$(U, \overline{u}, \sigma)$、$(U', \overline{u}', \sigma')$
定义茎中的同一个元素，当且仅当存在第三个 \'etale 邻域
$(U'', \overline{u}'')$, $U'' \in \Ob(X_\etale)$
以及 \'etale 邻域态射
$h : (U'', \overline{u}'') \to (U, \overline{u})$,
$h' : (U'', \overline{u}'') \to (U', \overline{u}')$，使得
$h^*\sigma = (h')^*\sigma'$ 在 $\mathcal{F}(U'')$ 中成立。
见“范畴”第 \ref{categories-section-directed-colimits} 节。

\medskip\noindent
这还说明，若 $\mathcal{F}'$ 是 $X_{spaces, \etale}$ 上对应于
$X_\etale$ 上 $\mathcal{F}$ 的层，则
\begin{equation}
\label{spaces-properties-equation-stalk-spaces-etale}
\mathcal{F}_{\overline{x}} = \colim_{(U, \overline{u})} \mathcal{F}'(U)
\end{equation}
其中的余极限现在遍历 $\overline{x}$ 的所有 \'etale 邻域。
以后常在使用 $X_\etale$ 与使用 $X_{spaces, \etale}$ 的观点之间切换，
而不再特别说明。

\medskip\noindent
特别地，若 $\mathcal{F}$ 是阿贝尔群、环等的预层，则按通常在阿贝尔群、
环等的有向余极限上定义群结构的方法，$\mathcal{F}_{\overline{x}}$
自然也是阿贝尔群、环等。

\begin{lemma}
\label{spaces-properties-lemma-stalk-gives-point}
\begin{slogan}
代数空间的一个几何点给出其 \'etale 拓扑斯的一个点。
\end{slogan}
设 $S$ 为概形，$X$ 为 $S$ 上的代数空间，并设 $\overline{x}$
为 $X$ 的几何点。考虑函子
$$
u : X_\etale \longrightarrow \textit{Sets}, \quad
U \longmapsto |U_{\overline{x}}|
$$
则 $u$ 定义站点 $X_\etale$ 的一个点 $p$
（“站点”定义 \ref{sites-definition-point}），且与之相伴的茎函子
$\mathcal{F} \mapsto \mathcal{F}_p$
（“站点”式 \ref{sites-equation-stalk}）正是上面定义的函子
$\mathcal{F} \mapsto \mathcal{F}_{\overline{x}}$。
\end{lemma}

\begin{proof}
在引理 \ref{spaces-properties-lemma-geometric-lift-to-usual} 的证明中已经看到，概形
$U_{\overline{x}} = \overline{x} \times_X U$ 是若干与
$\overline{x}$ 同构的概形的不交并。因此也可把 $|U_{\overline{x}}|$
看作 $U$ 中位于 $\overline{x}$ 上方的几何点之集，即所有能嵌入定义
\ref{spaces-properties-definition-etale-neighbourhood} 之图的态射
$\overline{u} : \overline{x} \to U$ 所成的集合。由此，$u(X)$ 为单点集，且
$u(U \times_V W) = u(U) \times_{u(V)} u(W)$
对 $X_\etale$ 中任意态射 $U \to V$ 与 $W \to V$ 都成立。
再者，给定 $X_\etale$ 中的覆盖 $\{U_i \to U\}_{i \in I}$，
由引理 \ref{spaces-properties-lemma-geometric-lift-to-cover} 可知
$\coprod u(U_i) \to u(U)$ 是满的。因此可应用“站点”命题
\ref{sites-proposition-point-limits}，故 $p$ 是站点 $X_\etale$ 的点。
最后，我们的函子 $\mathcal{F} \mapsto \mathcal{F}_{\overline{s}}$
与“站点”式 \ref{sites-equation-stalk} 中相伴于 $p$ 的函子
$\mathcal{F} \mapsto \mathcal{F}_p$ 由完全相同的余极限给出，
这证明了最后一个断言。
\end{proof}

\begin{lemma}
\label{spaces-properties-lemma-stalk-exact}
设 $S$ 为概形，$X$ 为 $S$ 上的代数空间，并设 $\overline{x}$
为 $X$ 的几何点。
\begin{enumerate}
\item 茎函子
$\textit{PAb}(X_\etale) \to \textit{Ab}$,
$\mathcal{F}  \mapsto  \mathcal{F}_{\overline{x}}$
是正合的。
\item 对 $X_\etale$ 上任意集合预层 $\mathcal{F}$，有
$(\mathcal{F}^\#)_{\overline{x}} = \mathcal{F}_{\overline{x}}$。
\item 函子
$\textit{Ab}(X_\etale) \to \textit{Ab}$,
$\mathcal{F} \mapsto \mathcal{F}_{\overline{x}}$ 是正合的。
\item 类似地，由茎函子
$\mathcal{F} \mapsto \mathcal{F}_{\overline{x}}$ 给出的函子
$\textit{PSh}(X_\etale) \to \textit{Sets}$ and
$\Sh(X_\etale) \to \textit{Sets}$ 也是正合的（见“范畴”定义
\ref{categories-definition-exact}），并与任意余极限交换。
\end{enumerate}
\end{lemma}

\begin{proof}
该结果由“站点上的模”第 \ref{sites-modules-section-stalks} 节的一般理论推出。
这是因为 $\mathcal{F} \mapsto \mathcal{F}_{\overline{x}}$ 来自
$X$ 的小 \'etale 站点的一个点，见引理 \ref{spaces-properties-lemma-stalk-gives-point}。
在概形小 \'etale 站点的情形，这些陈述中一部分的直接证明可见
“\'Etale 上同调”引理 \ref{etale-cohomology-lemma-stalk-exact} 的证明。
\end{proof}

\noindent
下面将看到，茎函子
$\mathcal{F} \mapsto \mathcal{F}_{\overline{x}}$
其实就是沿态射 $\overline{x}$ 的拉回。在此意义下，下述引理是上面引理的推广。

\begin{lemma}
\label{spaces-properties-lemma-stalk-pullback}
设 $S$ 为概形，$f : X \to Y$ 为 $S$ 上代数空间的态射。
\begin{enumerate}
\item 函子
$f_{small}^{-1} :
\textit{Ab}(Y_\etale)
\to
\textit{Ab}(X_\etale)$
是正合的。
\item 函子
$f_{small}^{-1} :
\Sh(Y_\etale)
\to
\Sh(X_\etale)$
是正合的，即它与有限极限及有限余极限交换；见“范畴”定义
\ref{categories-definition-exact}。
\item 对代数空间的任意 \'etale 态射 $V \to Y$，有
$f_{small}^{-1}h_V = h_{X \times_Y V}$。
\item 设 $\overline{x} \to X$ 为几何点，$\mathcal{G}$ 为
$Y_\etale$ 上的层。则有典范等同
$$
(f_{small}^{-1}\mathcal{G})_{\overline{x}} = \mathcal{G}_{\overline{y}}.
$$
其中 $\overline{y} = f \circ \overline{x}$。
\end{enumerate}
\end{lemma}

\begin{proof}
回忆，在引理 \ref{spaces-properties-lemma-functoriality-etale-site} 中，$f_{small}$
是通过 $f_{spaces, small}$ 定义的。(1)、(2)、(3) 是下述事实的一般推论：
$f_{spaces, \etale} :
X_{spaces, \etale}
\to
Y_{spaces, \etale}$
是站点态射。关于 (2) 见“站点”定义
\ref{sites-definition-morphism-sites}，关于 (1) 见“站点上的模”引理
\ref{sites-modules-lemma-flat-pullback-exact}，关于 (3) 见“站点”引理
\ref{sites-lemma-pullback-representable-sheaf}。

\medskip\noindent
证明 (4)。该陈述是“站点”引理 \ref{sites-lemma-point-morphism-sites}
通过引理 \ref{spaces-properties-lemma-stalk-gives-point} 得到的特殊情形。这里也给出直接证明。
注意，由引理 \ref{spaces-properties-lemma-stalk-exact}，取茎与层化交换。设
$\mathcal{G}'$ 为 $Y_{spaces, \etale}$ 上限制到 $Y_\etale$ 后等于
$\mathcal{G}$ 的层。回忆，$f_{spaces, \etale}^{-1}\mathcal{G}'$
是与下列预层相伴的层：
$$
U \longrightarrow \colim_{U \to X \times_Y V} \mathcal{G}'(V),
$$
见“站点”第 \ref{sites-section-continuous-functors} 节与第
\ref{sites-section-functoriality-PSh} 节。因此
\begin{align*}
(f_{spaces, \etale}^{-1}\mathcal{G}')_{\overline{x}}
& =
\colim_{(U, \overline{u})} f_{spaces, \etale}^{-1}\mathcal{G}'(U)
\\
& = \colim_{(U, \overline{u})}
\colim_{a : U \to X \times_Y V} \mathcal{G}'(V) \\
& = \colim_{(V, \overline{v})} \mathcal{G}'(V) \\
& = \mathcal{G}'_{\overline{y}}
\end{align*}
在第三个等号中，二元组 $(U, \overline{u})$ 与映射
$a : U \to X \times_Y V$ 对应于二元组 $(V, a \circ \overline{u})$。
由于 $\mathcal{G}'$ 的茎（相应地，
$f_{spaces, \etale}^{-1}\mathcal{G}'$ 的茎）等于 $\mathcal{G}$ 的茎
（相应地，$f_{small}^{-1}\mathcal{G}$ 的茎），见式
(\ref{spaces-properties-equation-stalk-spaces-etale})，结论成立。
\end{proof}

\begin{remark}
\label{spaces-properties-remark-stalk-pullback}
本注是“\'Etale 上同调”注
\ref{etale-cohomology-remark-stalk-pullback} 的对应版本。
设 $S$ 为概形，$X$ 为 $S$ 上的代数空间，并设
$\overline{x} : \Spec(k) \to X$ 为 $X$ 的几何点。
由“\'Etale 上同调”定理
\ref{etale-cohomology-theorem-equivalence-sheaves-point}，
$\Spec(k)_\etale$ 上层的范畴等价于集合范畴（把层送到其全局截面）。
因此，对态射 $\overline{x}$ 应用引理 \ref{spaces-properties-lemma-stalk-pullback} 的 (4)，
可知函子
$$
\Sh(X_\etale) \longrightarrow \textit{Sets}, \quad
\mathcal{F} \longmapsto \mathcal{F}_{\overline{x}}
$$
同构于函子
$$
\Sh(X_\etale)
\longrightarrow
\Sh(\Spec(k)_\etale) = \textit{Sets},
\quad
\mathcal{F} \longmapsto \overline{x}^*\mathcal{F}
$$
因此可把茎函子看作沿几何态射的拉回函子，而不只是引理
\ref{spaces-properties-lemma-stalk-gives-point} 的结果中那种抽象拓扑斯态射。
\end{remark}

\begin{remark}
\label{spaces-properties-remark-map-stalks}
设 $S$ 为概形，$X$ 为 $S$ 上的代数空间，并取 $x \in |X|$。
我们断言，对位于 $x$ 上方的任意两个几何点 $\overline{x}$ 与
$\overline{x}'$，相应茎函子彼此同构。由 $|X|$ 的定义，可找到第三个几何点
$\overline{x}''$，使得存在交换图
$$
\xymatrix{
\overline{x}'' \ar[r] \ar[d] \ar[rd]^{\overline{x}''} &
\overline{x}' \ar[d]^{\overline{x}'} \\
\overline{x} \ar[r]^{\overline{x}} & X.
}
$$
由于茎函子 $\mathcal{F} \mapsto \mathcal{F}_{\overline{x}}$
由沿态射 $\overline{x}$ 的拉回给出（其他两个也同样），
结论由拉回的函子性得到。
\end{remark}




\noindent
下述定理说明，代数空间的小 \'etale 站点有足够多的点。

\begin{theorem}
\label{spaces-properties-theorem-exactness-stalks}
设 $S$ 为概形，$X$ 为 $S$ 上的代数空间。集合层的映射
$a : \mathcal{F} \to \mathcal{G}$ 是单射（相应地，满射），当且仅当茎映射
$a_{\overline{x}} : \mathcal{F}_{\overline{x}} \to \mathcal{G}_{\overline{x}}$
对 $X$ 的所有几何点都是单射（相应地，满射）。$X_\etale$ 上的一个阿贝尔层列
是正合的，当且仅当它在 $S$ 的每个几何点处的茎上正合。
\end{theorem}

\begin{proof}
若 $X$ 为概形，本定理已知成立；见“\'Etale 上同调”定理
\ref{etale-cohomology-theorem-exactness-stalks}。选取满的 \'etale 态射
$f : U \to X$，其中 $U$ 为概形。由于 $\{U \to X\}$ 是覆盖
（在 $X_{spaces, \etale}$ 中），可通过限制到 $U$ 来检验层映射是否为单射或满射。
现在，若 $\overline{u} : \Spec(k) \to U$ 是 $U$ 的几何点，则
$(\mathcal{F}|_U)_{\overline{u}} = \mathcal{F}_{\overline{x}}$
，其中 $\overline{x} = f \circ \overline{u}$。（这从定义
$\overline{u}$ 与 $\overline{x}$ 处茎的余极限显然可见，也可由引理
\ref{spaces-properties-lemma-stalk-pullback} 推出。）故 $U$ 的结果蕴含 $X$ 的结果，证毕。
\end{proof}

\noindent
初读时应跳过下述引理。

\begin{lemma}
\label{spaces-properties-lemma-points-small-etale-site}
设 $S$ 为概形，$X$ 为 $S$ 上的代数空间，并设
$p : \Sh(pt) \to \Sh(X_\etale)$ 为 $X$ 的小 \'etale 拓扑斯的一个点。
则存在 $X$ 的几何点 $\overline{x}$，使得茎函子
$\mathcal{F} \mapsto \mathcal{F}_p$ 同构于茎函子
$\mathcal{F} \mapsto \mathcal{F}_{\overline{x}}$。
\end{lemma}

\begin{proof}
由“站点”引理 \ref{sites-lemma-point-site-topos}，站点的点与相伴拓扑斯的点
一一对应。因此可设 $p$ 由定义站点 $X_\etale$ 的一个点的函子
$u : X_\etale \to \textit{Sets}$ 给出。取对象
$U \in \Ob(X_\etale)$，使其结构态射 $j : U \to X$ 为满射。
注意，$h_U$ 是满射到终层的层。由于取茎正合，
$(h_U)_p = u(U)$ 非空（使用“站点”引理
\ref{sites-lemma-points-recover}）。取 $x \in u(U)$。由“站点”引理
\ref{sites-lemma-point-localize}，得到点
$q : \Sh(pt) \to \Sh(U_\etale)$，使得 $p = j_{small} \circ q$，
因而函子性地有 $\mathcal{F}_p = (\mathcal{F}|_U)_q$。
由“\'Etale 上同调”引理
\ref{etale-cohomology-lemma-points-small-etale-site}，存在 $U$ 的几何点
$\overline{u}$，并且对 $\mathcal{G} \in \Sh(U_\etale)$ 有函子同构
$\mathcal{G}_q = \mathcal{G}_{\overline{u}}$。令
$\overline{x} = j \circ \overline{u}$。于是引理
\ref{spaces-properties-lemma-stalk-pullback} 表明，对 $X_\etale$ 上的 $\mathcal{F}$，函子性地有
$\mathcal{F}_{\overline{x}} \cong (\mathcal{F}|_U)_{\overline{u}}$，证毕。
\end{proof}





\section{阿贝尔层的支撑}
\label{spaces-properties-section-support}

\noindent
先讨论局部截面的支撑。

\begin{lemma}
\label{spaces-properties-lemma-support-subsheaf-final}
设 $S$ 为概形，$X$ 为 $S$ 上的代数空间。设 $\mathcal{F}$ 为
$X$ 的 \'etale 拓扑斯之终对象的子层（见“站点”例
\ref{sites-example-singleton-sheaf}）。则存在唯一的开子空间
$W \subset X$，使得 $\mathcal{F} = h_W$。
\end{lemma}

\begin{proof}
该条件意味着，对 $\Ob(X_{spaces, \etale})$ 中每个
$\varphi : U \to X$，$\mathcal{F}(U)$ 或为单点集，或为空集。
特别地，局部截面总能胶合。若 $\mathcal{F}(U) \not = \emptyset$，则
$\mathcal{F}(\varphi(U)) \not = \emptyset$，因为
$\varphi(U) \subset X$ 是开子空间（引理 \ref{spaces-properties-lemma-etale-open}），且
$\{\varphi : U \to \varphi(U)\}$ 是 $X_{spaces, \etale}$ 中的覆盖。
取
$W = \bigcup_{\varphi : U \to S, \mathcal{F}(U) \not = \emptyset} \varphi(U)$
即得结论。
\end{proof}

\begin{lemma}
\label{spaces-properties-lemma-zero-over-image}
设 $S$ 为概形，$X$ 为 $S$ 上的代数空间，$\mathcal{F}$ 为
$X_{spaces, \etale}$ 上的阿贝尔层，并设
$\sigma \in \mathcal{F}(U)$ 为局部截面。则存在开子空间
$W \subset U$，使得
\begin{enumerate}
\item $W \subset U$ 是使得 $\sigma|_W = 0$ 的 $U$ 的最大开子空间；
\item 对 $X_{spaces, \etale}$ 中每个 $\varphi : V \to U$，有
$$
\sigma|_V = 0 \Leftrightarrow \varphi(V) \subset W,
$$
\item 对 $U$ 的每个几何点 $\overline{u}$，有
$$
(U, \overline{u}, \sigma) = 0\text{ 于 }\mathcal{F}_{\overline{x}}
\Leftrightarrow
\overline{u} \in W
$$
其中 $\overline{x} = (U \to X) \circ \overline{u}$。
\end{enumerate}
\end{lemma}

\begin{proof}
由于 $\mathcal{F}$ 是 \'etale 拓扑中的层，$\mathcal{F}$ 到 $U_{Zar}$ 的限制是
$U$ 上 Zariski 拓扑中的层。因此存在满足 (1) 的 Zariski 开集 $W$，
见“模”引理 \ref{modules-lemma-support-section-closed}。设
$\varphi : V \to U$ 为 $X_{spaces, \etale}$ 中的箭头。注意，
$\varphi(V) \subset U$ 是开子空间（引理 \ref{spaces-properties-lemma-etale-open}），
且 $\{V \to \varphi(V)\}$ 是 \'etale 覆盖。因此若
$\sigma|_V = 0$，则由 $\mathcal{F}$ 的层条件，
$\sigma|_{\varphi(V)} = 0$。这证明了 (2)。为证明 (3)，须证若
$(U, \overline{u}, \sigma)$ 定义 $\mathcal{F}_{\overline{x}}$ 的零元，
则 $\overline{u} \in W$。该假设意味着存在 \'etale 邻域态射
$(V, \overline{v}) \to (U, \overline{u})$，使得
$\sigma|_V = 0$。故由 (2)，$V \to U$ 的像包含于 $W$，
从而 $\overline{u} \in W$。
\end{proof}

\noindent
设 $S$ 为概形，$X$ 为 $S$ 上的代数空间，取 $x \in |X|$，并设
$\mathcal{F}$ 为 $X_\etale$ 上的层。由注 \ref{spaces-properties-remark-map-stalks}，
$\mathcal{F}$ 在位于 $x$ 上方的几何点处之茎的同构类是良定义的。

\begin{definition}
\label{spaces-properties-definition-support}
设 $S$ 为概形，$X$ 为 $S$ 上的代数空间，并设 $\mathcal{F}$
为 $X_\etale$ 上的阿贝尔层。
\begin{enumerate}
\item $\mathcal{F}$ 的{\it 支撑}是这样的点 $x \in |X|$ 所成的集合：
对位于 $x$ 上方的任意（等价地，某个）几何点 $\overline{x}$，有
$\mathcal{F}_{\overline{x}} \not = 0$。
\item 设 $\sigma \in \mathcal{F}(U)$ 为截面。$\sigma$ 的{\it 支撑}
是闭子集 $U \setminus W$，其中 $W \subset U$ 是 $U$ 中使
$\sigma$ 限制为零的最大开子集（见引理 \ref{spaces-properties-lemma-zero-over-image}）。
\end{enumerate}
\end{definition}

\begin{lemma}
\label{spaces-properties-lemma-support-section-closed}
设 $S$ 为概形，$X$ 为 $S$ 上的代数空间，$\mathcal{F}$ 为
$X_\etale$ 上的阿贝尔层，并取 $U \in \Ob(X_\etale)$ 与
$\sigma \in \mathcal{F}(U)$。
\begin{enumerate}
\item $\sigma$ 的支撑在 $|X|$ 中闭。
\item $\sigma + \sigma'$ 的支撑包含于
$\sigma, \sigma' \in \mathcal{F}(X)$ 的支撑之并。
\item 若 $\varphi : \mathcal{F} \to \mathcal{G}$ 是 $X_\etale$
上阿贝尔层的映射，则 $\varphi(\sigma)$ 的支撑包含于
$\sigma \in \mathcal{F}(U)$ 的支撑。
\item $\mathcal{F}$ 的支撑是 $\mathcal{F}$ 所有局部截面的支撑之像的并。
\item 若 $\mathcal{F} \to \mathcal{G}$ 满，则 $\mathcal{G}$ 的支撑是
$\mathcal{F}$ 的支撑的子集。
\item 若 $\mathcal{F} \to \mathcal{G}$ 单，则 $\mathcal{F}$ 的支撑是
$\mathcal{G}$ 的支撑的子集。
\end{enumerate}
\end{lemma}

\begin{proof}
(1) 由定义成立。(2) 与 (3) 成立，因为它们对 $\mathcal{F}$ 与
$\mathcal{G}$ 到 $U_{Zar}$ 的限制成立；见“模”引理
\ref{modules-lemma-support-section-closed}。(4) 是引理
\ref{spaces-properties-lemma-zero-over-image} 的 (3) 的直接推论。(5) 与 (6) 由其余各项推出。
\end{proof}

\begin{lemma}
\label{spaces-properties-lemma-support-sheaf-rings-closed}
代数空间小 \'etale 站点上的环层之支撑是闭的。
\end{lemma}

\begin{proof}
这是因为，按我们的约定，一个环等于 $0$ 当且仅当 $1 = 0$；
因此环层的支撑就是单位截面的支撑。
\end{proof}






\section{代数空间的结构层}
\label{spaces-properties-section-structure-sheaf}

\noindent
代数空间的结构层就是下述引理中的环层。

\begin{lemma}
\label{spaces-properties-lemma-sheaf-condition-holds}
设 $S$ 为概形，$X$ 为 $S$ 上的代数空间。规则
$U \mapsto \Gamma(U, \mathcal{O}_U)$ 定义 $X_\etale$ 上的环层。
\end{lemma}

\begin{proof}
由覆盖的定义及“下降”引理 \ref{descent-lemma-sheaf-condition-holds} 立即可得。
\end{proof}

\begin{definition}
\label{spaces-properties-definition-structure-sheaf}
设 $S$ 为概形，$X$ 为 $S$ 上的代数空间。$X$ 的{\it 结构层}是引理
\ref{spaces-properties-lemma-sheaf-condition-holds} 所述小 \'etale 站点
$X_\etale$ 上的环层 $\mathcal{O}_X$。
\end{definition}

\noindent
根据引理 \ref{spaces-properties-lemma-characterize-sheaf-small-etale}，层 $\mathcal{O}_X$
对应于一族 \'etale 层 $(\mathcal{O}_X)_U$，其中 $U$ 遍历
$X_\etale$ 的对象。由该引理的证明及我们的定义，显然直接有
$(\mathcal{O}_X)_U = \mathcal{O}_U$；这里 $\mathcal{O}_U$ 是“下降”定义
\ref{descent-definition-structure-sheaf} 中引入的 $U_\etale$ 的结构层。
特别地，若 $X$ 为概形，我们就恢复了 $X$ 的小 \'etale 站点上的层
$\mathcal{O}_X$。

\medskip\noindent
通过引理 \ref{spaces-properties-lemma-compare-etale-sites} 的等价
$\Sh(X_\etale) = \Sh(X_{spaces, \etale})$
，也可把 $\mathcal{O}_X$ 看作 $X_{spaces, \etale}$ 上的环层。
注 \ref{spaces-properties-remark-explain-equivalence} 说明了当 $Y \to X$ 是
$X_{spaces, \etale}$ 的对象时，如何计算 $\mathcal{O}_X(Y)$，
特别是 $\mathcal{O}_X(X)$。

\begin{lemma}
\label{spaces-properties-lemma-morphism-ringed-topoi}
设 $S$ 为概形，$f : X \to Y$ 为 $S$ 上代数空间的态射。
则存在典范映射
$f^\sharp : f_{small}^{-1}\mathcal{O}_Y \to \mathcal{O}_X$，使得
$$
(f_{small}, f^\sharp) :
(\Sh(X_\etale), \mathcal{O}_X)
\longrightarrow
(\Sh(Y_\etale), \mathcal{O}_Y)
$$
为环化拓扑斯态射。此外，
\begin{enumerate}
\item 构造 $f \mapsto (f_{small}, f^\sharp)$ 与复合相容。
\item 若 $f$ 是概形态射，则 $f^\sharp$ 是“下降”注
\ref{descent-remark-change-topologies-ringed} 所述映射。
\end{enumerate}
\end{lemma}

\begin{proof}
由引理 \ref{spaces-properties-lemma-f-map}，只须给出从 $\mathcal{O}_Y$ 到
$\mathcal{O}_X$ 的一个 $f$-映射。换言之，对每个交换图
$$
\xymatrix{
U \ar[d]_g \ar[r] & X \ar[d]^f \\
V \ar[r] & Y
}
$$
其中 $U \in X_\etale$、$V \in Y_\etale$，须给出环同态
$
(f^\sharp)_{(U, V, g)} :
\Gamma(V, \mathcal{O}_V)
\to
\Gamma(U, \mathcal{O}_U).
$
当然，只需取 $(f^\sharp)_{(U, V, g)} = g^\sharp$。这与限制映射相容是清楚的，
故确实给出一个 $f$-映射。略去它与复合相容，以及它与“下降”注
\ref{descent-remark-change-topologies-ringed} 中构造相同的验证。
\end{proof}

\begin{lemma}
\label{spaces-properties-lemma-reduced-space}
设 $S$ 为概形，$X$ 为 $S$ 上的代数空间。下列条件等价：
\begin{enumerate}
\item $X$ 既约；
\item 对每个 $x \in |X|$，$X$ 在 $x$ 处的局部环既约
（注 \ref{spaces-properties-remark-list-properties-local-ring-local-etale-topology}）。
\end{enumerate}
在此情形，$\Gamma(X, \mathcal{O}_X)$ 是既约环；而且若
$f \in \Gamma(X, \mathcal{O}_X)$ 满足 $X = V(f)$，则 $f = 0$。
\end{lemma}

\begin{proof}
对 $X$ 上 \'etale 的仿射概形应用“性质”引理
\ref{properties-lemma-characterize-reduced}，即得 (1) 与 (2) 等价。
最后两个陈述由所引引理以及如下事实推出：对某个在 $X$ 上 \'etale 的既约概形
$U$，$\Gamma(X, \mathcal{O}_X)$ 是 $\Gamma(U, \mathcal{O}_U)$ 的子环。
\end{proof}







\section{结构层的茎}
\label{spaces-properties-section-stalks-structure-sheaf}

\noindent
本节是“\'Etale 上同调”第
\ref{etale-cohomology-section-stalks-structure-sheaf} 节的对应版本。

\begin{lemma}
\label{spaces-properties-lemma-describe-etale-local-ring}
设 $S$ 为概形，$X$ 为 $S$ 上的代数空间，$\overline{x}$ 为 $X$
的几何点，并设 $(U, \overline{u})$ 为 $\overline{x}$ 的 \'etale 邻域，
其中 $U$ 为概形。则
$$
\mathcal{O}_{X, \overline{x}} =
\mathcal{O}_{U, \overline{u}} =
\mathcal{O}_{U, u}^{sh}
$$
其中左端是 $X$ 的结构层之茎，而右端是 $U$ 在点 $u$ 处局部环的严格
Hensel 化；$\overline{u}$ 即以该点为中心。
\end{lemma}

\begin{proof}
已知 $U_\etale$ 上的结构层 $\mathcal{O}_U$ 是 $X$ 的结构层的限制。
故第一个等号由引理 \ref{spaces-properties-lemma-stalk-pullback} 的 (4) 得到。
第二个等号见“\'Etale 上同调”引理
\ref{etale-cohomology-lemma-describe-etale-local-ring}。
\end{proof}

\begin{definition}
\label{spaces-properties-definition-etale-local-rings}
设 $S$ 为概形，$X$ 为 $S$ 上的代数空间，并设 $\overline{x}$
为 $X$ 的几何点，位于点 $x \in |X|$ 上方。
\begin{enumerate}
\item $X$ 在 $\overline{x}$ 处的{\it \'etale 局部环}是
$X_\etale$ 上结构层 $\mathcal{O}_X$ 在 $\overline{x}$ 处的茎。
记作 $\mathcal{O}_{X, \overline{x}}$。
\item $X$ 在 $\overline{x}$ 处的{\it 严格 Hensel 化}是概形
$\Spec(\mathcal{O}_{X, \overline{x}})$。
\end{enumerate}
\end{definition}

\noindent
$X$ 在 $\overline{x}$ 处的严格 Hensel 化作为 $X$ 上概形的同构类型，
仅依赖于点 $x \in |X|$，而与位于 $x$ 上方的几何点之选取无关；
见注 \ref{spaces-properties-remark-map-stalks}。

\begin{lemma}
\label{spaces-properties-lemma-etale-site-locally-ringed}
设 $S$ 为概形，$X$ 为 $S$ 上的代数空间。赋予结构层
$\mathcal{O}_X$ 的小 \'etale 站点 $X_\etale$ 是局部环化站点；
见“站点上的模”定义 \ref{sites-modules-definition-locally-ringed}。
\end{lemma}

\begin{proof}
这是因为各茎 $\mathcal{O}_{X, \overline{x}}$ 都是局部环，并且
$S_\etale$ 有足够多的点；见引理 \ref{spaces-properties-lemma-describe-etale-local-ring}
与定理 \ref{spaces-properties-theorem-exactness-stalks}。这些事实蕴含小 \'etale 站点局部环化，
见“站点上的模”引理 \ref{sites-modules-lemma-locally-ringed-stalk} 与
\ref{sites-modules-lemma-ringed-stalk-not-zero}。
\end{proof}

\begin{lemma}
\label{spaces-properties-lemma-dimension-local-ring}
设 $S$ 为概形，$X$ 为 $S$ 上的代数空间。取点 $x \in |X|$ 及
$d \in \{0, 1, 2, \ldots, \infty\}$。下列条件等价：
\begin{enumerate}
\item $X$ 在 $x$ 处局部环的维数（定义
\ref{spaces-properties-definition-dimension-local-ring}）为 $d$；
\item 对位于 $x$ 上方的某个几何点 $\overline{x}$，有
$\dim(\mathcal{O}_{X, \overline{x}}) = d$；
\item 对位于 $x$ 上方的任意几何点 $\overline{x}$，有
$\dim(\mathcal{O}_{X, \overline{x}}) = d$。
\end{enumerate}
\end{lemma}

\begin{proof}
(2) 与 (3) 的等价性来自如下事实：$\mathcal{O}_{X, \overline{x}}$
的同构类型只依赖于 $x \in |X|$；见注 \ref{spaces-properties-remark-map-stalks}。
利用引理 \ref{spaces-properties-lemma-describe-etale-local-ring}，(1) 与 (2)$+$(3) 的等价性
归结为下述陈述：对任意局部环 $R$，有
$\dim(R) = \dim(R^{sh})$。这正是“更多代数”引理
\ref{more-algebra-lemma-henselization-dimension}。
\end{proof}

\begin{lemma}
\label{spaces-properties-lemma-dimension-decent-invariant-under-etale}
设 $S$ 为概形，$f : X \to Y$ 为 $S$ 上代数空间的 \'etale 态射，
并取 $x \in X$。则 (1) $\dim_x(X) = \dim_{f(x)}(Y)$；
(2) $X$ 在 $x$ 处局部环的维数等于 $Y$ 在 $f(x)$ 处局部环的维数。
若 $f$ 满，则 (3) $\dim(X) = \dim(Y)$。
\end{lemma}

\begin{proof}
选取概形 $U$、点 $u \in U$ 以及把 $u$ 映到 $x$ 的 \'etale 态射
$U \to X$。则复合 $U \to Y$ 也为 \'etale，并把 $u$ 映到 $f(x)$。
相关数值是依照概形 $U$ 在 $u$ 处的性态定义的，故 (1) 与 (2) 成立。
关于 (1) 见定义 \ref{spaces-properties-definition-dimension-at-point}。(3) 是 (1) 的直接推论，
见定义 \ref{spaces-properties-definition-dimension}。
\end{proof}

\begin{lemma}
\label{spaces-properties-lemma-reduced-local-ring}
设 $S$ 为概形，$X$ 为 $S$ 上的代数空间，并取点 $x \in |X|$。
下列条件等价：
\begin{enumerate}
\item $X$ 在 $x$ 处的局部环既约
（注 \ref{spaces-properties-remark-list-properties-local-ring-local-etale-topology}）；
\item 对位于 $x$ 上方的某个几何点 $\overline{x}$，
$\mathcal{O}_{X, \overline{x}}$ 既约；
\item 对位于 $x$ 上方的任意几何点 $\overline{x}$，
$\mathcal{O}_{X, \overline{x}}$ 既约。
\end{enumerate}
\end{lemma}

\begin{proof}
(2) 与 (3) 的等价性来自如下事实：$\mathcal{O}_{X, \overline{x}}$
的同构类型只依赖于 $x \in |X|$；见注 \ref{spaces-properties-remark-map-stalks}。
利用引理 \ref{spaces-properties-lemma-describe-etale-local-ring}，(1) 与 (2)$+$(3) 的等价性
归结为：一个局部环既约，当且仅当它的严格 Hensel 化既约。
这正是“更多代数”引理 \ref{more-algebra-lemma-henselization-reduced}。
\end{proof}











\section{局部不可约性}
\label{spaces-properties-section-irreducible-local-ring}

\noindent
代数空间上的一个点具有良定义的 \'etale 局部环；在概形的情形，
它对应于局部环的严格 Hensel 化。一般而言，仅由 \'etale 局部环无法看出
概形或代数空间有多少个不可约分支通过给定点；我们只能计算几何分支的数目。

\begin{lemma}
\label{spaces-properties-lemma-irreducible-local-ring}
设 $S$ 为概形，$X$ 为 $S$ 上的代数空间，并取点 $x \in |X|$。
下列条件等价：
\begin{enumerate}
\item 对任意概形 $U$、\'etale 态射 $a : U \to X$ 及满足
$a(u) = x$ 的 $u \in U$，局部环 $\mathcal{O}_{U, u}$ 有唯一极小素理想；
\item 对任意概形 $U$、\'etale 态射 $a : U \to X$ 及满足
$a(u) = x$ 的 $u \in U$，恰有一个 $U$ 的不可约分支通过 $u$；
\item 对任意概形 $U$、\'etale 态射 $a : U \to X$ 及满足
$a(u) = x$ 的 $u \in U$，局部环 $\mathcal{O}_{U, u}$ 单枝；
\item 对任意概形 $U$、\'etale 态射 $a : U \to X$ 及满足
$a(u) = x$ 的 $u \in U$，局部环 $\mathcal{O}_{U, u}$ 几何单枝；
\item 对位于 $x$ 上方的任意几何点 $\overline{x}$，
$\mathcal{O}_{X, \overline{x}}$ 有唯一极小素理想。
\end{enumerate}
\end{lemma}

\begin{proof}
(1) 与 (2) 等价，因为通过 $u$ 的 $U$ 的不可约分支与 $U$ 在 $u$ 处局部环的
极小素理想成 $1$-$1$ 对应。取 (1) 中的 $a : U \to X$ 与 $u \in U$。
由引理 \ref{spaces-properties-lemma-describe-etale-local-ring}，$\mathcal{O}_{X, \overline{x}}$
是 $\mathcal{O}_{U, u}$ 的严格 Hensel 化。特别地，由“更多代数”引理
\ref{more-algebra-lemma-geometrically-unibranch}，(4) 与 (5) 等价。
由“更多态射”引理 \ref{more-morphisms-lemma-nr-branches}，
(2)、(3)、(4) 等价。
\end{proof}

\begin{definition}
\label{spaces-properties-definition-unibranch}
设 $S$ 为概形，$X$ 为 $S$ 上的代数空间，并取 $x \in |X|$。
若引理 \ref{spaces-properties-lemma-irreducible-local-ring} 中的等价条件成立，
则称 $X$ 在 $x$ 处{\it 几何单枝}。若 $X$ 在每个
$x \in |X|$ 处都几何单枝，则称 $X${\it 几何单枝}。
\end{definition}

\noindent
这与概形的定义相容（“性质”定义 \ref{properties-definition-unibranch}）。

\begin{lemma}
\label{spaces-properties-lemma-nr-branches-local-ring}
设 $S$ 为概形，$X$ 为 $S$ 上的代数空间，取点 $x \in |X|$，
并取整数 $n \in \{1, 2, \ldots\}$。下列条件等价：
\begin{enumerate}
\item 对任意概形 $U$、\'etale 态射 $a : U \to X$ 及满足
$a(u) = x$ 的 $u \in U$，局部环 $\mathcal{O}_{U, u}$ 的极小素理想数
为 $\leq n$，且至少对一组选取 $U, a, u$，该数为 $n$；
\item 对任意概形 $U$、\'etale 态射 $a : U \to X$ 及满足
$a(u) = x$ 的 $u \in U$，通过 $u$ 的 $U$ 的不可约分支数为
$\leq n$，且至少对一组选取 $U, a, u$，该数为 $n$；
\item 对任意概形 $U$、\'etale 态射 $a : U \to X$ 及满足
$a(u) = x$ 的 $u \in U$，$U$ 在 $u$ 处的分支数为 $\leq n$，
且至少对一组选取 $U, a, u$，该数为 $n$；
\item 对任意概形 $U$、\'etale 态射 $a : U \to X$ 及满足
$a(u) = x$ 的 $u \in U$，$U$ 在 $u$ 处的几何分支数为 $n$；
\item $\mathcal{O}_{X, \overline{x}}$ 的极小素理想数为 $n$。
\end{enumerate}
\end{lemma}

\begin{proof}
(1) 与 (2) 等价，因为通过 $u$ 的 $U$ 的不可约分支与 $U$ 在 $u$ 处局部环的
极小素理想成 $1$-$1$ 对应。取 (1) 中的 $a : U \to X$ 与 $u \in U$。
由引理 \ref{spaces-properties-lemma-describe-etale-local-ring}，$\mathcal{O}_{X, \overline{x}}$
是 $\mathcal{O}_{U, u}$ 的严格 Hensel 化。回忆，$U$ 在 $u$ 处的
（几何）分支数，就是 $\mathcal{O}_{U, u}$ 的（严格）Hensel 化之极小素理想数。
特别地，(4) 与 (5) 等价。由“更多态射”引理
\ref{more-morphisms-lemma-nr-branches}，(2)、(3)、(4) 等价。
\end{proof}

\begin{definition}
\label{spaces-properties-definition-number-of-geometric-branches}
设 $S$ 为概形，$X$ 为 $S$ 上的代数空间，并取 $x \in |X|$。
若引理 \ref{spaces-properties-lemma-nr-branches-local-ring} 的等价条件成立，则
{\it $X$ 在 $x$ 处的几何分支数}为 $n \in \mathbf{N}$；否则为 $\infty$。
\end{definition}








\section{Noether 代数空间}
\label{spaces-properties-section-noetherian}

\noindent
我们已在第 \ref{spaces-properties-section-types-properties} 节定义了局部 Noether 代数空间。

\begin{definition}
\label{spaces-properties-definition-noetherian}
设 $S$ 为概形，$X$ 为 $S$ 上的代数空间。若 $X$ 拟紧、拟分离且局部
Noether，则称 $X$ 是{\it Noether}的。
\end{definition}

\noindent
注意，Noether 代数空间 $X$ 不仅拟紧且局部 Noether，还须拟分离。
这与 Noether 概形的定义并不冲突，因为局部 Noether 概形必拟分离；
见“性质”引理 \ref{properties-lemma-locally-Noetherian-quasi-separated}。
对代数空间则未必如此：$X = \mathbf{A}^1_k/\mathbf{Z}$（见“空间”例
\ref{spaces-example-affine-line-translation}）局部 Noether 且拟紧，
但不拟分离（因而按我们的定义不是 Noether 的）。

\medskip\noindent
上述选择的一个后果是：Noether 代数空间上的有限型代数空间未必自动 Noether；
也就是说，“态射”引理 \ref{morphisms-lemma-finite-type-noetherian}
的类似命题不成立。正确陈述是，Noether 代数空间上的有限表示代数空间是
Noether 的（见“空间的态射”引理
\ref{spaces-morphisms-lemma-finite-presentation-noetherian}）。

\medskip\noindent
Noether 代数空间 $X$ 非常接近概形。本节余下部分收集若干引理来说明这一点。

\begin{lemma}
\label{spaces-properties-lemma-Noetherian-topology}
设 $S$ 为概形，$X$ 为 $S$ 上的代数空间。
\begin{enumerate}
\item 若 $X$ 局部 Noether，则 $|X|$ 是局部 Noether 拓扑空间。
\item 若 $X$ 拟紧且局部 Noether，则 $|X|$ 是 Noether 拓扑空间。
\end{enumerate}
\end{lemma}

\begin{proof}
设 $X$ 局部 Noether。选取概形 $U$ 及满的 \'etale 态射 $U \to X$。
由 $X$ 局部 Noether 可知 $U$ 局部 Noether。根据“性质”引理
\ref{properties-lemma-Noetherian-topology}，这意味着 $|U|$ 是局部
Noether 拓扑空间。由于 $|U| \to |X|$ 是开满射，“拓扑”引理
\ref{topology-lemma-image-Noetherian} 表明 $|X|$ 局部 Noether。
这证明了 (1)。若 $X$ 拟紧且局部 Noether，则 $|X|$ 拟紧且局部
Noether，故由“拓扑”引理
\ref{topology-lemma-quasi-compact-locally-Noetherian-Noetherian}，
$|X|$ 是 Noether 的。
\end{proof}

\begin{lemma}
\label{spaces-properties-lemma-Noetherian-sober}
设 $S$ 为概形，$X$ 为 $S$ 上的代数空间。若 $X$ Noether，
则 $|X|$ 是清醒 Noether 拓扑空间。
\end{lemma}

\begin{proof}
拟分离代数空间的底拓扑空间是清醒的，见引理
\ref{spaces-properties-lemma-quasi-separated-sober}；它由引理
\ref{spaces-properties-lemma-Noetherian-topology} 又是 Noether 的。
\end{proof}

\begin{lemma}
\label{spaces-properties-lemma-Noetherian-local-ring-Noetherian}
设 $S$ 为概形，$X$ 为 $S$ 上的 Noether 代数空间，并设
$\overline{x}$ 为 $X$ 的几何点。则 $\mathcal{O}_{X, \overline{x}}$
是 Noether 局部环。
\end{lemma}

\begin{proof}
选取 $\overline{x}$ 的 \'etale 邻域 $(U, \overline{u})$，其中 $U$ 为概形。
由引理 \ref{spaces-properties-lemma-describe-etale-local-ring}，
$\mathcal{O}_{X, \overline{x}}$ 是 $U$ 在 $u$ 处局部环的严格 Hensel 化。
按我们对 Noether 空间的定义，概形 $U$ 局部 Noether。故结论由“更多代数”引理
\ref{more-algebra-lemma-henselization-noetherian} 得到。
\end{proof}








\section{正则代数空间}
\label{spaces-properties-section-regular}

\noindent
我们已在第 \ref{spaces-properties-section-types-properties} 节定义了正则代数空间。

\begin{lemma}
\label{spaces-properties-lemma-regular}
设 $S$ 为概形，$X$ 为 $S$ 上局部 Noether 的代数空间。下列条件等价：
\begin{enumerate}
\item $X$ 正则；
\item 每个 \'etale 局部环 $\mathcal{O}_{X, \overline{x}}$ 都正则。
\end{enumerate}
\end{lemma}

\begin{proof}
设 $U$ 为概形，且 $U \to X$ 为满的 \'etale 态射。依假设，$U$ 局部
Noether。此外，每个 \'etale 局部环 $\mathcal{O}_{X, \overline{x}}$
都是 $U$ 上某个局部环的严格 Hensel 化，反之亦然；见引理
\ref{spaces-properties-lemma-describe-etale-local-ring}。故由“更多代数”引理
\ref{more-algebra-lemma-henselization-regular}，(2) 等价于 $U$ 的每个局部环
都正则，即 $U$ 为正则概形（见“性质”引理
\ref{properties-lemma-characterize-regular}）。由定义
\ref{spaces-properties-definition-type-property}，这又等价于 (1)。
\end{proof}

\noindent
可利用“下降”引理 \ref{descent-lemma-regular-local-ring-local}
定义代数空间 $X$ 在点 $x$ 处正则的含义。

\begin{definition}
\label{spaces-properties-definition-regular-at-point}
设 $S$ 为概形，$X$ 为 $S$ 上的代数空间，并取点 $x \in |X|$。
若对任意（等价地，某个）二元组 $(a : U \to X, u)$，
$\mathcal{O}_{U, u}$ 都是正则局部环，则称{\it $X$ 在 $x$ 处正则}；
这里 $a : U \to X$ 是从概形到 $X$ 的 \'etale 态射，且点
$u \in U$ 满足 $a(u) = x$。
\end{definition}

\noindent
见定义 \ref{spaces-properties-definition-property-at-point}、引理
\ref{spaces-properties-lemma-local-source-target-at-point} 以及“下降”引理
\ref{descent-lemma-regular-local-ring-local}。

\begin{lemma}
\label{spaces-properties-lemma-regular-at-x}
设 $S$ 为概形，$X$ 为 $S$ 上的代数空间，并取点 $x \in |X|$。
下列条件等价：
\begin{enumerate}
\item $X$ 在 $x$ 处正则；
\item 对位于 $x$ 上方的任意（等价地，某个）几何点 $\overline{x}$，
\'etale 局部环 $\mathcal{O}_{X, \overline{x}}$ 正则。
\end{enumerate}
\end{lemma}

\begin{proof}
设 $U$ 为概形，$u \in U$ 为点，并设 $a : U \to X$ 为把 $u$ 映到
$x$ 的 \'etale 态射。对 $X$ 中位于 $x$ 上方的任意几何点
$\overline{x}$，\'etale 局部环 $\mathcal{O}_{X, \overline{x}}$
是 $U$ 在 $u$ 处某个局部环的严格 Hensel 化，见引理
\ref{spaces-properties-lemma-describe-etale-local-ring}。故结论由“更多代数”引理
\ref{more-algebra-lemma-henselization-regular} 得到。
\end{proof}

\begin{lemma}
\label{spaces-properties-lemma-regular-normal}
正则代数空间是正规代数空间。
\end{lemma}

\begin{proof}
这由定义及概形的情形推出；见“性质”引理
\ref{properties-lemma-regular-normal}。
\end{proof}

 








\section{代数空间上的模层}
\label{spaces-properties-section-modules}

\noindent
若 $X$ 为代数空间，则 $X$ 上的模层是 $X$ 的小 \'etale 站点上的
$\mathcal{O}_X$-模层，其中 $\mathcal{O}_X$ 为 $X$ 的结构层。
模层范畴记作 $\textit{Mod}(\mathcal{O}_X)$。

\medskip\noindent
给定代数空间的态射 $f : X \to Y$，由引理
\ref{spaces-properties-lemma-morphism-ringed-topoi} 得到环化拓扑斯态射，进而由“站点上的模”定义
\ref{sites-modules-definition-pushforward} 得到良定义的拉回函子与直像函子
\begin{equation}
\label{spaces-properties-equation-push-pull}
f^* :
\textit{Mod}(\mathcal{O}_Y)
\longrightarrow
\textit{Mod}(\mathcal{O}_X), \quad
f_* :
\textit{Mod}(\mathcal{O}_X)
\longrightarrow
\textit{Mod}(\mathcal{O}_Y)
\end{equation}
，它们按通常方式互为伴随。若 $g : Y \to Z$ 是 $S$ 上代数空间的另一态射，
则
$(g \circ f)^* = f^* \circ g^*$，且 $(g \circ f)_* = g_* \circ f_*$
，这只因环化拓扑斯态射按相应方式复合（由上述引理）。

\begin{lemma}
\label{spaces-properties-lemma-etale-exact-pullback}
设 $S$ 为概形，$f : X \to Y$ 为 $S$ 上代数空间的 \'etale 态射。
则 $f^{-1}\mathcal{O}_Y = \mathcal{O}_X$，并且对任意
$\mathcal{O}_Y$-模层 $\mathcal{G}$，有
$f^*\mathcal{G} = f_{small}^{-1}\mathcal{G}$。特别地，
$f^* : \textit{Mod}(\mathcal{O}_Y) \to \textit{Mod}(\mathcal{O}_X)$
正合。
\end{lemma}

\begin{proof}
由引理 \ref{spaces-properties-lemma-etale-morphism-topoi} 中对逆像的描述及结构层的定义，
$f_{small}^{-1}\mathcal{O}_Y = \mathcal{O}_X$ 是清楚的。由于拉回按定义为
$$
f^*\mathcal{G} =
f_{small}^{-1}\mathcal{G} \otimes_{f_{small}^{-1}\mathcal{O}_Y}
\mathcal{O}_X
$$
，可得 $f^*\mathcal{G} = f_{small}^{-1}\mathcal{G}$。正合性也是清楚的，
因为 $f_{small}$ 是拓扑斯态射，所以 $f_{small}^{-1}$ 正合。
\end{proof}

\noindent
继续沿用式 (\ref{spaces-properties-equation-restrict}) 中引入的记号滥用，在上述引理的情形写成
\begin{equation}
\label{spaces-properties-equation-restrict-modules}
\mathcal{G}|_{X_\etale}
= f^*\mathcal{G}
= f_{small}^{-1}\mathcal{G}
\end{equation}
。第 \ref{spaces-properties-section-localize} 节将以更技术性的方式讨论此事。

\begin{lemma}
\label{spaces-properties-lemma-pushforward-etale-base-change-modules}
设 $S$ 为概形，并设
$$
\xymatrix{
X' \ar[r] \ar[d]_{f'} & X \ar[d]^f \\
Y' \ar[r]^g & Y
}
$$
为 $S$ 上代数空间的笛卡儿方块。取
$\mathcal{F} \in \textit{Mod}(\mathcal{O}_X)$。若 $g$ 是 \'etale 的，则
$f'_*(\mathcal{F}|_{X'}) = (f_*\mathcal{F})|_{Y'}$\footnote{还有
$(f')^*(\mathcal{G}|_{Y'}) = (f^*\mathcal{G})|_{X'}$
，这是因为该图交换且有式 (\ref{spaces-properties-equation-restrict-modules})}，并且
$R^if'_*(\mathcal{F}|_{X'}) = (R^if_*\mathcal{F})|_{Y'}$ 在
$\textit{Mod}(\mathcal{O}_{Y'})$ 中成立。
\end{lemma}

\begin{proof}
这是引理 \ref{spaces-properties-lemma-pushforward-etale-base-change} 在模情形的重新表述。
\end{proof}

\begin{lemma}
\label{spaces-properties-lemma-characterize-module-small-etale}
设 $S$ 为概形，$X$ 为 $S$ 上的代数空间。一个 $\mathcal{O}_X$-模层
$\mathcal{F}$ 由下列数据给出：
\begin{enumerate}
\item 对每个 $U \in \Ob(X_\etale)$，给定 $U_\etale$ 上的
$\mathcal{O}_U$-模层 $\mathcal{F}_U$；
\item 对 $X_\etale$ 中每个 $f : U' \to U$，给定同构
$c_f : f_{small}^*\mathcal{F}_U \to \mathcal{F}_{U'}$.
\end{enumerate}
这些数据须满足如下条件：对 $X_\etale$ 中任意 $f : U' \to U$ 与
$g : U'' \to U'$，复合 $c_g \circ g_{small}^*c_f$ 等于
$c_{f \circ g}$。
\end{lemma}

\begin{proof}
合用引理 \ref{spaces-properties-lemma-etale-exact-pullback} 与
\ref{spaces-properties-lemma-characterize-sheaf-small-etale}，并使用如下事实：
$X_\etale$ 的对象之间任意态射都是概形的 \'etale 态射。
\end{proof}







\section{\'Etale 局部化}
\label{spaces-properties-section-localize}

\noindent
应不惜一切代价避免阅读本节。

\medskip\noindent
设 $X \to Y$ 为代数空间的 \'etale 态射。则 $X$ 是
$Y_{spaces, \etale}$ 的对象；由定义立即可得（也见引理
\ref{spaces-properties-lemma-etale-morphism-topoi} 的证明）
\begin{equation}
\label{spaces-properties-equation-localize}
X_{spaces, \etale} = Y_{spaces, \etale}/X
\end{equation}
其中右端是站点 $Y_{spaces, \etale}$ 在对象 $X$ 处的局部化，
见“站点”定义 \ref{sites-definition-localize}。此外，由引理
\ref{spaces-properties-lemma-etale-exact-pullback}，此等同与结构层相容。
因此环化站点 $(X_{spaces, \etale}, \mathcal{O}_X)$ 等同于环化站点
$(Y_{spaces, \etale}, \mathcal{O}_Y)$ 在对象 $X$ 处的局部化：
\begin{equation}
\label{spaces-properties-equation-localize-ringed}
(X_{spaces, \etale}, \mathcal{O}_X) =
(Y_{spaces, \etale}/X, \mathcal{O}_Y|_{Y_{spaces, \etale}/X})
\end{equation}
右端所用的环化站点局部化定义于“站点上的模”定义
\ref{sites-modules-definition-localize-ringed-site}。

\medskip\noindent
现在设 $X \to Y$ 是代数空间的 \'etale 态射，且 $X$ 为概形。
则 $X$ 是 $Y_\etale$ 的对象，从而
\begin{equation}
\label{spaces-properties-equation-localize-at-scheme}
X_\etale = Y_\etale/X
\end{equation}
以及
\begin{equation}
\label{spaces-properties-equation-localize-at-scheme-ringed}
(X_\etale, \mathcal{O}_X) =
(Y_\etale/X, \mathcal{O}_Y|_{Y_\etale/X})
\end{equation}
，与上面相同。

\medskip\noindent
最后，若 $X \to Y$ 是代数空间的 \'etale 态射，且 $X$ 为仿射概形，
则 $X$ 是 $Y_{affine, \etale}$ 的对象，并有
\begin{equation}
\label{spaces-properties-equation-localize-at-affine}
X_{affine, \etale} = Y_{affine, \etale}/X
\end{equation}
以及
\begin{equation}
\label{spaces-properties-equation-localize-at-affine-ringed}
(X_{affine, \etale}, \mathcal{O}_X) =
(Y_{affine, \etale}/X, \mathcal{O}_Y|_{Y_{affine, \etale}/X})
\end{equation}
，与上面相同。

\medskip\noindent
下面证明这些局部化与态射相容。

\begin{lemma}
\label{spaces-properties-lemma-relocalize-morphism}
设 $S$ 为概形，并设
$$
\xymatrix{
U \ar[d]_p \ar[r]_g & V \ar[d]^q \\
X \ar[r]^f & Y
}
$$
为 $S$ 上代数空间的交换图，其中 $p$ 与 $q$ 都是 \'etale 的。
通过式 (\ref{spaces-properties-equation-localize-ringed}) 分别对 $U \to X$ 与
$V \to Y$ 给出的等同，环化拓扑斯态射
$$
(g_{spaces, \etale}, g^\sharp) :
(\Sh(U_{spaces, \etale}), \mathcal{O}_U)
\longrightarrow
(\Sh(V_{spaces, \etale}), \mathcal{O}_V)
$$
与下述态射 $2$-同构：从环化站点态射
$(f_{spaces, \etale}, f^\sharp)$ 以及对应于 $g$ 的映射
$c : U \to V \times_Y X$ 出发，按“站点上的模”引理
\ref{sites-modules-lemma-relocalize-morphism-ringed-sites} 构造的态射
$(f_{spaces, \etale, c}, f_c^\sharp)$。
\end{lemma}

\begin{proof}
态射 $(f_{spaces, \etale, c}, f_c^\sharp)$ 定义为局部化与基变换映射的复合
$f' \circ j$。类似地，$g$ 是复合 $U \to V \times_Y X \to V$。
故只须在下列两种情形证明本引理：(1) $f = \text{id}$；
(2) $U = X \times_Y V$。在情形 (1)，态射 $g : U \to V$ 是 \'etale 的，
见引理 \ref{spaces-properties-lemma-etale-permanence}。于是由式
(\ref{spaces-properties-equation-localize}) 与 (\ref{spaces-properties-equation-localize-ringed}) 周围的讨论，
$(g_{spaces, \etale}, g^\sharp)$ 是局部化态射；这恰是本引理在该情形的内容。
在情形 (2)，态射 $g_{spaces, \etale}$ 来自函子
$V_{spaces, \etale} \to U_{spaces, \etale}$,
$V'/V \mapsto V' \times_V U/U$
所给的环化站点态射，而态射 $f'$ 也是由此定义的；见“站点”引理
\ref{sites-lemma-localize-morphism}。略去在此情形验证
$(f')^\sharp = g^\sharp$（二者都是 $f^\sharp$ 到
$U_{spaces, \etale}$ 的限制）。
\end{proof}

\begin{lemma}
\label{spaces-properties-lemma-relocalize-morphism-at-schemes}
采用与引理 \ref{spaces-properties-lemma-relocalize-morphism} 相同的记号及假设，
但另假设 $U$ 与 $V$ 都是概形。通过式
(\ref{spaces-properties-equation-localize-at-scheme-ringed}) 分别对 $U \to X$ 与
$V \to Y$ 给出的等同，环化拓扑斯态射
$$
(g_{small}, g^\sharp) :
(\Sh(U_\etale), \mathcal{O}_U)
\longrightarrow
(\Sh(V_\etale), \mathcal{O}_V)
$$
与下述态射 $2$-同构：从 $(f_{small}, f^\sharp)$ 及对应于 $g$ 的映射
$s : h_U \to f_{small}^{-1}h_V$ 出发，按“站点上的模”引理
\ref{sites-modules-lemma-relocalize-morphism-ringed-topoi} 构造的态射
$(f_{small, s}, f_s^\sharp)$。
\end{lemma}

\begin{proof}
注意，作为环化拓扑斯态射，$(g_{small}, g^\sharp)$ 与相伴于环化站点态射
$(g_{spaces, \etale}, g^\sharp)$ 的环化拓扑斯态射 $2$-同构。
故结论由引理 \ref{spaces-properties-lemma-relocalize-morphism} 及“站点上的模”引理
\ref{sites-modules-lemma-relocalize-morphism-compare} 得到。
\end{proof}

\noindent
最后讨论代数空间 $Y$ 的小 \'etale 站点 $Y_\etale$ 上的集合层，
与 $Y$ 上 \'etale 的代数空间之间的关系。设 $S$ 为概形，$Y$ 为
$S$ 上的代数空间，且 $\mathcal{F}$ 为 $\Sh(Y_\etale)$ 的对象。
考虑函子
$$
X : (\Sch/S)_{fppf}^{opp} \longrightarrow \textit{Sets}
$$
，其定义规则为
$$
X(T) = \{(y, s) \mid y : T \to Y\text{ 是 }S\text{ 上的态射且 }
s \in \Gamma(T, y_{small}^{-1}\mathcal{F})\}
$$
给定态射 $g : T' \to T$，限制映射把 $(y, s)$ 送到
$(y \circ g, g_{small}^{-1}s)$。这是有意义的，因为由引理
\ref{spaces-properties-lemma-functoriality-etale-site}，
$y_{small} \circ g_{small} = (y \circ g)_{small}$。
存在典范映射 $X \to Y$，它把二元组 $(y, s)$ 送到 $y$。

\begin{lemma}
\label{spaces-properties-lemma-sheaf-gives-space}
设 $S$ 为概形，$Y$ 为 $S$ 上的代数空间，并设 $\mathcal{F}$
为 $Y_\etale$ 上的集合层。只要满足一个集合论条件（见证明），就有：
\begin{enumerate}
\item 上述与 $\mathcal{F}$ 相伴的函子 $X$ 是代数空间；
\item 映射 $X \to Y$ 是代数空间的 \'etale 态射；
\item 通过等同 $\Sh(Y_\etale) = \Sh(Y_{spaces, \etale})$，有
$\mathcal{F} \cong h_X$；
\item 有 $\mathcal{F} \cong f_{small, !}*$。这里 $*$ 是范畴
$\Sh(X_\etale)$ 的终对象，而 $f_{small, !}$ 由引理
\ref{spaces-properties-lemma-etale-morphism-topoi} 存在。
\end{enumerate}
\end{lemma}

\begin{proof}
先证明 $X$ 是 fppf 拓扑的层。设 $\{g_i : T_i \to T\}$ 是
$(\Sch/S)_{fppf}$ 的覆盖，且 $(y_i, s_i) \in X(T_i)$ 满足胶合条件，
即 $(y_i, s_i)$ 与 $(y_j, s_j)$ 到 $T_i \times_T T_j$ 的限制相同。
由于 $Y$ 是 fppf 拓扑的层，各 $y_i$ 胶合为唯一态射
$y : T \to Y$，使得 $y_i = y \circ g_i$。于是
$y_{i, small}^{-1}\mathcal{F} = g_{i, small}^{-1}y_{small}^{-1}\mathcal{F}$.
故由“\'Etale 上同调”引理
\ref{etale-cohomology-lemma-describe-pullback}，各截面 $s_i$
唯一胶合为 $y_{small}^{-1}\mathcal{F}$ 的一个截面。

\medskip\noindent
把 $\mathcal{F} \in \Ob(\Sh(Y_\etale))$ 送到
$X \in \Ob((\Sch/S)_{fppf})$ 的构造保持有限极限及所有余极限，
因为每个函子 $y_{small}^{-1}$ 都有此性质。当然，若
$V \in \Ob(Y_\etale)$，则该构造把 $Y_\etale$ 上的可表层 $h_V$
送到由 $V$ 表示的可表函子。

\medskip\noindent
由“站点”引理 \ref{sites-lemma-sheaf-coequalizer-representable}，
可找到集合 $I$、对每个 $i \in I$ 找到 $Y_\etale$ 的对象 $V_i$，
以及层的满射
$$
\coprod h_{V_i} \longrightarrow \mathcal{F}
$$
于 $Y_\etale$ 上。我们需要的集合论条件是指标集 $I$ 不太大\footnote{只须
$\mathcal{F}$ 在 $Y$ 的各几何点处之茎的基数上确界，不超过
$(\Sch/S)_{fppf}$ 某个对象的大小。}。于是
$V = \coprod V_i$ 是 $(\Sch/S)_{fppf}$ 的对象，从而也是
$Y_\etale$ 的对象，并且有满射 $h_V \to \mathcal{F}$。

\medskip\noindent
注意，$h_V$ 在 $\Sh(Y_\etale)$ 中与自身的积是
$h_{V \times_Y V}$。考虑纤维积
$$
h_V \times_\mathcal{F} h_V \subset h_{V \times_Y V}
$$
存在 $V \times_Y V$ 的开子概形 $R$，使得
$h_V \times_\mathcal{F} h_V = h_R$；见引理
\ref{spaces-properties-lemma-support-subsheaf-final}（略去一个小细节）。由 Yoneda 引理，
得到 $Y_\etale$ 中两个态射 $s, t : R \to V$，并得到余等化子图
$$
\xymatrix{
h_R \ar@<1ex>[r] \ar@<-1ex>[r] &
h_V \ar[r] &
\mathcal{F}
}
$$
于 $\Sh(Y_\etale)$ 中。当然，态射 $s, t$ 都是 \'etale 的，
并定义一个 \'etale 等价关系 $(t, s) : R \to V \times_S V$。

\medskip\noindent
由前两段的讨论，得到余等化子图
$$
\xymatrix{
R \ar@<1ex>[r] \ar@<-1ex>[r] &
V \ar[r] &
X
}
$$
于 $(\Sch/S)_{fppf}$ 中。因此，由“空间”定理
\ref{spaces-theorem-presentation}，$X = V/R$ 是代数空间。这证明了 (1)。
(2) 来自 $V \to Y$ 是 \'etale 的。(3) 由 $X$ 与 $h_X$ 的定义立即可得。
略去 (4) 的证明；它可由下述比较得到：一边是与引理
\ref{spaces-properties-lemma-etale-morphism-topoi} 的余连续函子 $j$ 相伴的态射，
另一边是上面对 $X_{spaces, \etale}$ 作为 $Y_{spaces, \etale}$
在 $X$ 处之局部化的描述以及“站点”引理
\ref{sites-lemma-describe-j-shriek-representable}。
\end{proof}







\section{恢复态射}
\label{spaces-properties-section-morphisms}

\noindent
本节证明，把一个代数空间与其局部环化的小 \'etale 拓扑斯对应起来的规则，
在适当意义下是全忠实的，参见
定理 \ref{spaces-properties-theorem-fully-faithful}。

\begin{lemma}
\label{spaces-properties-lemma-morphism-locally-ringed}
设 $S$ 为概形。
设 $f : X \to Y$ 为 $S$ 上代数空间的态射。
与 $f$ 关联的环化拓扑斯态射 $(f_{small}, f^\sharp)$
是局部环化拓扑斯的态射，参见
《站点上的模》，
定义 \ref{sites-modules-definition-morphism-locally-ringed-topoi}。
\end{lemma}

\begin{proof}
注意该断言是有意义的，因为我们已经看到
$(X_\etale, \mathcal{O}_{X_\etale})$ 以及
$(Y_\etale, \mathcal{O}_{Y_\etale})$
均为局部环化站点，参见
引理 \ref{spaces-properties-lemma-etale-site-locally-ringed}。
此外，我们知道 $X_\etale$ 有足够多的点，参见
定理 \ref{spaces-properties-theorem-exactness-stalks}。
因此只需证明 $(f_{small}, f^\sharp)$
满足《站点上的模》
引理 \ref{sites-modules-lemma-locally-ringed-morphism}
的条件 (3)。为此，取 $X_\etale$ 的一个点 $p$。由
引理 \ref{spaces-properties-lemma-points-small-etale-site}，
$p$ 对应于 $X$ 的一个几何点 $\overline{x}$。
由
引理 \ref{spaces-properties-lemma-stalk-pullback}，
点 $q = f_{small} \circ p$ 对应于 $Y$ 的几何点
$\overline{y} = f \circ \overline{x}$。
故我们需要证明的断言是，所诱导的 \'etale 局部环同态
$$
\mathcal{O}_{Y, \overline{y}} \longrightarrow \mathcal{O}_{X, \overline{x}}
$$
是局部环同态。可以直接证明这一点；这里我们改由概形的相应结果推出。
为此选择交换图
$$
\xymatrix{
U \ar[d] \ar[r]_\psi & V \ar[d] \\
X \ar[r] & Y
}
$$
其中 $U$ 和 $V$ 为概形，竖直箭头为满的 \'etale 态射（参见
《空间》，引理 \ref{spaces-lemma-lift-morphism-presentations}）。
选择一个提升 $\overline{u} : \overline{x} \to U$（由
引理 \ref{spaces-properties-lemma-geometric-lift-to-cover} 可知这是可能的）。
令 $\overline{v} = \psi \circ \overline{u}$。于是得到 \'etale 局部环的交换图
$$
\xymatrix{
\mathcal{O}_{U, \overline{u}} &
\mathcal{O}_{V, \overline{v}} \ar[l] \\
\mathcal{O}_{X, \overline{x}} \ar[u] &
\mathcal{O}_{Y, \overline{y}}. \ar[l] \ar[u]
}
$$
由
《\'Etale 上同调》，引理 \ref{etale-cohomology-lemma-morphism-locally-ringed}，
上方的水平箭头是局部环同态。最后，由
引理 \ref{spaces-properties-lemma-describe-etale-local-ring}，
竖直箭头均为同构。结论随即成立。
\end{proof}

\begin{lemma}
\label{spaces-properties-lemma-2-morphism}
设 $S$ 为概形。
设 $X$、$Y$ 为 $S$ 上的代数空间。
设 $f : X \to Y$ 为 $S$ 上代数空间的态射。
设 $t$ 是从 $(f_{small}, f^\sharp)$ 到其自身的一个 $2$-态射，参见
《站点上的模》，
定义 \ref{sites-modules-definition-2-morphism-ringed-topoi}。
则 $t = \text{id}$。
\end{lemma}

\begin{proof}
令 $X'$（分别地，$Y'$）表示按同样方式把 $X$ 视为
$\Spec(\mathbf{Z})$ 上的代数空间，参见
《空间》，定义 \ref{spaces-definition-base-change}。
由构造可知 $(X_{small}, \mathcal{O})$
等于 $(X'_{small}, \mathcal{O})$，对 $Y$ 也同样成立。
因此可以处理 $X'$ 和 $Y'$。换言之，我们可以假设
$S = \Spec(\mathbf{Z})$。

\medskip\noindent
假设 $S = \Spec(\mathbf{Z})$，且 $f : X \to Y$ 和 $t$ 如引理中所述。
这意味着 $t : f^{-1}_{small} \to f^{-1}_{small}$
是函子的一个变换，并且下图交换：
$$
\xymatrix{
f_{small}^{-1}\mathcal{O}_Y
\ar[rd]_{f^\sharp}  & &
f_{small}^{-1}\mathcal{O}_Y \ar[ll]^t \ar[ld]^{f^\sharp} \\
& \mathcal{O}_X
}
$$
设 $V \to Y$ 为 \'etale 态射且 $V$ 为仿射概形。
写作 $V = \Spec(B)$。选取 $b_j \in B$（$j \in J$）作为
$B$ 这个 $\mathbf{Z}$-代数的一组生成元。令
$T = \Spec(\mathbf{Z}[\{x_j\}_{j \in J}])$.
下文中我们将使用
$\Mor_{\Sch}(U, T) = \prod_{j \in J} \Gamma(U, \mathcal{O}_U)$
对任意概形 $U$ 都成立，而不再另行说明。
满环同态 $\mathbf{Z}[x_j] \to B$（$x_j \mapsto b_j$）
对应于闭浸入 $V \to T$。
于是得到一个单态射
$$
i : V \longrightarrow T_Y = T \times Y
$$
这是 $Y$ 上代数空间的单态射。就 $Y_\etale$ 上的层而言，
态射 $i$ 诱导层的单射
$h_i : h_V \to \prod_{j \in J} \mathcal{O}_Y$。
$i$ 到 $X$ 的基变换 $i' : X \times_Y V \to T_X$
也为单态射（《空间》，
引理 \ref{spaces-lemma-base-change-representable-transformations-property}）。
因此 $i' : X \times_Y V \to T_X$ 为单态射，这又意味着
$h_{i'} : h_{X \times_Y V} \to \prod_{j \in J} \mathcal{O}_X$
是层的单射。
通过引理 \ref{spaces-properties-lemma-stalk-pullback} 给出的恒等
$f_{small}^{-1}h_V = h_{X \times_Y V}$，
映射 $h_{i'}$ 等于
$$
\xymatrix{
f_{small}^{-1}h_V \ar[r]^-{f^{-1}h_i} &
\prod_{j \in J} f_{small}^{-1}\mathcal{O}_Y
\ar[r]^{\prod f^\sharp} &
\prod_{j \in J} \mathcal{O}_X
}
$$
(验证略去）。这意味着映射
$t : f_{small}^{-1}h_V \to f_{small}^{-1}h_V$
适合于下列交换图：
$$
\xymatrix{
f_{small}^{-1}h_V \ar[r]^-{f^{-1}h_i} \ar[d]^t &
\prod_{j \in J} f_{small}^{-1}\mathcal{O}_Y
\ar[r]^-{\prod f^\sharp} \ar[d]^{\prod t} &
\prod_{j \in J} \mathcal{O}_X \ar[d]^{\text{id}}\\
f_{small}^{-1}h_V \ar[r]^-{f^{-1}h_i} &
\prod_{j \in J} f_{small}^{-1}\mathcal{O}_Y
\ar[r]^-{\prod f^\sharp} &
\prod_{j \in J} \mathcal{O}_X
}
$$
右方方块的交换性由上文解释的关于 $t$ 的假设保证。
由于上面讨论的水平箭头复合是单射，我们也得到左方竖直箭头
是恒等映射。$Y_\etale$ 上的任意集合层都可以由一族形如
$h_V$（其中 $V$ 为仿射概形）的层的一个（很大的）余积满射到它
（结合引理 \ref{spaces-properties-lemma-alternative} 与
《站点》，引理 \ref{sites-lemma-sheaf-coequalizer-representable}）。
因此得到 $t : f_{small}^{-1} \to f_{small}^{-1}$
是恒等变换，正如所需。
\end{proof}

\begin{lemma}
\label{spaces-properties-lemma-faithful}
设 $S$ 为概形。
设 $X$、$Y$ 为 $S$ 上的代数空间。
对于 $S$ 上代数空间的任意两个态射 $a, b : X \to Y$，
若存在一个 $2$-同构
$(a_{small}, a^\sharp) \cong (b_{small}, b^\sharp)$
使得它们在环化拓扑斯的 $2$-范畴中相同，则二者相等。
\end{lemma}

\begin{proof}
设 $t : a_{small}^{-1} \to b_{small}^{-1}$ 为该 $2$-同构。
等价地，可以把 $t$ 看作变换
$t : a_{spaces, \etale}^{-1} \to b_{spaces, \etale}^{-1}$
因为 $X_\etale$ 上的层与 $X_{spaces, \etale}$ 上的层没有区别。
选取交换图
$$
\xymatrix{
U \ar[d]_p \ar[r]_\alpha  & V \ar[d]^q \\
X \ar[r]^a & Y
}
$$
其中 $U$、$V$ 为概形，且 $p$、$q$ 为满的 \'etale 态射。
考虑图
$$
\xymatrix{
h_U \ar[r]_-\alpha \ar@{=}[d] & a_{spaces, \etale}^{-1}h_V \ar[d]^t \\
h_U \ar@{..>}[r] & b_{spaces, \etale}^{-1}h_V
}
$$
由于层 $b_{spaces, \etale}^{-1}h_V$ 同构于
$h_{V \times_{Y, b} X}$，可知虚线箭头来自一个概形态射
$\beta : U \to V$，并适合交换图
$$
\xymatrix{
U \ar[d]_p \ar[r]_\beta  & V \ar[d]^q \\
X \ar[r]^b & Y
}
$$
我们断言存在如下 $2$-同构序列：
\begin{align*}
(\alpha_{small}, \alpha^\sharp)
& \cong
(\alpha_{spaces, \etale}, \alpha^\sharp) \\
& \cong
(a_{spaces, \etale, c}, a_c^\sharp) \\
& \cong
(b_{spaces, \etale, d}, b_d^\sharp) \\
& \cong
(\beta_{spaces, \etale}, \beta^\sharp) \\
& \cong
(\beta_{small}, \beta^\sharp)
\end{align*}
第一个和最后一个 $2$-同构来自
$U_{spaces, \etale}$ 上的层与 $U_\etale$ 上的层之间的恒等，
对 $V$ 也一样。第二个和第四个
$2$-同构来自
引理 \ref{spaces-properties-lemma-relocalize-morphism}，
其中 $c : U \to X \times_{a, Y} V$ 由 $\alpha$ 诱导，
$d : U \to X \times_{b, Y} V$ 由 $\beta$ 诱导。
中间的 $2$-同构来自变换 $t$。具体地，函子
$a_{spaces, \etale, c}^{-1}$ 对应于函子
$$
(\mathcal{H} \to h_V) \longmapsto
(a_{spaces, \etale}^{-1}\mathcal{H}
\times_{a_{spaces, \etale}^{-1}h_V, \alpha}
h_U \to
h_U)
$$
对 $b_{spaces, \etale, d}^{-1}$ 也同样，参见
《站点》，引理 \ref{sites-lemma-relocalize-morphism}。
这里使用了如下恒等：$Y_{spaces, \etale}/V$ 上的层
对应于 $\Sh(Y_{spaces, \etale})$ 中的箭头
$(\mathcal{H} \to h_V)$，对 $U/X$ 也一样；参见
《站点》，引理 \ref{sites-lemma-essential-image-j-shriek}。
在此恒等下，结构层 $\mathcal{O}_V$ 对应于
对 $(\mathcal{O}_Y \times h_V \to h_V)$，而 $\mathcal{O}_U$ 也类似；参见
《站点上的模》，
引理 \ref{sites-modules-lemma-localize-compare}。
由于 $t$ 交换 $\alpha$ 与 $\beta$，可知 $t$ 诱导同构
$$
t :
a_{spaces, \etale}^{-1}\mathcal{H}
\times_{a_{spaces, \etale}^{-1}h_V, \alpha}
h_U
\longrightarrow
b_{spaces, \etale}^{-1}\mathcal{H}
\times_{b_{spaces, \etale}^{-1}h_V, \beta}
h_U
$$
它在 $h_U$ 上，并且关于 $(\mathcal{H} \to h_V)$ 具有函子性。
此外，$t$ 与 $a_c^\sharp$、$b_d^\sharp$ 相容，因为按照上面对
结构层 $\mathcal{O}_U$、$\mathcal{O}_V$ 的描述，$t$ 与
$a^\sharp$、$b^\sharp$ 相容。因此，环化拓扑斯态射
$(\alpha_{small}, \alpha^\sharp)$ 和 $(\beta_{small}, \beta^\sharp)$
是 $2$-同构的。由
《\'Etale 上同调》，引理 \ref{etale-cohomology-lemma-faithful}，
得到 $\alpha = \beta$！由于 $p : U \to X$ 是层的满射，遂有 $a=b$。
\end{proof}

\noindent
下面是本节的主要结果。

\begin{theorem}
\label{spaces-properties-theorem-fully-faithful}
设 $X$、$Y$ 为 $\Spec(\mathbf{Z})$ 上的代数空间。
设
$$
(g, g^\sharp) :
(\Sh(X_\etale), \mathcal{O}_X)
\longrightarrow
(\Sh(Y_\etale), \mathcal{O}_Y)
$$
为局部环化拓扑斯的态射。则存在唯一的代数空间态射
$f : X \to Y$，使得
$(g, g^\sharp)$ 与 $(f_{small}, f^\sharp)$ 同构。
换言之，构造
$$
\textit{Spaces}/\Spec(\mathbf{Z})
\longrightarrow \textit{Locally ringed topoi},
\quad
X \longrightarrow (X_\etale, \mathcal{O}_X)
$$
是全忠实的（右侧的态射取至 $2$-同构）。
\end{theorem}

\begin{proof}
唯一性已在引理 \ref{spaces-properties-lemma-faithful} 中证明。
因此只需证明存在性。
在本证明中，我们将自由使用等式
（\ref{spaces-properties-equation-localize-at-scheme-ringed}）中的恒等，
以及引理 \ref{spaces-properties-lemma-relocalize-morphism-at-schemes} 的结果。

\medskip\noindent
令 $U \in \Ob(X_\etale)$，令
$V \in \Ob(Y_\etale)$
，并令 $s \in g^{-1}h_V(U)$ 为一个截面。可以把
$s$ 看作层的映射 $s : h_U \to g^{-1}h_V$。由
《站点上的模》，
引理 \ref{sites-modules-lemma-relocalize-morphism-ringed-topoi}，
得到环化拓扑斯态射的交换图
$$
\xymatrix{
(\Sh(X_\etale/U), \mathcal{O}_U)
\ar[rr]_-{(j, j^\sharp)} \ar[d]_{(g_s, g_s^\sharp)} & &
(\Sh(X_\etale), \mathcal{O}_X) \ar[d]^{(g, g^\sharp)} \\
(\Sh(V_\etale), \mathcal{O}_V) \ar[rr] & &
(\Sh(Y_\etale), \mathcal{O}_Y).
}
$$
由
《\'Etale 上同调》，定理 \ref{etale-cohomology-theorem-fully-faithful}，
得到唯一的概形态射 $f_s : U \to V$，使得
$(g_s, g_s^\sharp)$ 与 $(f_{s, small}, f_s^\sharp)$ 为 $2$-同构。
刚才说明的构造 $(U, V, s) \leadsto f_s$ 满足如下函子性：
设有 $X_\etale$ 中的态射 $a : U' \to U$ 以及 $Y_\etale$ 中的态射
$b : V' \to V$，
并有映射 $s' : h_{U'} \to g^{-1}h_{V'}$，使下图交换：
$$
\xymatrix{
h_{U'} \ar[d]_a \ar[r]_{s'} & g^{-1}h_{V'} \ar[d]^{g^{-1}b} \\
h_U \ar[r]^s & g^{-1}h_V
}
$$
则下图（概形的图）交换：
$$
\xymatrix{
U' \ar[r]_-{f_{s'}} \ar[d]_a & u(V') \ar[d]^{u(b)} \\
U \ar[r]^-{f_s} & u(V)
}
$$
之所以如此，是因为对由《站点上的模》
引理 \ref{sites-modules-lemma-relocalize-morphism-ringed-topoi}
构造的态射 $(g_s, g_s^\sharp)$，同样的条件成立；
再利用《\'Etale 上同调》定理
\ref{etale-cohomology-theorem-fully-faithful} 的唯一性即可。

\medskip\noindent
问题在于将态射 $f_s$ 胶合为代数空间的态射。
为此，先选取概形 $V$ 以及满的 \'etale 态射 $V \to Y$。
这意味着 $h_V \to *$ 为满射，因而 $g^{-1}h_V \to *$ 也为满射。
于是存在概形 $U$、满的 \'etale 态射 $U \to X$ 以及态射
$s : h_U \to g^{-1}h_V$。接着令 $R = V \times_Y V$，
$R' = U \times_X U$。由于 $g^{-1}$ 正合，有
$g^{-1}h_R = g^{-1}h_V \times g^{-1}h_V$。
因此 $s$ 诱导态射 $s \times s : h_{R'} \to g^{-1}h_R$。
应用上述构造可知得到概形态射的交换图
$$
\xymatrix{
R' \ar@<1ex>[d] \ar@<-1ex>[d] \ar[rr]_{f_{s \times s}} & &
R \ar@<1ex>[d] \ar@<-1ex>[d] \\
U \ar[rr]^{f_s} & &
V
}
$$
由于 $X = U/R'$ 且 $Y = V/R$（参见
《空间》，引理 \ref{spaces-lemma-space-presentation}），
可知该图定义了代数空间态射 $f : X \to Y$，并且落在一个显然的交换图中。
现在还需证明 $(f_{small}, f^\sharp)$ 与 $(g, g^\sharp)$
是 $2$-同构的。
令 $t_V : f_{s, small}^{-1} \to g_s^{-1}$ 以及
$t_R : f_{s \times s, small}^{-1} \to g_{s \times s}^{-1}$
为上述构造给出的 $2$-同构。
令 $\mathcal{G}$ 为 $Y_\etale$ 上的层。则 $t_V$ 定义同构
$$
f_{small}^{-1}\mathcal{G}|_{U_\etale}
=
f_{s, small}^{-1}\mathcal{G}|_{V_\etale}
\xrightarrow{t_V}
g_s^{-1}\mathcal{G}|_{V_\etale}
=
g^{-1}\mathcal{G}|_{U_\etale}.
$$
此外，经任一投影 $R' \to U$ 拉回到 $R'$ 后，该同构就是
$$
f_{small}^{-1}\mathcal{G}|_{R'_\etale}
=
f_{s \times s, small}^{-1}\mathcal{G}|_{R_\etale}
\xrightarrow{t_R}
g_{s \times s}^{-1}\mathcal{G}|_{R_\etale}
=
g^{-1}\mathcal{G}|_{R'_\etale}.
$$
由于 $\{U \to X\}$ 是站点 $X_{spaces, \etale}$ 中的覆盖，
这意味着第一个显示的同构下降为层的同构
$t : f_{small}^{-1}\mathcal{G} \to g^{-1}\mathcal{G}$
（细节略去）。由于 $t_V$ 和 $t_R$ 是函子变换，该同构关于
$\mathcal{G}$ 具有函子性。最后，$t$ 与 $f^\sharp$、$g^\sharp$ 相容，
因为 $t_V$ 和 $t_R$ 具有这一性质（细节略去）。
定理得证。
\end{proof}

\begin{lemma}
\label{spaces-properties-lemma-isomorphism-ringed-topoi}
设 $X$、$Y$ 为 $\mathbf{Z}$ 上的代数空间。若
$$
(g, g^\sharp) :
(\Sh(X_\etale), \mathcal{O}_X)
\longrightarrow
(\Sh(Y_\etale), \mathcal{O}_Y)
$$
是环化拓扑斯的同构，则存在唯一的代数空间态射
$f : X \to Y$，使得 $(g, g^\sharp)$ 与 $(f_{small}, f^\sharp)$
同构，并且 $f$ 是代数空间的同构。
\end{lemma}

\begin{proof}
由定理 \ref{spaces-properties-theorem-fully-faithful}，只需证明 $(g, g^\sharp)$ 是
局部环化拓扑斯的态射。由《站点上的模》，
引理 \ref{sites-modules-lemma-locally-ringed-morphism}
（并且因为站点 $X_\etale$ 有足够多的点），只需检查
映射
$\mathcal{O}_{Y, q} \to \mathcal{O}_{X, p}$
由 $g^\sharp$ 诱导，在 $q = f \circ p$ 且 $p$ 为 $X_\etale$ 的任意点时，
是局部环同态。由于它是同构，结论显然。
\end{proof}














\section{代数空间上的拟凝聚层}
\label{spaces-properties-section-quasi-coherent}

\noindent
在《下降》的各节
\ref{descent-section-quasi-coherent-sheaves}、
\ref{descent-section-quasi-coherent-cohomology} 以及
\ref{descent-section-quasi-coherent-sheaves-bis} 中，我们已经看到，对于概形 $U$，
$U$ 上的拟凝聚 $\mathcal{O}_U$-模与 $U$ 的小 \'etale 站点上的拟凝聚
$\mathcal{O}$-模没有区别。因此，将以下定义应用于可表示代数空间时，
它与我们原先对概形上的拟凝聚层的概念相容
（《概形》，第 \ref{schemes-section-quasi-coherent} 节）。

\begin{definition}
\label{spaces-properties-definition-quasi-coherent}
设 $S$ 为概形，$X$ 为 $S$ 上的代数空间。
一个{\it 拟凝聚} $\mathcal{O}_X$-模
是环化站点 $(X_\etale, \mathcal{O}_X)$ 上的拟凝聚模，
其意义参见《站点上的模》，
定义 \ref{sites-modules-definition-site-local}。
$X$ 上拟凝聚层的范畴记作
$\QCoh(\mathcal{O}_X)$.
\end{definition}

\noindent
注意，拟凝聚性是内在的概念（参见《站点上的模》，
引理 \ref{sites-modules-lemma-special-locally-free}），
所以这等价于说，在 $X_{spaces, \etale}$ 上相应的 $\mathcal{O}_X$-模是拟凝聚的。

\medskip\noindent
如通常情形，拟凝聚层关于拉回具有良好性质。

\begin{lemma}
\label{spaces-properties-lemma-pullback-quasi-coherent}
设 $S$ 为概形。
设 $f : X \to Y$ 为 $S$ 上代数空间的态射。
拉回函子
$f^* : \textit{Mod}(\mathcal{O}_Y) \to \textit{Mod}(\mathcal{O}_X)$
保持拟凝聚层。
\end{lemma}

\begin{proof}
这是一般事实，参见《站点上的模》，引理
\ref{sites-modules-lemma-local-pullback}。
\end{proof}

\noindent
注意，如果 $X$、$Y$ 恰为概形，该拉回函子与模的拟凝聚层之间的通常拉回函子
一致，参见《下降》，命题
\ref{descent-proposition-equivalence-quasi-coherent-functorial}。
下面给出一个必要引理，将其与 $X$ 的小 \'etale 站点的对象上的拟凝聚层作比较。

\begin{lemma}
\label{spaces-properties-lemma-characterize-quasi-coherent-small-etale}
设 $S$ 为概形，$X$ 为 $S$ 上的代数空间。
一个拟凝聚 $\mathcal{O}_X$-模 $\mathcal{F}$
由以下数据给出：
\begin{enumerate}
\item 对每个 $U \in \Ob(X_\etale)$，在 $U_\etale$ 上的一个拟凝聚
$\mathcal{O}_U$-模 $\mathcal{F}_U$，
\item 对 $X_\etale$ 中的每个 $f : U' \to U$，一个同构
$c_f : f_{small}^*\mathcal{F}_U \to \mathcal{F}_{U'}$.
\end{enumerate}
这些数据须满足如下条件：给定 $X_\etale$ 中的任意
$f : U' \to U$ 和 $g : U'' \to U'$，复合
$c_g \circ g_{small}^*c_f$ 等于 $c_{f \circ g}$。
\end{lemma}

\begin{proof}
结合引理 \ref{spaces-properties-lemma-pullback-quasi-coherent} 与
\ref{spaces-properties-lemma-characterize-module-small-etale} 即得。
\end{proof}

\begin{lemma}
\label{spaces-properties-lemma-stalk-quasi-coherent}
设 $S$ 为概形。
设 $X$ 为 $S$ 上的代数空间。
设 $\mathcal{F}$ 为拟凝聚 $\mathcal{O}_X$-模。
设 $x \in |X|$ 为一点，$\overline{x}$ 为位于 $x$ 上方的几何点。
最后，设
$\varphi : (U, \overline{u}) \to (X, \overline{x})$
为 \'etale 邻域，其中 $U$ 为概形。则
$$
(\varphi^*\mathcal{F})_u \otimes_{\mathcal{O}_{U, u}}
\mathcal{O}_{X, \overline{x}} =
\mathcal{F}_{\overline{x}}
$$
其中 $u \in U$ 是 $\overline{u}$ 的像。
\end{lemma}

\begin{proof}
注意，由引理 \ref{spaces-properties-lemma-describe-etale-local-ring}，
$\mathcal{O}_{X, \overline{x}} = \mathcal{O}_{U, u}^{sh}$，
所以张量积是有意义的。此外，由定义 \ref{spaces-properties-definition-stalk} 可知
$$
\mathcal{F}_{\overline{u}} = \colim (\varphi^*\mathcal{F})_u
$$
其中余极限遍历引理中所述的
$\varphi : (U, \overline{u}) \to (X, \overline{x})$。
因此，引理陈述中的左侧到右侧有一个典范映射。
$\mathcal{O}_{X, \overline{x}}$ 也有类似的余极限描述，并且由
引理 \ref{spaces-properties-lemma-characterize-quasi-coherent-small-etale}，有
$$
((\varphi')^*\mathcal{F})_{u'} =
(\varphi^*\mathcal{F})_u \otimes_{\mathcal{O}_{U, u}} \mathcal{O}_{U', u'}
$$
只要 $(U', \overline{u}') \to (U, \overline{u})$ 是 \'etale 邻域的态射。
最后使用 $\otimes$ 与余极限交换即可完成证明。
\end{proof}

\begin{lemma}
\label{spaces-properties-lemma-stalk-pullback-quasi-coherent}
设 $S$ 为概形。设 $f : X \to Y$ 为 $S$ 上代数空间的态射。
设 $\mathcal{G}$ 为拟凝聚 $\mathcal{O}_Y$-模。
设 $\overline{x}$ 为 $X$ 的几何点，且
$\overline{y} = f \circ \overline{x}$ 为其在 $Y$ 中的像。
则存在典范同构
$$
(f^*\mathcal{G})_{\overline{x}} =
\mathcal{G}_{\overline{y}} \otimes_{\mathcal{O}_{Y, \overline{y}}}
\mathcal{O}_{X, \overline{x}}
$$
把拉回的茎同构于该茎与 $\overline{x}$ 处 $X$ 的局部环的张量积。
\end{lemma}

\begin{proof}
由于
$f^*\mathcal{G} = f_{small}^{-1}\mathcal{G}
\otimes_{f_{small}^{-1}\mathcal{O}_Y} \mathcal{O}_X$，
由引理 \ref{spaces-properties-lemma-stalk-pullback} 中对拉回茎的描述以及取茎与张量积交换的事实，
即得结论。也可以更直接地如下看出这一点。
选取交换图
$$
\xymatrix{
U \ar[d]_p \ar[r]_\alpha  & V \ar[d]^q \\
X \ar[r]^a & Y
}
$$
其中 $U$、$V$ 为概形，且 $p$、$q$ 为满的 \'etale 态射。
由引理 \ref{spaces-properties-lemma-geometric-lift-to-usual}，可以选取 $U$ 的几何点
$\overline{u}$，使得 $\overline{x} = p \circ \overline{u}$。
令 $\overline{v} = \alpha \circ \overline{u}$。于是
\begin{align*}
(f^*\mathcal{G})_{\overline{x}} & =
(p^*f^*\mathcal{G})_u \otimes_{\mathcal{O}_{U, u}}
\mathcal{O}_{X, \overline{x}} \\
& = (\alpha^*q^*\mathcal{G})_u \otimes_{\mathcal{O}_{U, u}}
\mathcal{O}_{X, \overline{x}} \\
& = (q^*\mathcal{G})_v \otimes_{\mathcal{O}_{V, v}}
\mathcal{O}_{U, u} \otimes_{\mathcal{O}_{U, u}}
\mathcal{O}_{X, \overline{x}} \\
& = (q^*\mathcal{G})_v \otimes_{\mathcal{O}_{V, v}}
\mathcal{O}_{X, \overline{x}} \\
& = (q^*\mathcal{G})_v \otimes_{\mathcal{O}_{V, v}}
\mathcal{O}_{Y, \overline{y}} \otimes_{\mathcal{O}_{Y, \overline{y}}}
\mathcal{O}_{X, \overline{x}} \\
& = \mathcal{G}_{\overline{y}} \otimes_{\mathcal{O}_{Y, \overline{y}}}
\mathcal{O}_{X, \overline{x}}
\end{align*}
这里使用了引理 \ref{spaces-properties-lemma-stalk-quasi-coherent}（两次）以及概形上拟凝聚层拉回的相应结果，
参见《层》，引理 \ref{sheaves-lemma-stalk-pullback-modules}。
\end{proof}

\begin{lemma}
\label{spaces-properties-lemma-characterize-quasi-coherent}
设 $S$ 为概形，$X$ 为 $S$ 上的代数空间。
设 $\mathcal{F}$ 为 $\mathcal{O}_X$-模层。以下条件等价：
\begin{enumerate}
\item $\mathcal{F}$ 是拟凝聚 $\mathcal{O}_X$-模；
\item 存在 $S$ 上代数空间的 \'etale 态射 $f : Y \to X$，其
$|f| : |Y| \to |X|$ 为满射，并且 $f^*\mathcal{F}$ 在 $Y$ 上拟凝聚；
\item 存在概形 $U$ 及满的 \'etale 态射
$\varphi : U \to X$，使得 $\varphi^*\mathcal{F}$ 是拟凝聚
$\mathcal{O}_U$-模；并且
\item 对每个仿射概形 $U$ 及 \'etale 态射 $\varphi : U \to X$，
限制 $\varphi^*\mathcal{F}$ 是拟凝聚 $\mathcal{O}_U$-模。
\end{enumerate}
\end{lemma}

\begin{proof}
考虑 $\text{id}_X$，显然 (1) 蕴含 (2)。
假设 $f : Y \to X$ 如 (2) 中所述，且 $V \to Y$ 是来自概形到 $Y$ 的满
\'etale 态射。则复合 $V \to X$ 也为满的 \'etale 态射，
并且由引理 \ref{spaces-properties-lemma-pullback-quasi-coherent}，$\mathcal{F}$
拉回到 $V$ 上仍拟凝聚。因此 (2) 蕴含 (3)。

\medskip\noindent
令 $U \to X$ 如 (3) 中所述。沿用等式
(\ref{spaces-properties-equation-restrict-modules}) 中引入的记号滥用。
由于 $\mathcal{F}|_{U_\etale}$ 拟凝聚，存在一个 \'etale 覆盖
$\{U_i \to U\}$，使得 $\mathcal{F}|_{U_{i, \etale}}$ 有全局表示，参见
《站点上的模》，定义 \ref{sites-modules-definition-global} 与
引理 \ref{sites-modules-lemma-local-final-object}。
令 $V \to X$ 为 $X_\etale$ 的一个对象。由于 $U \to X$ 满且 \'etale，
映射族 $\{U_i \times_X V \to V\}$ 是 $V$ 的 \'etale 覆盖。
通过态射 $U_i \times_X V \to U_i$，可以将
$\mathcal{F}|_{U_{i, \etale}}$ 的全局表示限制为
$\mathcal{F}|_{(U_i \times_X V)_\etale}$ 的全局表示。
因此 $X_\etale$ 上的层 $\mathcal{F}$ 满足《站点上的模》，定义
\ref{sites-modules-definition-site-local} 的条件，故为拟凝聚。

\medskip\noindent
 (3) 与 (4) 的等价性来自每个概形都有仿射开覆盖这一事实。
\end{proof}

\begin{lemma}
\label{spaces-properties-lemma-properties-quasi-coherent}
设 $S$ 为概形，$X$ 为 $S$ 上的代数空间。
$X$ 上拟凝聚层的范畴 $\QCoh(\mathcal{O}_X)$
具有以下性质：
\begin{enumerate}
\item 任意拟凝聚层的直和仍为拟凝聚。
\item 任意拟凝聚层的余极限仍为拟凝聚。
\item 拟凝聚层态射的核与余核是拟凝聚的。
\item 给定 $\mathcal{O}_X$-模的短正合序列
$0 \to \mathcal{F}_1 \to \mathcal{F}_2 \to \mathcal{F}_3 \to 0$
若其中任意两个是拟凝聚的，则第三个也是拟凝聚的。
\item 给定两个拟凝聚 $\mathcal{O}_X$-模，它们的张量积是拟凝聚的。
\item 给定两个拟凝聚 $\mathcal{O}_X$-模
$\mathcal{F}$、$\mathcal{G}$，其中 $\mathcal{F}$
是有限表示的（参见第 \ref{spaces-properties-section-properties-modules} 节），
则内部 Hom
$\SheafHom_{\mathcal{O}_X}(\mathcal{F}, \mathcal{G})$
是拟凝聚的。
\end{enumerate}
\end{lemma}

\begin{proof}
若 $X$ 为概形，这就是《下降》，引理
\ref{descent-lemma-properties-quasi-coherent} 的结论。
我们将通过 \'etale 局部化把本引理化归到这种情形。

\medskip\noindent
选取概形 $U$ 及满的 \'etale 态射 $\varphi : U \to X$。
我们约定记号 $\textit{Mod}(\mathcal{O}_U) =
\textit{Mod}(U_\etale, \mathcal{O}_U)$ 以及
$\QCoh(\mathcal{O}_U) = \QCoh(U_\etale, \mathcal{O}_U)$；换言之，
尽管 $U$ 是概形，我们仍将 $U$ 上的拟凝聚模视为 $U$ 的小 \'etale 站点上的模。
由引理 \ref{spaces-properties-lemma-pullback-quasi-coherent} 有交换图
$$
\xymatrix{
\QCoh(\mathcal{O}_X) \ar[r]_{\varphi^*} \ar[d] &
\QCoh(\mathcal{O}_U) \ar[d] \\
\textit{Mod}(\mathcal{O}_X) \ar[r]^{\varphi^*} &
\textit{Mod}(\mathcal{O}_U)
}
$$
下方水平箭头是限制函子
(\ref{spaces-properties-equation-restrict-modules})
$\mathcal{G} \mapsto \mathcal{G}|_{U_\etale}$。
该函子同时有左伴随和右伴随，参见《站点上的模》，第
\ref{sites-modules-section-localize} 节，故与所有极限和余极限交换。
此外，已知 $\textit{Mod}(\mathcal{O}_X)$ 的对象属于
$\QCoh(\mathcal{O}_X)$ 当且仅当其限制到 $U$ 后属于
$\QCoh(\mathcal{O}_U)$，参见引理 \ref{spaces-properties-lemma-characterize-quasi-coherent}。
完成这些准备后即可开始证明。

\medskip\noindent
性质 (1) 的证明。设 $\mathcal{F}_i$（$i \in I$）为一族拟凝聚
$\mathcal{O}_X$-模。由上面的讨论，有
$$
\Big(\bigoplus \mathcal{F}_i\Big)|_{U_\etale} =
\bigoplus \mathcal{F}_i|_{U_\etale}
$$
每个模 $\mathcal{F}_i|_{U_\etale}$ 都是拟凝聚的。
因此由概形情形，直和是拟凝聚的。
于是 $\bigoplus \mathcal{F}_i$ 作为限制到 $U$ 后为拟凝聚模的模，
也是拟凝聚的。

\medskip\noindent
性质 (2) 的证明。设 $\mathcal{I} \to \QCoh(\mathcal{O}_X)$、
$i \mapsto \mathcal{F}_i$ 为一个图表。则
$$
(\colim \mathcal{F}_i)|_{U_\etale} = \colim \mathcal{F}_i|_{U_\etale}
$$
由上面的讨论，这与上述相同方式给出结论。

\medskip\noindent
性质 (3) 的证明。设 $a : \mathcal{F} \to \mathcal{F}'$
为 $\QCoh(\mathcal{O}_X)$ 中的箭头。则
有 $\Ker(a)|_{U_\etale} = \Ker(a|_{U_\etale})$
以及 $\Coker(a)|_{U_\etale} = \Coker(a|_{U_\etale})$，
同样得到结论。

\medskip\noindent
性质 (4) 的证明。限制
$0 \to \mathcal{F}_1|_{U_\etale} \to \mathcal{F}_2|_{U_\etale}
\to \mathcal{F}_3|_{U_\etale} \to 0$
是短正合的。因此该序列具有二取三性质，结论如前。

\medskip\noindent
性质 (5) 的证明。设 $\mathcal{F}$、$\mathcal{G}$ 属于 $\QCoh(\mathcal{O}_X)$。
则有
$$
(\mathcal{F} \otimes_{\mathcal{O}_X} \mathcal{G})_{U_\etale} =
\mathcal{F}|_{U_\etale} \otimes_{\mathcal{O}_U} \mathcal{G}|_{U_\etale}
$$
结论如前。

\medskip\noindent
性质 (6) 的证明。设 $\mathcal{F}$、$\mathcal{G}$
属于 $\QCoh(\mathcal{O}_X)$，且 $\mathcal{F}$
有限表示。则有
$$
\SheafHom_{\mathcal{O}_X}(\mathcal{F}, \mathcal{G})|_{U_\etale} =
\SheafHom_{\mathcal{O}_U}(\mathcal{F}|_{U_\etale}, \mathcal{G}|_{U_\etale})
$$
事实上，限制是局部化，参见第 \ref{spaces-properties-section-localize} 节，尤其是公式
(\ref{spaces-properties-equation-localize-at-scheme-ringed}))；内部 Hom 的形成与局部化交换，参见
《站点上的模》，引理 \ref{sites-modules-lemma-internal-hom-restriction}。
因此结论如前。
\end{proof}

\noindent
一般而言，沿代数空间态射的拟凝聚层推前不一定是拟凝聚的。
我们将在《空间的态射》，第 \ref{spaces-morphisms-section-pushforward} 节中回到这一问题。




\section{模的性质}
\label{spaces-properties-section-properties-modules}

\noindent
在《站点上的模》的各节
\ref{sites-modules-section-global}、\ref{sites-modules-section-local} 以及
定义 \ref{sites-modules-definition-flat} 中，我们为任意环化拓扑斯上的
$\mathcal{O}$-模定义了一些内在性质。若 $X$ 为代数空间，我们将自由地把这些概念
应用于环化站点 $(X_\etale, \mathcal{O}_X)$ 上的模，或等价地应用于环化站点
$(X_{spaces, \etale}, \mathcal{O}_X)$ 上的模。

\medskip\noindent
全局性质 $\mathcal{P}$：
\begin{enumerate}
\item[(a)] {\it 自由}；
\item[(b)] {\it 有限自由}；
\item[(c)] {\it 由全局截面生成}；
\item[(d)] {\it 由有限多个全局截面生成}；
\item[(e)] 具有{\it 全局表示}；以及
\item[(f)] 具有{\it 全局有限表示}。
\end{enumerate}
局部性质 $\mathcal{P}$：
\begin{enumerate}
\item[(g)] {\it 局部自由}；
\item[(f)] {\it 有限局部自由}；
\item[(h)] {\it 局部由截面生成}；
\item[(i)] {\it 局部由 $r$ 个截面生成}；
\item[(j)] {\it 有限型}；
\item[(k)] {\it 拟凝聚}（参见第 \ref{spaces-properties-section-quasi-coherent} 节）；
\item[(l)] {\it 有限表示}；
\item[(m)] {\it 相干}；以及
\item[(n)] {\it 平坦}。
\end{enumerate}
下面是一些由定义立即得到的结果：
\begin{enumerate}
\item 除 $\mathcal{P}=$“相干”外，在每种情形下该性质都在拉回下保持，参见
《站点上的模》，引理 \ref{sites-modules-lemma-global-pullback}、
\ref{sites-modules-lemma-local-pullback} 以及
\ref{sites-modules-lemma-pullback-flat}。
\item 上述每个性质（包括相干性）都在代数空间的 \'etale 态射拉回下保持
（因为此时拉回由限制给出，参见引理 \ref{spaces-properties-lemma-etale-morphism-topoi}）。
\item 假设 $f : Y \to X$ 是代数空间的满 \'etale 态射。对于局部性质
(g)--(m)，若 $f^*\mathcal{F}$ 具有 $\mathcal{P}$，则 $\mathcal{F}$
也具有 $\mathcal{P}$。这是因为 $\{Y \to X\}$ 是
$X_{spaces, \etale}$ 中的覆盖，并应用《站点上的模》引理
\ref{sites-modules-lemma-local-final-object}。
\item 若 $X$ 为概形，$\mathcal{F}$ 为 $X_\etale$ 上的拟凝聚模，且
$\mathcal{P}$ 为除“相干”或“局部自由”外的任一性质，则 $X_\etale$ 上
$\mathcal{F}$ 的 $\mathcal{P}$ 等价于 $\mathcal{F}|_{X_{Zar}}$ 的相应性质；
也就是说，把它看作概形 $X$ 上的拟凝聚层时，$\mathcal{F}$ 具有 $\mathcal{P}$。
参见《下降》，引理 \ref{descent-lemma-equivalence-quasi-coherent-properties}。
\item 若 $X$ 为局部 Noether 概形，$\mathcal{F}$ 为 $X_\etale$ 上的
拟凝聚模，则 $\mathcal{F}$ 在 $X_\etale$ 上相干，当且仅当
$\mathcal{F}|_{X_{Zar}}$ 相干；也就是说，这正对应于通常意义下
概形 $X$ 上的相干层。参见《下降》，引理
\ref{descent-lemma-equivalence-quasi-coherent-properties}。
\end{enumerate}



\section{局部投射模}
\label{spaces-properties-section-locally-projective}

\noindent
回忆一下，在《性质》，第 \ref{properties-section-locally-projective} 节中，
我们定义了局部投射拟凝聚模的概念。

\begin{lemma}
\label{spaces-properties-lemma-locally-projective}
设 $S$ 为概形，$X$ 为 $S$ 上的代数空间。
设 $\mathcal{F}$ 为拟凝聚 $\mathcal{O}_X$-模。以下条件等价：
\begin{enumerate}
\item 存在概形 $U$ 及满的 \'etale 态射 $U \to X$，使得限制
$\mathcal{F}|_U$ 在 $U$ 上局部投射；并且
\item 对任意概形 $U$ 及任意 \'etale 态射 $U \to X$，限制
$\mathcal{F}|_U$ 在 $U$ 上局部投射。
\end{enumerate}
\end{lemma}

\begin{proof}
令 $U \to X$ 如 (1) 中所述，并令 $V \to X$ 为 \'etale 态射，其中
$V$ 为概形。则 $\{U \times_X V \to V\}$ 是概形的 fppf 覆盖。
因此，若 $\mathcal{F}|_U$ 局部投射，则 $\mathcal{F}|_{U \times_X V}$
局部投射（参见《性质》，引理
\ref{properties-lemma-locally-projective-pullback}），
从而 $\mathcal{F}|_V$ 局部投射，参见《下降》，引理
\ref{descent-lemma-locally-projective-descends}。
\end{proof}

\begin{definition}
\label{spaces-properties-definition-locally-projective}
设 $S$ 为概形，$X$ 为 $S$ 上的代数空间。
设 $\mathcal{F}$ 为拟凝聚 $\mathcal{O}_X$-模。
若引理 \ref{spaces-properties-lemma-locally-projective} 的等价条件成立，
则称 $\mathcal{F}$ 为{\it 局部投射的}。
\end{definition}

\begin{lemma}
\label{spaces-properties-lemma-locally-projective-pullback}
设 $S$ 为概形。
设 $f : X \to Y$ 为 $S$ 上代数空间的态射。
设 $\mathcal{G}$ 为拟凝聚 $\mathcal{O}_Y$-模。
若 $\mathcal{G}$ 在 $Y$ 上局部投射，则 $f^*\mathcal{G}$
在 $X$ 上局部投射。
\end{lemma}

\begin{proof}
选取满的 \'etale 态射 $V \to Y$，其中 $V$ 为概形。
选取满的 \'etale 态射 $U \to V \times_Y X$，其中
$U$ 为概形。记诱导的态射为 $\psi : U \to V$。则
$$
f^*\mathcal{G}|_U = \psi^*(\mathcal{G}|_V)
$$
因此由定义以及概形情形的结果可得引理，参见《性质》，引理
\ref{properties-lemma-locally-projective-pullback}。
\end{proof}





\section{拟凝聚层与表示}
\label{spaces-properties-section-quasi-coherent-presentation}

\noindent
设 $S$ 为概形，$X$ 为 $S$ 上的代数空间。
令 $X = U/R$ 为由任意满的 \'etale 态射 $\varphi : U \to X$
得到的 $X$ 的一个表示，参见《空间》，定义
\ref{spaces-definition-presentation}。
特别地，我们得到群胚 $(U, R, s, t, c)$，其中
$j = (t, s) : R \to U \times_S U$，参见《群胚》，引理
\ref{groupoids-lemma-equivalence-groupoid}。
在《群胚》，定义 \ref{groupoids-definition-groupoid-module} 中，
我们定义了任意群胚上的拟凝聚层概念。具备这些记号后，有如下观察。

\begin{proposition}
\label{spaces-properties-proposition-quasi-coherent}
设 $S$、$\varphi : U \to X$ 以及 $(U, R, s, t, c)$ 如上。
对任意拟凝聚 $\mathcal{O}_X$-模 $\mathcal{F}$，层
$\varphi^*\mathcal{F}$ 带有如下典范同构：
$$
\alpha : t^*\varphi^*\mathcal{F} \longrightarrow s^*\varphi^*\mathcal{F}
$$
它满足《群胚》，定义 \ref{groupoids-definition-groupoid-module} 的条件，
因而在 $(U, R, s, t, c)$ 上定义拟凝聚层。
函子 $\mathcal{F} \mapsto (\varphi^*\mathcal{F}, \alpha)$
定义范畴等价
$$
\begin{matrix}
\text{Quasi-coherent} \\
\mathcal{O}_X\text{-模}
\end{matrix}
\longleftrightarrow
\begin{matrix}
\text{Quasi-coherent modules}\\
\text{定义于 }(U, R, s, t, c)
\end{matrix}
$$
\end{proposition}

\begin{proof}
在命题陈述及本证明中，我们把概形上的拟凝聚层视为该概形小
\'etale 站点上的拟凝聚层。这由《下降》各节的结果是允许的：
\ref{descent-section-quasi-coherent-sheaves}、
\ref{descent-section-quasi-coherent-cohomology} 以及
\ref{descent-section-quasi-coherent-sheaves-bis}。

\medskip\noindent
由于 $\varphi \circ t = \varphi \circ s$ 且拉回关于态射具有函子性，
便得到 $\alpha$ 的存在性，参见等式 (\ref{spaces-properties-equation-push-pull}) 周围的讨论。
完全同样地，即利用拉回的函子性，可知同构 $\alpha$ 满足《群胚》，
定义 \ref{groupoids-definition-groupoid-module} 的条件 (1)。
要验证定义的条件 (2)，只需注意 $\alpha$ 是同构，这显然成立。
构造 $\mathcal{F} \mapsto (\varphi^*\mathcal{F}, \alpha)$
显然关于拟凝聚层 $\mathcal{F}$ 具有函子性。
因此得到引理显示公式中从左到右的函子。

\medskip\noindent
反之，设 $(\mathcal{F}, \alpha)$ 是 $(U, R, s, t, c)$ 上的拟凝聚层。
令 $V \to X$ 为 $X_\etale$ 的一个对象。此时态射
$V' = U \times_X V \to V$ 是概形的满 \'etale 态射，故
$\{V' \to V\}$ 是 $V$ 的 \'etale 覆盖。此外，拟凝聚层 $\mathcal{F}$
拉回为 $V'$ 上的拟凝聚层 $\mathcal{F}'$。
由于 $R = U \times_X U$ 且 $t = \text{pr}_0$、$s = \text{pr}_0$，
可知 $V' \times_V V' = R \times_X V$，且投影映射
$V' \times_V V' \to V'$ 等于 $t$、$s$ 的拉回。因此 $\alpha$ 拉回为同构
$\alpha' : \text{pr}_0^*\mathcal{F}' \to \text{pr}_1^*\mathcal{F}'$，
并且二元组 $(\mathcal{F}', \alpha')$ 是关于 $\{V' \to V\}$ 的
拟凝聚层下降数据。由《下降》，命题
\ref{descent-proposition-fpqc-descent-quasi-coherent}，
该下降数据有效，从而得到 $V_\etale$ 上的拟凝聚
$\mathcal{O}_V$-模 $\mathcal{F}_V$。
要看出这给出了 $X_\etale$ 上的拟凝聚层，需要证明（由
引理 \ref{spaces-properties-lemma-characterize-quasi-coherent-small-etale}）
对于 $X_\etale$ 中的任意态射 $f : V_1 \to V_2$，存在典范同构
$c_f : \mathcal{F}_{V_1} \to \mathcal{F}_{V_2}$
且与态射复合相容。验证略去。
我们也略去如下验证：这定义了一个从右侧范畴到左侧范畴的函子，
并且它是上述函子的逆。
\end{proof}

\begin{proposition}
\label{spaces-properties-proposition-coherator}
设 $S$ 为概形，$X$ 为 $S$ 上的代数空间。
\begin{enumerate}
\item 范畴 $\QCoh(\mathcal{O}_X)$ 是 Grothendieck 阿贝尔范畴。
因此 $\QCoh(\mathcal{O}_X)$ 有足够多的内射对象并且有所有极限。
\item 包含函子
$\QCoh(\mathcal{O}_X) \to \textit{Mod}(\mathcal{O}_X)$
有右伴随\footnote{该函子有时称为{\it 相干化函子}。}
$$
Q : \textit{Mod}(\mathcal{O}_X) \longrightarrow \QCoh(\mathcal{O}_X)
$$
使得对每个拟凝聚层 $\mathcal{F}$，伴随映射
$Q(\mathcal{F}) \to \mathcal{F}$ 是同构。
\end{enumerate}
\end{proposition}

\begin{proof}
本证明重复概形情形的证明，参见《性质》，命题
\ref{properties-proposition-coherator}。
建议读者先阅读该证明。

\medskip\noindent
性质 (1) 意味着 $\QCoh(\mathcal{O}_X)$：(a) 有所有余极限；
(b) 滤余极限正合；(c) 有生成元，参见《内射》，第
\ref{injectives-section-grothendieck-conditions} 节。
由引理 \ref{spaces-properties-lemma-properties-quasi-coherent}，
$\QCoh(\mathcal{O}_X)$ 中的余极限存在并且与
$\textit{Mod}(\mathcal{O}_X)$ 中的余极限一致。由《站点上的模》，引理
\ref{sites-modules-lemma-limits-colimits}，滤余极限正合。
因此 (a)、(b) 成立。

\medskip\noindent
为构造一个生成元，选取表示 $X = U/R$，使得
$(U, R, s, t, c)$ 是一个
\'etale 群胚概形，特别地 $s$ 和 $t$ 是概形的平坦态射。选取《群胚》引理
\ref{groupoids-lemma-colimit-kappa} 中的一个基数 $\kappa$。
选取一族 $(\mathcal{E}_t, \alpha_t)_{t \in T}$，其中的对象是
$(U, R, s, t, c)$ 上的 $\kappa$-生成拟凝聚模，且如《群胚》引理
\ref{groupoids-lemma-set-of-iso-classes} 所述。
令 $\mathcal{F}_t$ 为 $X$ 上的拟凝聚模，它通过范畴等价
命题 \ref{spaces-properties-proposition-quasi-coherent} 对应于拟凝聚模
$(\mathcal{E}_t, \alpha_t)$。
于是可见，每个拟凝聚模 $\mathcal{H}$ 都是其拟凝聚子模的
有向余极限，而这些子模同构于某个 $\mathcal{F}_t$。因此
$\bigoplus_t \mathcal{F}_t$ 是 $\QCoh(\mathcal{O}_X)$ 的一个生成元，
从而 (c) 成立。关于极限和内射对象的断言在任意
Grothendieck 阿贝尔范畴中成立，参见《内射》定理
\ref{injectives-theorem-injective-embedding-grothendieck} 以及
引理 \ref{injectives-lemma-grothendieck-products}。

\medskip\noindent
\medskip\noindent
(2) 的证明。为构造 $Q$，我们使用如下的一般过程。
给定 $\textit{Mod}(\mathcal{O}_X)$ 中的一个对象 $\mathcal{F}$，
考虑函子
$$
\QCoh(\mathcal{O}_X)^{opp} \longrightarrow \textit{Sets},\quad
\mathcal{G} \longmapsto \Hom_X(\mathcal{G}, \mathcal{F})
$$
该函子把余极限变为极限，因而是可表示的，参见《内射》引理
\ref{injectives-lemma-grothendieck-brown}。
因此存在一个拟凝聚层 $Q(\mathcal{F})$ 以及函子性同构
$\Hom_X(\mathcal{G}, \mathcal{F}) = \Hom_X(\mathcal{G}, Q(\mathcal{F}))$
对 $\QCoh(\mathcal{O}_X)$ 中的 $\mathcal{G}$ 成立。由 Yoneda 引理
（《范畴》，引理 \ref{categories-lemma-yoneda}），构造
$\mathcal{F} \leadsto Q(\mathcal{F})$ 关于 $\mathcal{F}$ 具有函子性。
按构造，$Q$ 是包含函子的右伴随。
当 $\mathcal{F}$ 为拟凝聚层时，$Q(\mathcal{F}) \to \mathcal{F}$
为同构这一事实，在形式上源于包含函子
$\QCoh(\mathcal{O}_X) \to \textit{Mod}(\mathcal{O}_X)$
是全忠实的。
\end{proof}






\section{指向概形的态射}
\label{spaces-properties-section-morphisms-to-schemes}

\noindent
这是《概形》引理 \ref{schemes-lemma-morphism-into-affine} 的类似结果。

\begin{lemma}
\label{spaces-properties-lemma-morphism-to-affine-scheme}
设 $X$ 是 $\mathbf{Z}$ 上的代数空间。
设 $T$ 是仿射概形。映射
$$
\Mor(X, T)
\longrightarrow
\Hom(\Gamma(T, \mathcal{O}_T), \Gamma(X, \mathcal{O}_X))
$$
将 $f$ 映到 $f^\sharp$（在全局截面上）的映射是双射。
\end{lemma}

\begin{proof}
我们构造该映射的逆。令
$\varphi : \Gamma(T, \mathcal{O}_T) \to \Gamma(X, \mathcal{O}_X)$
为环映射。选取表示 $X = U/R$，参见《空间》定义
\ref{spaces-definition-presentation}。由
《概形》引理 \ref{schemes-lemma-morphism-into-affine}，复合映射
$$
\Gamma(T, \mathcal{O}_T) \to \Gamma(X, \mathcal{O}_X) \to
\Gamma(U, \mathcal{O}_U)
$$
对应于唯一的概形态射 $g : U \to T$。由同一引理，两个复合
$R \to U \to T$ 相等。因此得到态射
$f : X = U/R \to T$，使得 $U \to X \to T$ 等于 $g$。按构造，图表
$$
\xymatrix{
\Gamma(U, \mathcal{O}_U) & \Gamma(X, \mathcal{O}_X) \ar[l] \\
& \Gamma(T, \mathcal{O}_T) \ar[lu]^{g^\sharp} \ar[u]^{\varphi}_{f^\sharp}
}
$$
交换。因此，由于 $U \to X$ 是 \'etale 覆盖且 $\mathcal{O}_X$ 是
$X_\etale$ 上的层，$f^\sharp$ 等于 $\varphi$。
$f$ 的唯一性由 $g$ 的唯一性推出。
\end{proof}





\section{自由作用的商}
\label{spaces-properties-section-quotient-by-free}

\noindent
设 $S$ 为概形。
设 $X$ 为 $S$ 上的代数空间。
设 $G$ 为抽象群。
令 $a : G \to \text{Aut}(X)$ 为同态，即 $a$ 是 $G$ 在 $X$ 上的
{\it 作用}。若对每个 $S$ 上的概形 $T$，映射
$$
G \times X(T) \longrightarrow X(T)
$$
是自由的，则称该作用为 {\it 自由}。（我们不能使用《空间》引理
\ref{spaces-lemma-quotient} 中的判据，因为点可能没有定义良好的剩余域。）
在作用自由的情形，我们将把商 $X/G$ 构造为代数空间。这是一般结果
《引导》引理 \ref{bootstrap-lemma-quotient-free-action} 的特殊情形；
该引理将在后文证明。

\begin{lemma}
\label{spaces-properties-lemma-quotient}
设 $S$ 为概形。
设 $X$ 为 $S$ 上的代数空间。
设 $G$ 为在 $X$ 上自由作用的抽象群。
则商层 $X/G$ 是代数空间。
\end{lemma}

\begin{proof}
该陈述意味着，与预层
$$
T \longmapsto X(T)/G
$$
相关联的层 $F$
是一个代数空间。为看出这一点，我们构造一个表示。选取概形 $U$ 及满的
\'etale 态射 $\varphi : U \to X$。令 $V = \coprod_{g \in G} U$，并令
$\psi : V \to X$ 在对应于 $g \in G$ 的分支上等于
$a(g) \circ \varphi$。令 $G$ 通过置换分支作用在 $V$ 上，即
$g_0 \in G$ 通过 $U$ 的恒等态射，将对应于 $g$ 的分支映到对应于
$g_0g$ 的分支。于是 $\psi$ 是 $G$-等变态射，也就是说我们归结为
下一段所处理的情形。

\medskip\noindent
假设 $U$ 上存在 $G$-作用，且 $U \to X$ 是满的、\'etale 且 $G$-等变的。
此时 $R = U \times_X U$ 上诱导出一个 $G$-作用，并且与投影映射
$t, s : R \to U$ 相容。现在我们断言
$$
X/G = U/\coprod\nolimits_{g \in G} R
$$
其中映射
$$
j : \coprod\nolimits_{g \in G} R
\longrightarrow
U \times_S U
$$
由 $(r, g) \mapsto (t(r), g(s(r)))$ 给出。注意 $j$ 是单态射：
若 $(t(r), g(s(r))) = (t(r'), g'(s(r')))$，则
$t(r) = t(r')$，所以 $r$ 与 $r'$ 在 $X$ 中分别经 $s$ 和 $t$ 的像相同；
于是 $g = g'$（因为 $G$ 在 $X$ 上自由作用），从而
$s(r) = s(r')$，最后 $r = r'$（因为 $R$ 是 $U$ 上的等价关系）。
此外，$j$ 是等价关系（细节略）。
两个投影 $\coprod\nolimits_{g \in G} R \to U$ 都是 \'etale 的，因为
$s$ 和 $t$ 是 \'etale 的。因此 $j$ 是 \'etale 等价关系，并且由
《空间》定理 \ref{spaces-theorem-presentation}，
$U/\coprod\nolimits_{g \in G} R$ 是代数空间。
存在映射
$$
U/\coprod\nolimits_{g \in G} R \longrightarrow X/G
$$
由映射 $U \to X$ 诱导。它是层的同构；证明略。
\end{proof}
